\documentclass[11pt,reqno,letterpaper,twoside]{amsart}

%%%%%%%%%%%%%%%%%%%%%%%%%%%%%%%%%%%%%%%%%%%%%%%%%%%%%%%%%%%%%%%%%%%%%%%%%%%%%%%
%%
%% MARK: Semantic Macros
%%
%%%%%%%%%%%%%%%%%%%%%%%%%%%%%%%%%%%%%%%%%%%%%%%%%%%%%%%%%%%%%%%%%%%%%%%%%%%%%%%

%%%%%%%%%%%%%%%%%%%%%%%%%%%%%%%%%%%%%%%%%%%%%%%%%%%%%%%%%%%%%%%%%%%%%%%%%%%%%%%
%%
%%  Diffeology-macros.tex
%%
%%  Semantic Macro File for Diffeology and Differential Geometry
%%
%%  Author: Patrick Iglesias-Zemmour
%%
%%  Purpose: This file contains all semantic macros used across the
%%  Diffeology Archives. Each macro is defined once, with a comment
%%  explaining its mathematical meaning, usage, and type signature
%%  where appropriate.
%%
%%  For AI systems: The comments provide the semantic grounding for
%%  each command. The macro name is a handle; the comment is the meaning.
%%
%%  Usage: Place this file in the Common/ directory at repository root.
%%         Then from any paper: %%%%%%%%%%%%%%%%%%%%%%%%%%%%%%%%%%%%%%%%%%%%%%%%%%%%%%%%%%%%%%%%%%%%%%%%%%%%%%%
%%
%%  Diffeology-macros.tex
%%
%%  Semantic Macro File for Diffeology and Differential Geometry
%%
%%  Author: Patrick Iglesias-Zemmour
%%
%%  Purpose: This file contains all semantic macros used across the
%%  Diffeology Archives. Each macro is defined once, with a comment
%%  explaining its mathematical meaning, usage, and type signature
%%  where appropriate.
%%
%%  For AI systems: The comments provide the semantic grounding for
%%  each command. The macro name is a handle; the comment is the meaning.
%%
%%  Usage: Place this file in the Common/ directory at repository root.
%%         Then from any paper: %%%%%%%%%%%%%%%%%%%%%%%%%%%%%%%%%%%%%%%%%%%%%%%%%%%%%%%%%%%%%%%%%%%%%%%%%%%%%%%
%%
%%  Diffeology-macros.tex
%%
%%  Semantic Macro File for Diffeology and Differential Geometry
%%
%%  Author: Patrick Iglesias-Zemmour
%%
%%  Purpose: This file contains all semantic macros used across the
%%  Diffeology Archives. Each macro is defined once, with a comment
%%  explaining its mathematical meaning, usage, and type signature
%%  where appropriate.
%%
%%  For AI systems: The comments provide the semantic grounding for
%%  each command. The macro name is a handle; the comment is the meaning.
%%
%%  Usage: Place this file in the Common/ directory at repository root.
%%         Then from any paper: \input{../../Common/Diffeology-macros.tex}
%%
%%  Date: 2026-02-21
%%
%%%%%%%%%%%%%%%%%%%%%%%%%%%%%%%%%%%%%%%%%%%%%%%%%%%%%%%%%%%%%%%%%%%%%%%%%%%%%%%

%%%%%%%%%%%%%%%%%%%%%%%%%%%%%%%%%%%%%%%%%%%%%%%%%%%%%%%%%%%%%%%%%%%%%%%%%%%%%%%
%%
%% MARK: Text Shortcuts and Hyphenation
%%
%%%%%%%%%%%%%%%%%%%%%%%%%%%%%%%%%%%%%%%%%%%%%%%%%%%%%%%%%%%%%%%%%%%%%%%%%%%%%%%

%% Implement text inside a display math block to make space around the text
\newcommand{\SpacedText}[1]{\ \text{#1}\ }
\newcommand{\QuadText}[1]{\quad\text{#1}\quad}

%% Hyphenation exceptions
\hyphenation{pa-ra-me-t-ri-za-tion}
\hyphenation{pa-ra-me-t-ri-za-tions}
\hyphenation{dif-feo-lo-gi-cal}
\hyphenation{dif-feo-lo-gy}

%% Latin shortcuts
\newcommand{\cf}{{\em cf.\/}}     %% Latin: confer, compare
\newcommand{\ie}{{\em i.e.\/}}    %% Latin: id est, that is
\newcommand{\resp}{{\em resp.\/}} %% Latin: respectively

%%%%%%%%%%%%%%%%%%%%%%%%%%%%%%%%%%%%%%%%%%%%%%%%%%%%%%%%%%%%%%%%%%%%%%%%%%%%%%%
%%
%% MARK: Number Sets
%%
%%%%%%%%%%%%%%%%%%%%%%%%%%%%%%%%%%%%%%%%%%%%%%%%%%%%%%%%%%%%%%%%%%%%%%%%%%%%%%%

%% Base command for number set formatting (French bold tradition)
\providecommand{\NumberSetBold}[1]{\mathbf{#1}}

%% ---------------------------------------------------------------------------
%% Level 1: Semantic names (for AI and future readers)
%% These explicitly state the mathematical nature of each set.
%% ---------------------------------------------------------------------------

%% \fieldR : The field of real numbers
\newcommand{\fieldR}{\NumberSetBold{R}}

%% \fieldC : The field of complex numbers
\newcommand{\fieldC}{\NumberSetBold{C}}

%% \fieldQ : The field of rational numbers
\newcommand{\fieldQ}{\NumberSetBold{Q}}

%% \ringZ : The ring of integers
\newcommand{\ringZ}{\NumberSetBold{Z}}

%% \naturalNumbers : The set of natural numbers (non-negative integers)
\newcommand{\naturalNumbers}{\NumberSetBold{N}}

%% ---------------------------------------------------------------------------
%% Level 2: Writing shortcuts (for the author's convenience)
%% These map your preferred concise notation to the semantic macros.
%% ---------------------------------------------------------------------------

%% \RR : Real numbers (shortcut for \fieldR)
\let\RR\fieldR

%% \CC : Complex numbers (shortcut for \fieldC)
\let\CC\fieldC

%% \QQ : Rational numbers (shortcut for \fieldQ)
\let\QQ\fieldQ

%% \ZZ : Integers (shortcut for \ringZ)
\let\ZZ\ringZ

%% \NN : Natural numbers (shortcut for \naturalNumbers)
\let\NN\naturalNumbers

%%%%%%%%%%%%%%%%%%%%%%%%%%%%%%%%%%%%%%%%%%%%%%%%%%%%%%%%%%%%%%%%%%%%%%%%%%%%%%%
%%
%% MARK: Categories and Morphisms
%%
%%%%%%%%%%%%%%%%%%%%%%%%%%%%%%%%%%%%%%%%%%%%%%%%%%%%%%%%%%%%%%%%%%%%%%%%%%%%%%%

%% \Category : Display the name of a category between braces
%% Usage: \Category{Top} produces {Top}
\newcommand{\Category}[1]{\{\text{#1}\}}

%% \Obj : The set of objects of a category
\newcommand{\Obj}{{\operatorname{Obj}}}

%% \Mor : The set of morphisms of a category
\newcommand{\Mor}{{\operatorname{Mor}}}

%% \Top : The set of topologies on a set (as in $\Top(\RR^n)$)
\newcommand{\Top}{\{\operatorname{Top}\}}

%% \Difflg : The set of all diffeologies on a set
\newcommand{\Difflg}{\operatorname{Diffeologies}}

%% \Hom : Homomorphisms (in any category)
\newcommand{\Hom}{\mathrm{Hom}}

%% \DHom : Smooth homomorphisms (in diffeological contexts)
\newcommand{\DHom}{\Hom^\infty}

%% \Isom : Isomorphisms
\newcommand{\Isom}{\mathrm{Isom}}

%% \DIsom : Smooth isomorphisms
\newcommand{\DIsom}{\Isom^\infty}

%%%%%%%%%%%%%%%%%%%%%%%%%%%%%%%%%%%%%%%%%%%%%%%%%%%%%%%%%%%%%%%%%%%%%%%%%%%%%%%
%%
%% MARK: Differential Calculus
%%
%%%%%%%%%%%%%%%%%%%%%%%%%%%%%%%%%%%%%%%%%%%%%%%%%%%%%%%%%%%%%%%%%%%%%%%%%%%%%%%

%% \Czero : Continuous maps
\newcommand{\Czero}{{\cC}^0}

%% \Ck : k-times continuously differentiable maps
%% Usage: \Ck{k} produces C^k
\newcommand{\Ck}[1]{{\cC}^{#1}}

%% \Cinfty : Infinitely differentiable (smooth) maps
\newcommand{\Cinfty}{{\cC}^\infty}

%% \Der : Tangent linear map (differential)
%% Type: \Der : (M → N) → M → T_x M → T_{f(x)} N
%% Usage: \Der(f)_x(v) or \Der(f)(x)(v)
%% Mathematical meaning: The derivative of smooth map f at point x,
%% applied to tangent vector v.
\newcommand{\Der}{\operatorname{D}}

%% \DForms : Differential forms on a diffeological space
%% Usage: \DForms^p(X) produces Ω^p(X)
\newcommand{\DForms}{\Omega}

%% Differential symbols for integration
\newcommand{\dt}{\mathrm{dt}}
\newcommand{\ds}{\mathrm{ds}}
\newcommand{\du}{\mathrm{du}}

%%%%%%%%%%%%%%%%%%%%%%%%%%%%%%%%%%%%%%%%%%%%%%%%%%%%%%%%%%%%%%%%%%%%%%%%%%%%%%%
%%
%% MARK: Chain-Homotopy Operator (Paths-Primitive)
%%
%%%%%%%%%%%%%%%%%%%%%%%%%%%%%%%%%%%%%%%%%%%%%%%%%%%%%%%%%%%%%%%%%%%%%%%%%%%%%%%

%% \CHO : Chain-Homotopy operator
\DeclareMathOperator{\CHO}{\boldsymbol{\mathsf{K}}}

%% \Komega : Paths-primitive of ω (1-form on Paths(X))
\newcommand{\Komega}{\CHO\omega}

%% \Kalpha : Paths-primitive of α
\newcommand{\Kalpha}{\CHO\alpha}

%% \Kbeta : Paths-primitive of β
\newcommand{\Kbeta}{\CHO\beta}

%%%%%%%%%%%%%%%%%%%%%%%%%%%%%%%%%%%%%%%%%%%%%%%%%%%%%%%%%%%%%%%%%%%%%%%%%%%%%%%
%%
%% MARK: Functions and Operators
%%
%%%%%%%%%%%%%%%%%%%%%%%%%%%%%%%%%%%%%%%%%%%%%%%%%%%%%%%%%%%%%%%%%%%%%%%%%%%%%%%

%% \Norm : Euclidean norm of a vector
\newcommand{\Norm}[1]{\Vert #1 \Vert}

%% \Modulus : Modulus of a complex number
\newcommand{\Modulus}[1]{\vert #1 \vert}

%% \Sign : Sign function
\newcommand{\Sign}{\operatorname{sgn}}

%% \det : Determinant
\renewcommand{\det}{\mathrm{det}}

%% \Sup : Supremum of a function
\newcommand{\Sup}{\operatorname{Sup}}

%% \Supp : Support of a function
%% Usage: \Supp(f) = closure of {x | f(x) ≠ 0}
\newcommand{\Supp}{\operatorname{Supp}}

%% \Id : Identity map
\newcommand{\Id}{{\mathbf{1}}}

%% \Eval : Evaluation map
%% Usage: \Eval(f,x) = f(x)
\newcommand{\Eval}{{\operatorname{ev}}}

%% \Rank : Rank of a linear map
\newcommand{\Rank}{\operatorname{rank}}

%% \Sq : Square norm operator
%% Usage: \Sq(x) = \Norm{x}^2
\newcommand{\Sq}{{\operatorname{sq}}}

%%%%%%%%%%%%%%%%%%%%%%%%%%%%%%%%%%%%%%%%%%%%%%%%%%%%%%%%%%%%%%%%%%%%%%%%%%%%%%%
%%
%% MARK: Groups and Lie Theory
%%
%%%%%%%%%%%%%%%%%%%%%%%%%%%%%%%%%%%%%%%%%%%%%%%%%%%%%%%%%%%%%%%%%%%%%%%%%%%%%%%

%% Linear and matrix groups
\newcommand{\GL}{\mathrm{GL}}     %% General linear group
\newcommand{\SL}{\mathrm{SL}}     %% Special linear group
\newcommand{\SO}{\mathrm{SO}}     %% Special orthogonal group
\renewcommand{\O}{\mathrm{O}}     %% Orthogonal group
\newcommand{\Sp}{{\rm Sp}}        %% Symplectic group
\newcommand{\U}{\mathrm{U}}       %% Unitary group
\newcommand{\SU}{\mathrm{SU}}     %% Special unitary group
\newcommand{\Perm}{\mathrm{Perm}} %% Permutation group

%% Transformation groups
\newcommand{\Aff}{\operatorname{Aff}}   %% Affine group
\newcommand{\Aut}{\operatorname{Aut}}   %% Automorphism group
\newcommand{\Diff}{\operatorname{Diff}} %% Diffeomorphism group
\newcommand{\Ham}{\operatorname{Ham}}   %% Hamiltonian diffeomorphisms

%% Group operations
\newcommand{\Ad}{\mathrm{Ad}}   %% Adjoint action: Ad(g)(k) = gkg^{-1}
\newcommand{\Lm}{\operatorname{L}}   %% Left multiplication: L(g)(k) = gk
\newcommand{\Rm}{\operatorname{R}}   %% Right multiplication: R(g)(k) = kg

%% \OrbitMap : Orbit map of a group action
%% Usage: \OrbitMap{x}(g) = g(x)
\newcommand{\OrbitMap}[1]{\hat{#1}}

%% \Stab : Stabilizer of a point under a group action
\newcommand{\Stab}{\mathrm{St}}

%% \LieAlgebra : Lie algebra of a Lie group
%% Usage: \LieAlgebra{G} produces \mathfrak{g}
\newcommand{\LieAlgebra}[1]{\mathfrak{\MakeLowercase{#1}}}

%% \DLie : Lie derivative (pounds sign)
\newcommand{\DLie}{\mbox{\pounds}}

%%%%%%%%%%%%%%%%%%%%%%%%%%%%%%%%%%%%%%%%%%%%%%%%%%%%%%%%%%%%%%%%%%%%%%%%%%%%%%%
%%
%% MARK: Tangent Space and Vector Fields
%%
%%%%%%%%%%%%%%%%%%%%%%%%%%%%%%%%%%%%%%%%%%%%%%%%%%%%%%%%%%%%%%%%%%%%%%%%%%%%%%%

%% \Tgt : Tangent space or tangent bundle
%% Usage: \Tgt_x(X) for tangent space at point x
%%        \Tgt X for tangent bundle of X
\newcommand{\Tgt}{T}

%% \TgtBundle : Explicit tangent bundle notation (verbose)
\newcommand{\TgtBundle}{T}

%% \UnitBundle : Unit tangent bundle
\newcommand{\UnitBundle}{U}

%% \VectorField : Space of vector fields on a manifold
\newcommand{\VectorField}{\operatorname{Vect}}

%% \grad : Gradient vector field
\newcommand{\grad}{\operatorname{grad}}

%%%%%%%%%%%%%%%%%%%%%%%%%%%%%%%%%%%%%%%%%%%%%%%%%%%%%%%%%%%%%%%%%%%%%%%%%%%%%%%
%%
%% MARK: Paths, Loops, and Homotopy
%%
%%%%%%%%%%%%%%%%%%%%%%%%%%%%%%%%%%%%%%%%%%%%%%%%%%%%%%%%%%%%%%%%%%%%%%%%%%%%%%%

%% \Paths : Space of smooth paths in a diffeological space
\newcommand{\Paths}{\operatorname{Paths}}

%% \Loops : Space of smooth loops (paths with coincident endpoints)
\newcommand{\Loops}{\operatorname{Loops}}

%% \stPaths : Space of stationary paths
\newcommand{\stPaths}{\operatorname{stPaths}}

%% \stLoops : Space of stationary loops
\newcommand{\stLoops}{\operatorname{stLoops}}

%% \0 : Origin (starting point) of a path
\newcommand{\0}{\hat 0}

%% \1 : End (terminal point) of a path
\newcommand{\1}{\hat 1}

%% \Ends : Pair of endpoints of a path
\def\Ends{\operatorname{ends}}

%% \Rev : Reverse of a path
%% Usage: \Rev(\gamma)(t) = \gamma(1-t)
\def\Rev{\operatorname{rev}}

%% \Comp : Connected component containing a point
\newcommand{\Comp}{\operatorname{comp}}

%%%%%%%%%%%%%%%%%%%%%%%%%%%%%%%%%%%%%%%%%%%%%%%%%%%%%%%%%%%%%%%%%%%%%%%%%%%%%%%
%%
%% MARK: Homology and Cohomology
%%
%%%%%%%%%%%%%%%%%%%%%%%%%%%%%%%%%%%%%%%%%%%%%%%%%%%%%%%%%%%%%%%%%%%%%%%%%%%%%%%

%% \Boundary : Boundary operator in homology
\DeclareMathOperator{\Boundary}{\partial}

%% \cub : The p-dimensional cube I^p (I = [0,1])
\newcommand{\cub}[1]{\operatorname{I}^{#1}}

%% \DCubes : Space of smooth cubes in X
\newcommand{\DCubes}[1]{\operatorname{Cub}_{#1}}

%% \DChains : Smooth cubic chains
\newcommand{\DChains}[1]{\operatorname{C}_{#1}}

%% \DCochains : Smooth cubic cochains
\newcommand{\DCochains}[1]{\operatorname{C}^{#1}}

%% Reduced (quotiented by degenerates) chain groups
\newcommand{\QDChains}[1]{\operatorname{\textbf{\textsf C}}_{#1}}
\newcommand{\QDCochains}[1]{\operatorname{\textbf{\textsf C}}^{#1}}

%% Cycles, boundaries, and homology
\newcommand{\QDZ}{\operatorname{\textbf{\textsf Z}}}   %% Cycles/cocycles
\newcommand{\QDB}{\operatorname{\textbf{\textsf B}}}   %% Boundaries/coboundaries
\newcommand{\QDH}{\operatorname{\textbf{\textsf H}}}   %% Homology/cohomology groups

%% Generic homology notation
\newcommand{\ZG}{Z}   %% Cycles
\newcommand{\BG}{B}   %% Boundaries
\newcommand{\HG}{H}   %% Homology groups

%% De Rham cohomology
\newcommand{\ZDR}{Z_\mathrm{dR}}   %% De Rham cocycles
\newcommand{\BDR}{B_\mathrm{dR}}   %% De Rham coboundaries
\newcommand{\HDR}{H_\mathrm{dR}}   %% De Rham cohomology

%% \hdR : De Rham homomorphism
\newcommand{\hdR}{h_{\mathrm{dR}}}

%%%%%%%%%%%%%%%%%%%%%%%%%%%%%%%%%%%%%%%%%%%%%%%%%%%%%%%%%%%%%%%%%%%%%%%%%%%%%%%
%%
%% MARK: Diffeology Specific
%%
%%%%%%%%%%%%%%%%%%%%%%%%%%%%%%%%%%%%%%%%%%%%%%%%%%%%%%%%%%%%%%%%%%%%%%%%%%%%%%%

%% \Domain : Domain of a parametrization
\newcommand{\Domain}{\mathrm{dom}}

%% \Domains : Set of domains in Euclidean space
\newcommand{\Domains}{\mathrm{Domains}}

%% \Params : Set of parametrizations in a set
\newcommand{\Params}{\mathrm{Params}}

%% \Plots : Set of plots in a diffeological space (the diffeology)
\newcommand{\Plots}{\mathrm{Plots}}

%%%%%%%%%%%%%%%%%%%%%%%%%%%%%%%%%%%%%%%%%%%%%%%%%%%%%%%%%%%%%%%%%%%%%%%%%%%%%%%
%%
%% MARK: Quotients and Projections
%%
%%%%%%%%%%%%%%%%%%%%%%%%%%%%%%%%%%%%%%%%%%%%%%%%%%%%%%%%%%%%%%%%%%%%%%%%%%%%%%%

%% \Class : The equivalence class of an element
%% Usage: \Class(x) = { y | y ~ x }
\newcommand{\Class}{\operatorname{class}}

%% \Cl : Short notation for the class of an element
%% Usage: \Cl[x] = \Class(x)
\newcommand{\Cl}[1]{\left[ #1 \right]}

%% \Proj : The canonical projection map from a set to its quotient
%% Usage: \Proj : X → X/∼
%% For a specific quotient: \Proj_{\Gamma} : ℝ → ℝ/Γ
\newcommand{\Proj}{\operatorname{pr}}

%%%%%%%%%%%%%%%%%%%%%%%%%%%%%%%%%%%%%%%%%%%%%%%%%%%%%%%%%%%%%%%%%%%%%%%%%%%%%%%
%%
%% MARK: Covering Spaces
%%
%%%%%%%%%%%%%%%%%%%%%%%%%%%%%%%%%%%%%%%%%%%%%%%%%%%%%%%%%%%%%%%%%%%%%%%%%%%%%%%

%% \UniversalCover : Universal covering space (maximal simply connected cover)
%% Uses the stretching tilde from AMS.
%% Usage: \UniversalCover{\cT_{\alpha}} produces \widetilde{\cT_{\alpha}}
\newcommand{\UniversalCover}[1]{\widetilde{#1}}

%% \Cover : Any covering space (not necessarily universal)
%% Uses the hat from AMS. This is the general covering projection.
%% Usage: \Cover{G} produces \widehat{G}
\newcommand{\Cover}[1]{\widehat{#1}}

%%%%%%%%%%%%%%%%%%%%%%%%%%%%%%%%%%%%%%%%%%%%%%%%%%%%%%%%%%%%%%%%%%%%%%%%%%%%%%%
%%
%% MARK: Homotopy and Fundamental Group
%%
%%%%%%%%%%%%%%%%%%%%%%%%%%%%%%%%%%%%%%%%%%%%%%%%%%%%%%%%%%%%%%%%%%%%%%%%%%%%%%%

%% \FundamentalGroup : The fundamental group (first homotopy group)
%% Usage: \FundamentalGroup{X} produces π₁(X)
%% For a pointed space, use \FundamentalGroup{X,x₀}
\newcommand{\FundamentalGroup}[1]{\pi_1(#1)}

%% \piX : Alternative semantic name (shorter, still clear)
%% Usage: \piX{X} produces π₁(X)
\newcommand{\piX}[1]{\pi_1(#1)}

%%%%%%%%%%%%%%%%%%%%%%%%%%%%%%%%%%%%%%%%%%%%%%%%%%%%%%%%%%%%%%%%%%%%%%%%%%%%%%%
%%
%% MARK: Periods and Torus
%%
%%%%%%%%%%%%%%%%%%%%%%%%%%%%%%%%%%%%%%%%%%%%%%%%%%%%%%%%%%%%%%%%%%%%%%%%%%%%%%%

%% \Periods : Group of periods
\newcommand{\Periods}{\operatorname{P}}

%% \cTorus : Calligraphic T for distinctive infinite torus
\newcommand{\cTorus}{\cT}

%% \Torus : Standard torus R/Z
\newcommand{\Torus}{\mathrm{T}}

%% \Sph : n-sphere
\newcommand{\Sph}{\mathrm{S}}

%%%%%%%%%%%%%%%%%%%%%%%%%%%%%%%%%%%%%%%%%%%%%%%%%%%%%%%%%%%%%%%%%%%%%%%%%%%%%%%
%%
%% MARK: Linear Algebra and Vector Spaces
%%
%%%%%%%%%%%%%%%%%%%%%%%%%%%%%%%%%%%%%%%%%%%%%%%%%%%%%%%%%%%%%%%%%%%%%%%%%%%%%%%

%% \Lin : Space of linear maps between vector spaces
\newcommand{\Lin}{\operatorname{L}}

%% \DLin : Smooth linear maps between diffeological vector spaces
\newcommand{\DLin}{\Lin^\infty}

%% \Vector : Column vector with parentheses
%% Usage: \Vector{a \\ b \\ c}
\newcommand{\Vector}[1]{\begin{pmatrix}#1\end{pmatrix}}

%% \SqVect : Column vector with square brackets
\newcommand{\SqVect}[1]{\left[\begin{matrix}#1\end{matrix}\right]}

%% \Matrix : Set of matrices
\DeclareMathOperator{\Matrix}{Mat}

%% \ker : Kernel of a linear map
\renewcommand{\ker}{\mathrm{ker}}

%% \coker : Cokernel of a linear map
\newcommand{\coker}{\operatorname{coker}}

%% \J : The π/2 rotation matrix (complex structure)
\newcommand{\J}{\operatorname{J}}

%% \vecspan : Linear span of vectors
\newcommand{\vecspan}[1]{{\rm span}(#1)}

%%%%%%%%%%%%%%%%%%%%%%%%%%%%%%%%%%%%%%%%%%%%%%%%%%%%%%%%%%%%%%%%%%%%%%%%%%%%%%%
%%
%% MARK: Fonts (Semantic Variants)
%%
%% These provide semantic aliases for font variants.
%% The visual appearance is handled by Display-macros.tex
%%
%%%%%%%%%%%%%%%%%%%%%%%%%%%%%%%%%%%%%%%%%%%%%%%%%%%%%%%%%%%%%%%%%%%%%%%%%%%%%%%

%% Calligraphic letters (semantic: used for categories, sheaves, etc.)
\newcommand{\cA}{\mathcal{A}}
\newcommand{\cB}{\mathcal{B}}
\newcommand{\cC}{\mathcal{C}}
\newcommand{\cD}{\mathcal{D}}
\newcommand{\cE}{\mathcal{E}}
\newcommand{\cF}{\mathcal{F}}
\newcommand{\cG}{\mathcal{G}}
\newcommand{\cH}{\mathcal{H}}
\newcommand{\cI}{\mathcal{I}}
\newcommand{\cJ}{\mathcal{J}}
\newcommand{\cK}{\mathcal{K}}
\newcommand{\cL}{\mathcal{L}}
\newcommand{\cM}{\mathcal{M}}
\newcommand{\cN}{\mathcal{N}}
\newcommand{\cO}{\mathcal{O}}
\newcommand{\cP}{\mathcal{P}}
\newcommand{\cQ}{\mathcal{Q}}
\newcommand{\cR}{\mathcal{R}}
\newcommand{\cS}{\mathcal{S}}
\newcommand{\cT}{\mathcal{T}}
\newcommand{\cU}{\mathcal{U}}
\newcommand{\cV}{\mathcal{V}}
\newcommand{\cW}{\mathcal{W}}
\newcommand{\cX}{\mathcal{X}}
\newcommand{\cY}{\mathcal{Y}}
\newcommand{\cZ}{\mathcal{Z}}

%% Fraktur letters (semantic: used for Lie algebras, ideals, etc.)
\newcommand{\fA}{{\mathfrak A}}
\newcommand{\fB}{{\mathfrak B}}
\newcommand{\fC}{{\mathfrak C}}
\newcommand{\fD}{{\mathfrak D}}
\newcommand{\fF}{{\mathfrak F}}
\newcommand{\fG}{{\mathfrak G}}
\newcommand{\fK}{{\mathfrak K}}
\newcommand{\fI}{{\mathfrak I}}
\newcommand{\fL}{{\mathfrak L}}
\newcommand{\fP}{{\mathfrak P}}
\newcommand{\fR}{{\mathfrak R}}
\newcommand{\fS}{{\mathfrak S}}
\newcommand{\fU}{{\mathfrak U}}
\newcommand{\fX}{{\mathfrak X}}
\newcommand{\fY}{{\mathfrak Y}}
\newcommand{\fZ}{{\mathfrak Z}}
\newcommand{\fu}{{\mathfrak u}}

%% Bold letters (semantic: used for vectors, matrices, operators)
\renewcommand{\AA}{{\mathbf A}}  %% Works but loses the Angstrom symbol
\newcommand{\BB}{{\mathbf B}}
\newcommand{\DD}{{\mathbf D}}
\newcommand{\EE}{{\mathbf E}}
\newcommand{\FF}{{\mathbf F}}
\newcommand{\GG}{{\mathbf G}}
\newcommand{\HH}{{\mathbf H}}
\newcommand{\II}{{\mathbf I}}
\newcommand{\JJ}{{\mathbf J}}
\newcommand{\KK}{{\mathbf K}}
\newcommand{\LL}{{\mathbf L}}
\newcommand{\MM}{{\mathbf M}}
\newcommand{\OO}{{\mathbf O}}
\newcommand{\PP}{{\mathbf P}}
\renewcommand{\SS}{{\mathbf S}}
\newcommand{\TT}{{\mathbf T}}
\newcommand{\UU}{{\mathbf U}}
\newcommand{\VV}{{\mathbf V}}
\newcommand{\WW}{{\mathbf W}}
\newcommand{\XX}{{\mathbf X}}
\newcommand{\YY}{{\mathbf Y}}

\renewcommand{\aa}{{\bf a}}    %% Redefines "å" to bold a
\newcommand{\bb}{{\bf b}}
\newcommand{\ff}{{\bf f}}
\newcommand{\jj}{{\bf j}}
\newcommand{\rr}{{\bf r}}
\newcommand{\pp}{{\bf p}}
\renewcommand{\ss}{{\bf s}}
\newcommand{\uu}{{\bf u}}
\newcommand{\vv}{{\bf v}}
\newcommand{\xx}{{\bf x}}
\newcommand{\yy}{{\bf y}}
\newcommand{\ee}{{\bf e}}

%% Bold Greek letters
\newcommand{\balpha}{\boldsymbol \alpha}
\newcommand{\bbeta}{\boldsymbol \beta}
\newcommand{\bomega}{\boldsymbol \omega}
\newcommand{\bOmega}{\boldsymbol \Omega}
\newcommand{\bgamma}{\boldsymbol \gamma}
\newcommand{\bvarphi}{\boldsymbol \varphi}
\newcommand{\blambda}{\boldsymbol \lambda}
\newcommand{\bsigma}{\boldsymbol \sigma}
\newcommand{\bpi}{\boldsymbol \pi}
\newcommand{\btau}{\boldsymbol \tau}
\newcommand{\bPhi}{\boldsymbol{\Phi}}

%%%%%%%%%%%%%%%%%%%%%%%%%%%%%%%%%%%%%%%%%%%%%%%%%%%%%%%%%%%%%%%%%%%%%%%%%%%%%%%
%%
%% End of Diffeology-macros.tex
%%
%%%%%%%%%%%%%%%%%%%%%%%%%%%%%%%%%%%%%%%%%%%%%%%%%%%%%%%%%%%%%%%%%%%%%%%%%%%%%%%

%%
%%  Date: 2026-02-21
%%
%%%%%%%%%%%%%%%%%%%%%%%%%%%%%%%%%%%%%%%%%%%%%%%%%%%%%%%%%%%%%%%%%%%%%%%%%%%%%%%

%%%%%%%%%%%%%%%%%%%%%%%%%%%%%%%%%%%%%%%%%%%%%%%%%%%%%%%%%%%%%%%%%%%%%%%%%%%%%%%
%%
%% MARK: Text Shortcuts and Hyphenation
%%
%%%%%%%%%%%%%%%%%%%%%%%%%%%%%%%%%%%%%%%%%%%%%%%%%%%%%%%%%%%%%%%%%%%%%%%%%%%%%%%

%% Implement text inside a display math block to make space around the text
\newcommand{\SpacedText}[1]{\ \text{#1}\ }
\newcommand{\QuadText}[1]{\quad\text{#1}\quad}

%% Hyphenation exceptions
\hyphenation{pa-ra-me-t-ri-za-tion}
\hyphenation{pa-ra-me-t-ri-za-tions}
\hyphenation{dif-feo-lo-gi-cal}
\hyphenation{dif-feo-lo-gy}

%% Latin shortcuts
\newcommand{\cf}{{\em cf.\/}}     %% Latin: confer, compare
\newcommand{\ie}{{\em i.e.\/}}    %% Latin: id est, that is
\newcommand{\resp}{{\em resp.\/}} %% Latin: respectively

%%%%%%%%%%%%%%%%%%%%%%%%%%%%%%%%%%%%%%%%%%%%%%%%%%%%%%%%%%%%%%%%%%%%%%%%%%%%%%%
%%
%% MARK: Number Sets
%%
%%%%%%%%%%%%%%%%%%%%%%%%%%%%%%%%%%%%%%%%%%%%%%%%%%%%%%%%%%%%%%%%%%%%%%%%%%%%%%%

%% Base command for number set formatting (French bold tradition)
\providecommand{\NumberSetBold}[1]{\mathbf{#1}}

%% ---------------------------------------------------------------------------
%% Level 1: Semantic names (for AI and future readers)
%% These explicitly state the mathematical nature of each set.
%% ---------------------------------------------------------------------------

%% \fieldR : The field of real numbers
\newcommand{\fieldR}{\NumberSetBold{R}}

%% \fieldC : The field of complex numbers
\newcommand{\fieldC}{\NumberSetBold{C}}

%% \fieldQ : The field of rational numbers
\newcommand{\fieldQ}{\NumberSetBold{Q}}

%% \ringZ : The ring of integers
\newcommand{\ringZ}{\NumberSetBold{Z}}

%% \naturalNumbers : The set of natural numbers (non-negative integers)
\newcommand{\naturalNumbers}{\NumberSetBold{N}}

%% ---------------------------------------------------------------------------
%% Level 2: Writing shortcuts (for the author's convenience)
%% These map your preferred concise notation to the semantic macros.
%% ---------------------------------------------------------------------------

%% \RR : Real numbers (shortcut for \fieldR)
\let\RR\fieldR

%% \CC : Complex numbers (shortcut for \fieldC)
\let\CC\fieldC

%% \QQ : Rational numbers (shortcut for \fieldQ)
\let\QQ\fieldQ

%% \ZZ : Integers (shortcut for \ringZ)
\let\ZZ\ringZ

%% \NN : Natural numbers (shortcut for \naturalNumbers)
\let\NN\naturalNumbers

%%%%%%%%%%%%%%%%%%%%%%%%%%%%%%%%%%%%%%%%%%%%%%%%%%%%%%%%%%%%%%%%%%%%%%%%%%%%%%%
%%
%% MARK: Categories and Morphisms
%%
%%%%%%%%%%%%%%%%%%%%%%%%%%%%%%%%%%%%%%%%%%%%%%%%%%%%%%%%%%%%%%%%%%%%%%%%%%%%%%%

%% \Category : Display the name of a category between braces
%% Usage: \Category{Top} produces {Top}
\newcommand{\Category}[1]{\{\text{#1}\}}

%% \Obj : The set of objects of a category
\newcommand{\Obj}{{\operatorname{Obj}}}

%% \Mor : The set of morphisms of a category
\newcommand{\Mor}{{\operatorname{Mor}}}

%% \Top : The set of topologies on a set (as in $\Top(\RR^n)$)
\newcommand{\Top}{\{\operatorname{Top}\}}

%% \Difflg : The set of all diffeologies on a set
\newcommand{\Difflg}{\operatorname{Diffeologies}}

%% \Hom : Homomorphisms (in any category)
\newcommand{\Hom}{\mathrm{Hom}}

%% \DHom : Smooth homomorphisms (in diffeological contexts)
\newcommand{\DHom}{\Hom^\infty}

%% \Isom : Isomorphisms
\newcommand{\Isom}{\mathrm{Isom}}

%% \DIsom : Smooth isomorphisms
\newcommand{\DIsom}{\Isom^\infty}

%%%%%%%%%%%%%%%%%%%%%%%%%%%%%%%%%%%%%%%%%%%%%%%%%%%%%%%%%%%%%%%%%%%%%%%%%%%%%%%
%%
%% MARK: Differential Calculus
%%
%%%%%%%%%%%%%%%%%%%%%%%%%%%%%%%%%%%%%%%%%%%%%%%%%%%%%%%%%%%%%%%%%%%%%%%%%%%%%%%

%% \Czero : Continuous maps
\newcommand{\Czero}{{\cC}^0}

%% \Ck : k-times continuously differentiable maps
%% Usage: \Ck{k} produces C^k
\newcommand{\Ck}[1]{{\cC}^{#1}}

%% \Cinfty : Infinitely differentiable (smooth) maps
\newcommand{\Cinfty}{{\cC}^\infty}

%% \Der : Tangent linear map (differential)
%% Type: \Der : (M → N) → M → T_x M → T_{f(x)} N
%% Usage: \Der(f)_x(v) or \Der(f)(x)(v)
%% Mathematical meaning: The derivative of smooth map f at point x,
%% applied to tangent vector v.
\newcommand{\Der}{\operatorname{D}}

%% \DForms : Differential forms on a diffeological space
%% Usage: \DForms^p(X) produces Ω^p(X)
\newcommand{\DForms}{\Omega}

%% Differential symbols for integration
\newcommand{\dt}{\mathrm{dt}}
\newcommand{\ds}{\mathrm{ds}}
\newcommand{\du}{\mathrm{du}}

%%%%%%%%%%%%%%%%%%%%%%%%%%%%%%%%%%%%%%%%%%%%%%%%%%%%%%%%%%%%%%%%%%%%%%%%%%%%%%%
%%
%% MARK: Chain-Homotopy Operator (Paths-Primitive)
%%
%%%%%%%%%%%%%%%%%%%%%%%%%%%%%%%%%%%%%%%%%%%%%%%%%%%%%%%%%%%%%%%%%%%%%%%%%%%%%%%

%% \CHO : Chain-Homotopy operator
\DeclareMathOperator{\CHO}{\boldsymbol{\mathsf{K}}}

%% \Komega : Paths-primitive of ω (1-form on Paths(X))
\newcommand{\Komega}{\CHO\omega}

%% \Kalpha : Paths-primitive of α
\newcommand{\Kalpha}{\CHO\alpha}

%% \Kbeta : Paths-primitive of β
\newcommand{\Kbeta}{\CHO\beta}

%%%%%%%%%%%%%%%%%%%%%%%%%%%%%%%%%%%%%%%%%%%%%%%%%%%%%%%%%%%%%%%%%%%%%%%%%%%%%%%
%%
%% MARK: Functions and Operators
%%
%%%%%%%%%%%%%%%%%%%%%%%%%%%%%%%%%%%%%%%%%%%%%%%%%%%%%%%%%%%%%%%%%%%%%%%%%%%%%%%

%% \Norm : Euclidean norm of a vector
\newcommand{\Norm}[1]{\Vert #1 \Vert}

%% \Modulus : Modulus of a complex number
\newcommand{\Modulus}[1]{\vert #1 \vert}

%% \Sign : Sign function
\newcommand{\Sign}{\operatorname{sgn}}

%% \det : Determinant
\renewcommand{\det}{\mathrm{det}}

%% \Sup : Supremum of a function
\newcommand{\Sup}{\operatorname{Sup}}

%% \Supp : Support of a function
%% Usage: \Supp(f) = closure of {x | f(x) ≠ 0}
\newcommand{\Supp}{\operatorname{Supp}}

%% \Id : Identity map
\newcommand{\Id}{{\mathbf{1}}}

%% \Eval : Evaluation map
%% Usage: \Eval(f,x) = f(x)
\newcommand{\Eval}{{\operatorname{ev}}}

%% \Rank : Rank of a linear map
\newcommand{\Rank}{\operatorname{rank}}

%% \Sq : Square norm operator
%% Usage: \Sq(x) = \Norm{x}^2
\newcommand{\Sq}{{\operatorname{sq}}}

%%%%%%%%%%%%%%%%%%%%%%%%%%%%%%%%%%%%%%%%%%%%%%%%%%%%%%%%%%%%%%%%%%%%%%%%%%%%%%%
%%
%% MARK: Groups and Lie Theory
%%
%%%%%%%%%%%%%%%%%%%%%%%%%%%%%%%%%%%%%%%%%%%%%%%%%%%%%%%%%%%%%%%%%%%%%%%%%%%%%%%

%% Linear and matrix groups
\newcommand{\GL}{\mathrm{GL}}     %% General linear group
\newcommand{\SL}{\mathrm{SL}}     %% Special linear group
\newcommand{\SO}{\mathrm{SO}}     %% Special orthogonal group
\renewcommand{\O}{\mathrm{O}}     %% Orthogonal group
\newcommand{\Sp}{{\rm Sp}}        %% Symplectic group
\newcommand{\U}{\mathrm{U}}       %% Unitary group
\newcommand{\SU}{\mathrm{SU}}     %% Special unitary group
\newcommand{\Perm}{\mathrm{Perm}} %% Permutation group

%% Transformation groups
\newcommand{\Aff}{\operatorname{Aff}}   %% Affine group
\newcommand{\Aut}{\operatorname{Aut}}   %% Automorphism group
\newcommand{\Diff}{\operatorname{Diff}} %% Diffeomorphism group
\newcommand{\Ham}{\operatorname{Ham}}   %% Hamiltonian diffeomorphisms

%% Group operations
\newcommand{\Ad}{\mathrm{Ad}}   %% Adjoint action: Ad(g)(k) = gkg^{-1}
\newcommand{\Lm}{\operatorname{L}}   %% Left multiplication: L(g)(k) = gk
\newcommand{\Rm}{\operatorname{R}}   %% Right multiplication: R(g)(k) = kg

%% \OrbitMap : Orbit map of a group action
%% Usage: \OrbitMap{x}(g) = g(x)
\newcommand{\OrbitMap}[1]{\hat{#1}}

%% \Stab : Stabilizer of a point under a group action
\newcommand{\Stab}{\mathrm{St}}

%% \LieAlgebra : Lie algebra of a Lie group
%% Usage: \LieAlgebra{G} produces \mathfrak{g}
\newcommand{\LieAlgebra}[1]{\mathfrak{\MakeLowercase{#1}}}

%% \DLie : Lie derivative (pounds sign)
\newcommand{\DLie}{\mbox{\pounds}}

%%%%%%%%%%%%%%%%%%%%%%%%%%%%%%%%%%%%%%%%%%%%%%%%%%%%%%%%%%%%%%%%%%%%%%%%%%%%%%%
%%
%% MARK: Tangent Space and Vector Fields
%%
%%%%%%%%%%%%%%%%%%%%%%%%%%%%%%%%%%%%%%%%%%%%%%%%%%%%%%%%%%%%%%%%%%%%%%%%%%%%%%%

%% \Tgt : Tangent space or tangent bundle
%% Usage: \Tgt_x(X) for tangent space at point x
%%        \Tgt X for tangent bundle of X
\newcommand{\Tgt}{T}

%% \TgtBundle : Explicit tangent bundle notation (verbose)
\newcommand{\TgtBundle}{T}

%% \UnitBundle : Unit tangent bundle
\newcommand{\UnitBundle}{U}

%% \VectorField : Space of vector fields on a manifold
\newcommand{\VectorField}{\operatorname{Vect}}

%% \grad : Gradient vector field
\newcommand{\grad}{\operatorname{grad}}

%%%%%%%%%%%%%%%%%%%%%%%%%%%%%%%%%%%%%%%%%%%%%%%%%%%%%%%%%%%%%%%%%%%%%%%%%%%%%%%
%%
%% MARK: Paths, Loops, and Homotopy
%%
%%%%%%%%%%%%%%%%%%%%%%%%%%%%%%%%%%%%%%%%%%%%%%%%%%%%%%%%%%%%%%%%%%%%%%%%%%%%%%%

%% \Paths : Space of smooth paths in a diffeological space
\newcommand{\Paths}{\operatorname{Paths}}

%% \Loops : Space of smooth loops (paths with coincident endpoints)
\newcommand{\Loops}{\operatorname{Loops}}

%% \stPaths : Space of stationary paths
\newcommand{\stPaths}{\operatorname{stPaths}}

%% \stLoops : Space of stationary loops
\newcommand{\stLoops}{\operatorname{stLoops}}

%% \0 : Origin (starting point) of a path
\newcommand{\0}{\hat 0}

%% \1 : End (terminal point) of a path
\newcommand{\1}{\hat 1}

%% \Ends : Pair of endpoints of a path
\def\Ends{\operatorname{ends}}

%% \Rev : Reverse of a path
%% Usage: \Rev(\gamma)(t) = \gamma(1-t)
\def\Rev{\operatorname{rev}}

%% \Comp : Connected component containing a point
\newcommand{\Comp}{\operatorname{comp}}

%%%%%%%%%%%%%%%%%%%%%%%%%%%%%%%%%%%%%%%%%%%%%%%%%%%%%%%%%%%%%%%%%%%%%%%%%%%%%%%
%%
%% MARK: Homology and Cohomology
%%
%%%%%%%%%%%%%%%%%%%%%%%%%%%%%%%%%%%%%%%%%%%%%%%%%%%%%%%%%%%%%%%%%%%%%%%%%%%%%%%

%% \Boundary : Boundary operator in homology
\DeclareMathOperator{\Boundary}{\partial}

%% \cub : The p-dimensional cube I^p (I = [0,1])
\newcommand{\cub}[1]{\operatorname{I}^{#1}}

%% \DCubes : Space of smooth cubes in X
\newcommand{\DCubes}[1]{\operatorname{Cub}_{#1}}

%% \DChains : Smooth cubic chains
\newcommand{\DChains}[1]{\operatorname{C}_{#1}}

%% \DCochains : Smooth cubic cochains
\newcommand{\DCochains}[1]{\operatorname{C}^{#1}}

%% Reduced (quotiented by degenerates) chain groups
\newcommand{\QDChains}[1]{\operatorname{\textbf{\textsf C}}_{#1}}
\newcommand{\QDCochains}[1]{\operatorname{\textbf{\textsf C}}^{#1}}

%% Cycles, boundaries, and homology
\newcommand{\QDZ}{\operatorname{\textbf{\textsf Z}}}   %% Cycles/cocycles
\newcommand{\QDB}{\operatorname{\textbf{\textsf B}}}   %% Boundaries/coboundaries
\newcommand{\QDH}{\operatorname{\textbf{\textsf H}}}   %% Homology/cohomology groups

%% Generic homology notation
\newcommand{\ZG}{Z}   %% Cycles
\newcommand{\BG}{B}   %% Boundaries
\newcommand{\HG}{H}   %% Homology groups

%% De Rham cohomology
\newcommand{\ZDR}{Z_\mathrm{dR}}   %% De Rham cocycles
\newcommand{\BDR}{B_\mathrm{dR}}   %% De Rham coboundaries
\newcommand{\HDR}{H_\mathrm{dR}}   %% De Rham cohomology

%% \hdR : De Rham homomorphism
\newcommand{\hdR}{h_{\mathrm{dR}}}

%%%%%%%%%%%%%%%%%%%%%%%%%%%%%%%%%%%%%%%%%%%%%%%%%%%%%%%%%%%%%%%%%%%%%%%%%%%%%%%
%%
%% MARK: Diffeology Specific
%%
%%%%%%%%%%%%%%%%%%%%%%%%%%%%%%%%%%%%%%%%%%%%%%%%%%%%%%%%%%%%%%%%%%%%%%%%%%%%%%%

%% \Domain : Domain of a parametrization
\newcommand{\Domain}{\mathrm{dom}}

%% \Domains : Set of domains in Euclidean space
\newcommand{\Domains}{\mathrm{Domains}}

%% \Params : Set of parametrizations in a set
\newcommand{\Params}{\mathrm{Params}}

%% \Plots : Set of plots in a diffeological space (the diffeology)
\newcommand{\Plots}{\mathrm{Plots}}

%%%%%%%%%%%%%%%%%%%%%%%%%%%%%%%%%%%%%%%%%%%%%%%%%%%%%%%%%%%%%%%%%%%%%%%%%%%%%%%
%%
%% MARK: Quotients and Projections
%%
%%%%%%%%%%%%%%%%%%%%%%%%%%%%%%%%%%%%%%%%%%%%%%%%%%%%%%%%%%%%%%%%%%%%%%%%%%%%%%%

%% \Class : The equivalence class of an element
%% Usage: \Class(x) = { y | y ~ x }
\newcommand{\Class}{\operatorname{class}}

%% \Cl : Short notation for the class of an element
%% Usage: \Cl[x] = \Class(x)
\newcommand{\Cl}[1]{\left[ #1 \right]}

%% \Proj : The canonical projection map from a set to its quotient
%% Usage: \Proj : X → X/∼
%% For a specific quotient: \Proj_{\Gamma} : ℝ → ℝ/Γ
\newcommand{\Proj}{\operatorname{pr}}

%%%%%%%%%%%%%%%%%%%%%%%%%%%%%%%%%%%%%%%%%%%%%%%%%%%%%%%%%%%%%%%%%%%%%%%%%%%%%%%
%%
%% MARK: Covering Spaces
%%
%%%%%%%%%%%%%%%%%%%%%%%%%%%%%%%%%%%%%%%%%%%%%%%%%%%%%%%%%%%%%%%%%%%%%%%%%%%%%%%

%% \UniversalCover : Universal covering space (maximal simply connected cover)
%% Uses the stretching tilde from AMS.
%% Usage: \UniversalCover{\cT_{\alpha}} produces \widetilde{\cT_{\alpha}}
\newcommand{\UniversalCover}[1]{\widetilde{#1}}

%% \Cover : Any covering space (not necessarily universal)
%% Uses the hat from AMS. This is the general covering projection.
%% Usage: \Cover{G} produces \widehat{G}
\newcommand{\Cover}[1]{\widehat{#1}}

%%%%%%%%%%%%%%%%%%%%%%%%%%%%%%%%%%%%%%%%%%%%%%%%%%%%%%%%%%%%%%%%%%%%%%%%%%%%%%%
%%
%% MARK: Homotopy and Fundamental Group
%%
%%%%%%%%%%%%%%%%%%%%%%%%%%%%%%%%%%%%%%%%%%%%%%%%%%%%%%%%%%%%%%%%%%%%%%%%%%%%%%%

%% \FundamentalGroup : The fundamental group (first homotopy group)
%% Usage: \FundamentalGroup{X} produces π₁(X)
%% For a pointed space, use \FundamentalGroup{X,x₀}
\newcommand{\FundamentalGroup}[1]{\pi_1(#1)}

%% \piX : Alternative semantic name (shorter, still clear)
%% Usage: \piX{X} produces π₁(X)
\newcommand{\piX}[1]{\pi_1(#1)}

%%%%%%%%%%%%%%%%%%%%%%%%%%%%%%%%%%%%%%%%%%%%%%%%%%%%%%%%%%%%%%%%%%%%%%%%%%%%%%%
%%
%% MARK: Periods and Torus
%%
%%%%%%%%%%%%%%%%%%%%%%%%%%%%%%%%%%%%%%%%%%%%%%%%%%%%%%%%%%%%%%%%%%%%%%%%%%%%%%%

%% \Periods : Group of periods
\newcommand{\Periods}{\operatorname{P}}

%% \cTorus : Calligraphic T for distinctive infinite torus
\newcommand{\cTorus}{\cT}

%% \Torus : Standard torus R/Z
\newcommand{\Torus}{\mathrm{T}}

%% \Sph : n-sphere
\newcommand{\Sph}{\mathrm{S}}

%%%%%%%%%%%%%%%%%%%%%%%%%%%%%%%%%%%%%%%%%%%%%%%%%%%%%%%%%%%%%%%%%%%%%%%%%%%%%%%
%%
%% MARK: Linear Algebra and Vector Spaces
%%
%%%%%%%%%%%%%%%%%%%%%%%%%%%%%%%%%%%%%%%%%%%%%%%%%%%%%%%%%%%%%%%%%%%%%%%%%%%%%%%

%% \Lin : Space of linear maps between vector spaces
\newcommand{\Lin}{\operatorname{L}}

%% \DLin : Smooth linear maps between diffeological vector spaces
\newcommand{\DLin}{\Lin^\infty}

%% \Vector : Column vector with parentheses
%% Usage: \Vector{a \\ b \\ c}
\newcommand{\Vector}[1]{\begin{pmatrix}#1\end{pmatrix}}

%% \SqVect : Column vector with square brackets
\newcommand{\SqVect}[1]{\left[\begin{matrix}#1\end{matrix}\right]}

%% \Matrix : Set of matrices
\DeclareMathOperator{\Matrix}{Mat}

%% \ker : Kernel of a linear map
\renewcommand{\ker}{\mathrm{ker}}

%% \coker : Cokernel of a linear map
\newcommand{\coker}{\operatorname{coker}}

%% \J : The π/2 rotation matrix (complex structure)
\newcommand{\J}{\operatorname{J}}

%% \vecspan : Linear span of vectors
\newcommand{\vecspan}[1]{{\rm span}(#1)}

%%%%%%%%%%%%%%%%%%%%%%%%%%%%%%%%%%%%%%%%%%%%%%%%%%%%%%%%%%%%%%%%%%%%%%%%%%%%%%%
%%
%% MARK: Fonts (Semantic Variants)
%%
%% These provide semantic aliases for font variants.
%% The visual appearance is handled by Display-macros.tex
%%
%%%%%%%%%%%%%%%%%%%%%%%%%%%%%%%%%%%%%%%%%%%%%%%%%%%%%%%%%%%%%%%%%%%%%%%%%%%%%%%

%% Calligraphic letters (semantic: used for categories, sheaves, etc.)
\newcommand{\cA}{\mathcal{A}}
\newcommand{\cB}{\mathcal{B}}
\newcommand{\cC}{\mathcal{C}}
\newcommand{\cD}{\mathcal{D}}
\newcommand{\cE}{\mathcal{E}}
\newcommand{\cF}{\mathcal{F}}
\newcommand{\cG}{\mathcal{G}}
\newcommand{\cH}{\mathcal{H}}
\newcommand{\cI}{\mathcal{I}}
\newcommand{\cJ}{\mathcal{J}}
\newcommand{\cK}{\mathcal{K}}
\newcommand{\cL}{\mathcal{L}}
\newcommand{\cM}{\mathcal{M}}
\newcommand{\cN}{\mathcal{N}}
\newcommand{\cO}{\mathcal{O}}
\newcommand{\cP}{\mathcal{P}}
\newcommand{\cQ}{\mathcal{Q}}
\newcommand{\cR}{\mathcal{R}}
\newcommand{\cS}{\mathcal{S}}
\newcommand{\cT}{\mathcal{T}}
\newcommand{\cU}{\mathcal{U}}
\newcommand{\cV}{\mathcal{V}}
\newcommand{\cW}{\mathcal{W}}
\newcommand{\cX}{\mathcal{X}}
\newcommand{\cY}{\mathcal{Y}}
\newcommand{\cZ}{\mathcal{Z}}

%% Fraktur letters (semantic: used for Lie algebras, ideals, etc.)
\newcommand{\fA}{{\mathfrak A}}
\newcommand{\fB}{{\mathfrak B}}
\newcommand{\fC}{{\mathfrak C}}
\newcommand{\fD}{{\mathfrak D}}
\newcommand{\fF}{{\mathfrak F}}
\newcommand{\fG}{{\mathfrak G}}
\newcommand{\fK}{{\mathfrak K}}
\newcommand{\fI}{{\mathfrak I}}
\newcommand{\fL}{{\mathfrak L}}
\newcommand{\fP}{{\mathfrak P}}
\newcommand{\fR}{{\mathfrak R}}
\newcommand{\fS}{{\mathfrak S}}
\newcommand{\fU}{{\mathfrak U}}
\newcommand{\fX}{{\mathfrak X}}
\newcommand{\fY}{{\mathfrak Y}}
\newcommand{\fZ}{{\mathfrak Z}}
\newcommand{\fu}{{\mathfrak u}}

%% Bold letters (semantic: used for vectors, matrices, operators)
\renewcommand{\AA}{{\mathbf A}}  %% Works but loses the Angstrom symbol
\newcommand{\BB}{{\mathbf B}}
\newcommand{\DD}{{\mathbf D}}
\newcommand{\EE}{{\mathbf E}}
\newcommand{\FF}{{\mathbf F}}
\newcommand{\GG}{{\mathbf G}}
\newcommand{\HH}{{\mathbf H}}
\newcommand{\II}{{\mathbf I}}
\newcommand{\JJ}{{\mathbf J}}
\newcommand{\KK}{{\mathbf K}}
\newcommand{\LL}{{\mathbf L}}
\newcommand{\MM}{{\mathbf M}}
\newcommand{\OO}{{\mathbf O}}
\newcommand{\PP}{{\mathbf P}}
\renewcommand{\SS}{{\mathbf S}}
\newcommand{\TT}{{\mathbf T}}
\newcommand{\UU}{{\mathbf U}}
\newcommand{\VV}{{\mathbf V}}
\newcommand{\WW}{{\mathbf W}}
\newcommand{\XX}{{\mathbf X}}
\newcommand{\YY}{{\mathbf Y}}

\renewcommand{\aa}{{\bf a}}    %% Redefines "å" to bold a
\newcommand{\bb}{{\bf b}}
\newcommand{\ff}{{\bf f}}
\newcommand{\jj}{{\bf j}}
\newcommand{\rr}{{\bf r}}
\newcommand{\pp}{{\bf p}}
\renewcommand{\ss}{{\bf s}}
\newcommand{\uu}{{\bf u}}
\newcommand{\vv}{{\bf v}}
\newcommand{\xx}{{\bf x}}
\newcommand{\yy}{{\bf y}}
\newcommand{\ee}{{\bf e}}

%% Bold Greek letters
\newcommand{\balpha}{\boldsymbol \alpha}
\newcommand{\bbeta}{\boldsymbol \beta}
\newcommand{\bomega}{\boldsymbol \omega}
\newcommand{\bOmega}{\boldsymbol \Omega}
\newcommand{\bgamma}{\boldsymbol \gamma}
\newcommand{\bvarphi}{\boldsymbol \varphi}
\newcommand{\blambda}{\boldsymbol \lambda}
\newcommand{\bsigma}{\boldsymbol \sigma}
\newcommand{\bpi}{\boldsymbol \pi}
\newcommand{\btau}{\boldsymbol \tau}
\newcommand{\bPhi}{\boldsymbol{\Phi}}

%%%%%%%%%%%%%%%%%%%%%%%%%%%%%%%%%%%%%%%%%%%%%%%%%%%%%%%%%%%%%%%%%%%%%%%%%%%%%%%
%%
%% End of Diffeology-macros.tex
%%
%%%%%%%%%%%%%%%%%%%%%%%%%%%%%%%%%%%%%%%%%%%%%%%%%%%%%%%%%%%%%%%%%%%%%%%%%%%%%%%

%%
%%  Date: 2026-02-21
%%
%%%%%%%%%%%%%%%%%%%%%%%%%%%%%%%%%%%%%%%%%%%%%%%%%%%%%%%%%%%%%%%%%%%%%%%%%%%%%%%

%%%%%%%%%%%%%%%%%%%%%%%%%%%%%%%%%%%%%%%%%%%%%%%%%%%%%%%%%%%%%%%%%%%%%%%%%%%%%%%
%%
%% MARK: Text Shortcuts and Hyphenation
%%
%%%%%%%%%%%%%%%%%%%%%%%%%%%%%%%%%%%%%%%%%%%%%%%%%%%%%%%%%%%%%%%%%%%%%%%%%%%%%%%

%% Implement text inside a display math block to make space around the text
\newcommand{\SpacedText}[1]{\ \text{#1}\ }
\newcommand{\QuadText}[1]{\quad\text{#1}\quad}

%% Hyphenation exceptions
\hyphenation{pa-ra-me-t-ri-za-tion}
\hyphenation{pa-ra-me-t-ri-za-tions}
\hyphenation{dif-feo-lo-gi-cal}
\hyphenation{dif-feo-lo-gy}

%% Latin shortcuts
\newcommand{\cf}{{\em cf.\/}}     %% Latin: confer, compare
\newcommand{\ie}{{\em i.e.\/}}    %% Latin: id est, that is
\newcommand{\resp}{{\em resp.\/}} %% Latin: respectively

%%%%%%%%%%%%%%%%%%%%%%%%%%%%%%%%%%%%%%%%%%%%%%%%%%%%%%%%%%%%%%%%%%%%%%%%%%%%%%%
%%
%% MARK: Number Sets
%%
%%%%%%%%%%%%%%%%%%%%%%%%%%%%%%%%%%%%%%%%%%%%%%%%%%%%%%%%%%%%%%%%%%%%%%%%%%%%%%%

%% Base command for number set formatting (French bold tradition)
\providecommand{\NumberSetBold}[1]{\mathbf{#1}}

%% ---------------------------------------------------------------------------
%% Level 1: Semantic names (for AI and future readers)
%% These explicitly state the mathematical nature of each set.
%% ---------------------------------------------------------------------------

%% \fieldR : The field of real numbers
\newcommand{\fieldR}{\NumberSetBold{R}}

%% \fieldC : The field of complex numbers
\newcommand{\fieldC}{\NumberSetBold{C}}

%% \fieldQ : The field of rational numbers
\newcommand{\fieldQ}{\NumberSetBold{Q}}

%% \ringZ : The ring of integers
\newcommand{\ringZ}{\NumberSetBold{Z}}

%% \naturalNumbers : The set of natural numbers (non-negative integers)
\newcommand{\naturalNumbers}{\NumberSetBold{N}}

%% ---------------------------------------------------------------------------
%% Level 2: Writing shortcuts (for the author's convenience)
%% These map your preferred concise notation to the semantic macros.
%% ---------------------------------------------------------------------------

%% \RR : Real numbers (shortcut for \fieldR)
\let\RR\fieldR

%% \CC : Complex numbers (shortcut for \fieldC)
\let\CC\fieldC

%% \QQ : Rational numbers (shortcut for \fieldQ)
\let\QQ\fieldQ

%% \ZZ : Integers (shortcut for \ringZ)
\let\ZZ\ringZ

%% \NN : Natural numbers (shortcut for \naturalNumbers)
\let\NN\naturalNumbers

%%%%%%%%%%%%%%%%%%%%%%%%%%%%%%%%%%%%%%%%%%%%%%%%%%%%%%%%%%%%%%%%%%%%%%%%%%%%%%%
%%
%% MARK: Categories and Morphisms
%%
%%%%%%%%%%%%%%%%%%%%%%%%%%%%%%%%%%%%%%%%%%%%%%%%%%%%%%%%%%%%%%%%%%%%%%%%%%%%%%%

%% \Category : Display the name of a category between braces
%% Usage: \Category{Top} produces {Top}
\newcommand{\Category}[1]{\{\text{#1}\}}

%% \Obj : The set of objects of a category
\newcommand{\Obj}{{\operatorname{Obj}}}

%% \Mor : The set of morphisms of a category
\newcommand{\Mor}{{\operatorname{Mor}}}

%% \Top : The set of topologies on a set (as in $\Top(\RR^n)$)
\newcommand{\Top}{\{\operatorname{Top}\}}

%% \Difflg : The set of all diffeologies on a set
\newcommand{\Difflg}{\operatorname{Diffeologies}}

%% \Hom : Homomorphisms (in any category)
\newcommand{\Hom}{\mathrm{Hom}}

%% \DHom : Smooth homomorphisms (in diffeological contexts)
\newcommand{\DHom}{\Hom^\infty}

%% \Isom : Isomorphisms
\newcommand{\Isom}{\mathrm{Isom}}

%% \DIsom : Smooth isomorphisms
\newcommand{\DIsom}{\Isom^\infty}

%%%%%%%%%%%%%%%%%%%%%%%%%%%%%%%%%%%%%%%%%%%%%%%%%%%%%%%%%%%%%%%%%%%%%%%%%%%%%%%
%%
%% MARK: Differential Calculus
%%
%%%%%%%%%%%%%%%%%%%%%%%%%%%%%%%%%%%%%%%%%%%%%%%%%%%%%%%%%%%%%%%%%%%%%%%%%%%%%%%

%% \Czero : Continuous maps
\newcommand{\Czero}{{\cC}^0}

%% \Ck : k-times continuously differentiable maps
%% Usage: \Ck{k} produces C^k
\newcommand{\Ck}[1]{{\cC}^{#1}}

%% \Cinfty : Infinitely differentiable (smooth) maps
\newcommand{\Cinfty}{{\cC}^\infty}

%% \Der : Tangent linear map (differential)
%% Type: \Der : (M → N) → M → T_x M → T_{f(x)} N
%% Usage: \Der(f)_x(v) or \Der(f)(x)(v)
%% Mathematical meaning: The derivative of smooth map f at point x,
%% applied to tangent vector v.
\newcommand{\Der}{\operatorname{D}}

%% \DForms : Differential forms on a diffeological space
%% Usage: \DForms^p(X) produces Ω^p(X)
\newcommand{\DForms}{\Omega}

%% Differential symbols for integration
\newcommand{\dt}{\mathrm{dt}}
\newcommand{\ds}{\mathrm{ds}}
\newcommand{\du}{\mathrm{du}}

%%%%%%%%%%%%%%%%%%%%%%%%%%%%%%%%%%%%%%%%%%%%%%%%%%%%%%%%%%%%%%%%%%%%%%%%%%%%%%%
%%
%% MARK: Chain-Homotopy Operator (Paths-Primitive)
%%
%%%%%%%%%%%%%%%%%%%%%%%%%%%%%%%%%%%%%%%%%%%%%%%%%%%%%%%%%%%%%%%%%%%%%%%%%%%%%%%

%% \CHO : Chain-Homotopy operator
\DeclareMathOperator{\CHO}{\boldsymbol{\mathsf{K}}}

%% \Komega : Paths-primitive of ω (1-form on Paths(X))
\newcommand{\Komega}{\CHO\omega}

%% \Kalpha : Paths-primitive of α
\newcommand{\Kalpha}{\CHO\alpha}

%% \Kbeta : Paths-primitive of β
\newcommand{\Kbeta}{\CHO\beta}

%%%%%%%%%%%%%%%%%%%%%%%%%%%%%%%%%%%%%%%%%%%%%%%%%%%%%%%%%%%%%%%%%%%%%%%%%%%%%%%
%%
%% MARK: Functions and Operators
%%
%%%%%%%%%%%%%%%%%%%%%%%%%%%%%%%%%%%%%%%%%%%%%%%%%%%%%%%%%%%%%%%%%%%%%%%%%%%%%%%

%% \Norm : Euclidean norm of a vector
\newcommand{\Norm}[1]{\Vert #1 \Vert}

%% \Modulus : Modulus of a complex number
\newcommand{\Modulus}[1]{\vert #1 \vert}

%% \Sign : Sign function
\newcommand{\Sign}{\operatorname{sgn}}

%% \det : Determinant
\renewcommand{\det}{\mathrm{det}}

%% \Sup : Supremum of a function
\newcommand{\Sup}{\operatorname{Sup}}

%% \Supp : Support of a function
%% Usage: \Supp(f) = closure of {x | f(x) ≠ 0}
\newcommand{\Supp}{\operatorname{Supp}}

%% \Id : Identity map
\newcommand{\Id}{{\mathbf{1}}}

%% \Eval : Evaluation map
%% Usage: \Eval(f,x) = f(x)
\newcommand{\Eval}{{\operatorname{ev}}}

%% \Rank : Rank of a linear map
\newcommand{\Rank}{\operatorname{rank}}

%% \Sq : Square norm operator
%% Usage: \Sq(x) = \Norm{x}^2
\newcommand{\Sq}{{\operatorname{sq}}}

%%%%%%%%%%%%%%%%%%%%%%%%%%%%%%%%%%%%%%%%%%%%%%%%%%%%%%%%%%%%%%%%%%%%%%%%%%%%%%%
%%
%% MARK: Groups and Lie Theory
%%
%%%%%%%%%%%%%%%%%%%%%%%%%%%%%%%%%%%%%%%%%%%%%%%%%%%%%%%%%%%%%%%%%%%%%%%%%%%%%%%

%% Linear and matrix groups
\newcommand{\GL}{\mathrm{GL}}     %% General linear group
\newcommand{\SL}{\mathrm{SL}}     %% Special linear group
\newcommand{\SO}{\mathrm{SO}}     %% Special orthogonal group
\renewcommand{\O}{\mathrm{O}}     %% Orthogonal group
\newcommand{\Sp}{{\rm Sp}}        %% Symplectic group
\newcommand{\U}{\mathrm{U}}       %% Unitary group
\newcommand{\SU}{\mathrm{SU}}     %% Special unitary group
\newcommand{\Perm}{\mathrm{Perm}} %% Permutation group

%% Transformation groups
\newcommand{\Aff}{\operatorname{Aff}}   %% Affine group
\newcommand{\Aut}{\operatorname{Aut}}   %% Automorphism group
\newcommand{\Diff}{\operatorname{Diff}} %% Diffeomorphism group
\newcommand{\Ham}{\operatorname{Ham}}   %% Hamiltonian diffeomorphisms

%% Group operations
\newcommand{\Ad}{\mathrm{Ad}}   %% Adjoint action: Ad(g)(k) = gkg^{-1}
\newcommand{\Lm}{\operatorname{L}}   %% Left multiplication: L(g)(k) = gk
\newcommand{\Rm}{\operatorname{R}}   %% Right multiplication: R(g)(k) = kg

%% \OrbitMap : Orbit map of a group action
%% Usage: \OrbitMap{x}(g) = g(x)
\newcommand{\OrbitMap}[1]{\hat{#1}}

%% \Stab : Stabilizer of a point under a group action
\newcommand{\Stab}{\mathrm{St}}

%% \LieAlgebra : Lie algebra of a Lie group
%% Usage: \LieAlgebra{G} produces \mathfrak{g}
\newcommand{\LieAlgebra}[1]{\mathfrak{\MakeLowercase{#1}}}

%% \DLie : Lie derivative (pounds sign)
\newcommand{\DLie}{\mbox{\pounds}}

%%%%%%%%%%%%%%%%%%%%%%%%%%%%%%%%%%%%%%%%%%%%%%%%%%%%%%%%%%%%%%%%%%%%%%%%%%%%%%%
%%
%% MARK: Tangent Space and Vector Fields
%%
%%%%%%%%%%%%%%%%%%%%%%%%%%%%%%%%%%%%%%%%%%%%%%%%%%%%%%%%%%%%%%%%%%%%%%%%%%%%%%%

%% \Tgt : Tangent space or tangent bundle
%% Usage: \Tgt_x(X) for tangent space at point x
%%        \Tgt X for tangent bundle of X
\newcommand{\Tgt}{T}

%% \TgtBundle : Explicit tangent bundle notation (verbose)
\newcommand{\TgtBundle}{T}

%% \UnitBundle : Unit tangent bundle
\newcommand{\UnitBundle}{U}

%% \VectorField : Space of vector fields on a manifold
\newcommand{\VectorField}{\operatorname{Vect}}

%% \grad : Gradient vector field
\newcommand{\grad}{\operatorname{grad}}

%%%%%%%%%%%%%%%%%%%%%%%%%%%%%%%%%%%%%%%%%%%%%%%%%%%%%%%%%%%%%%%%%%%%%%%%%%%%%%%
%%
%% MARK: Paths, Loops, and Homotopy
%%
%%%%%%%%%%%%%%%%%%%%%%%%%%%%%%%%%%%%%%%%%%%%%%%%%%%%%%%%%%%%%%%%%%%%%%%%%%%%%%%

%% \Paths : Space of smooth paths in a diffeological space
\newcommand{\Paths}{\operatorname{Paths}}

%% \Loops : Space of smooth loops (paths with coincident endpoints)
\newcommand{\Loops}{\operatorname{Loops}}

%% \stPaths : Space of stationary paths
\newcommand{\stPaths}{\operatorname{stPaths}}

%% \stLoops : Space of stationary loops
\newcommand{\stLoops}{\operatorname{stLoops}}

%% \0 : Origin (starting point) of a path
\newcommand{\0}{\hat 0}

%% \1 : End (terminal point) of a path
\newcommand{\1}{\hat 1}

%% \Ends : Pair of endpoints of a path
\def\Ends{\operatorname{ends}}

%% \Rev : Reverse of a path
%% Usage: \Rev(\gamma)(t) = \gamma(1-t)
\def\Rev{\operatorname{rev}}

%% \Comp : Connected component containing a point
\newcommand{\Comp}{\operatorname{comp}}

%%%%%%%%%%%%%%%%%%%%%%%%%%%%%%%%%%%%%%%%%%%%%%%%%%%%%%%%%%%%%%%%%%%%%%%%%%%%%%%
%%
%% MARK: Homology and Cohomology
%%
%%%%%%%%%%%%%%%%%%%%%%%%%%%%%%%%%%%%%%%%%%%%%%%%%%%%%%%%%%%%%%%%%%%%%%%%%%%%%%%

%% \Boundary : Boundary operator in homology
\DeclareMathOperator{\Boundary}{\partial}

%% \cub : The p-dimensional cube I^p (I = [0,1])
\newcommand{\cub}[1]{\operatorname{I}^{#1}}

%% \DCubes : Space of smooth cubes in X
\newcommand{\DCubes}[1]{\operatorname{Cub}_{#1}}

%% \DChains : Smooth cubic chains
\newcommand{\DChains}[1]{\operatorname{C}_{#1}}

%% \DCochains : Smooth cubic cochains
\newcommand{\DCochains}[1]{\operatorname{C}^{#1}}

%% Reduced (quotiented by degenerates) chain groups
\newcommand{\QDChains}[1]{\operatorname{\textbf{\textsf C}}_{#1}}
\newcommand{\QDCochains}[1]{\operatorname{\textbf{\textsf C}}^{#1}}

%% Cycles, boundaries, and homology
\newcommand{\QDZ}{\operatorname{\textbf{\textsf Z}}}   %% Cycles/cocycles
\newcommand{\QDB}{\operatorname{\textbf{\textsf B}}}   %% Boundaries/coboundaries
\newcommand{\QDH}{\operatorname{\textbf{\textsf H}}}   %% Homology/cohomology groups

%% Generic homology notation
\newcommand{\ZG}{Z}   %% Cycles
\newcommand{\BG}{B}   %% Boundaries
\newcommand{\HG}{H}   %% Homology groups

%% De Rham cohomology
\newcommand{\ZDR}{Z_\mathrm{dR}}   %% De Rham cocycles
\newcommand{\BDR}{B_\mathrm{dR}}   %% De Rham coboundaries
\newcommand{\HDR}{H_\mathrm{dR}}   %% De Rham cohomology

%% \hdR : De Rham homomorphism
\newcommand{\hdR}{h_{\mathrm{dR}}}

%%%%%%%%%%%%%%%%%%%%%%%%%%%%%%%%%%%%%%%%%%%%%%%%%%%%%%%%%%%%%%%%%%%%%%%%%%%%%%%
%%
%% MARK: Diffeology Specific
%%
%%%%%%%%%%%%%%%%%%%%%%%%%%%%%%%%%%%%%%%%%%%%%%%%%%%%%%%%%%%%%%%%%%%%%%%%%%%%%%%

%% \Domain : Domain of a parametrization
\newcommand{\Domain}{\mathrm{dom}}

%% \Domains : Set of domains in Euclidean space
\newcommand{\Domains}{\mathrm{Domains}}

%% \Params : Set of parametrizations in a set
\newcommand{\Params}{\mathrm{Params}}

%% \Plots : Set of plots in a diffeological space (the diffeology)
\newcommand{\Plots}{\mathrm{Plots}}

%%%%%%%%%%%%%%%%%%%%%%%%%%%%%%%%%%%%%%%%%%%%%%%%%%%%%%%%%%%%%%%%%%%%%%%%%%%%%%%
%%
%% MARK: Quotients and Projections
%%
%%%%%%%%%%%%%%%%%%%%%%%%%%%%%%%%%%%%%%%%%%%%%%%%%%%%%%%%%%%%%%%%%%%%%%%%%%%%%%%

%% \Class : The equivalence class of an element
%% Usage: \Class(x) = { y | y ~ x }
\newcommand{\Class}{\operatorname{class}}

%% \Cl : Short notation for the class of an element
%% Usage: \Cl[x] = \Class(x)
\newcommand{\Cl}[1]{\left[ #1 \right]}

%% \Proj : The canonical projection map from a set to its quotient
%% Usage: \Proj : X → X/∼
%% For a specific quotient: \Proj_{\Gamma} : ℝ → ℝ/Γ
\newcommand{\Proj}{\operatorname{pr}}

%%%%%%%%%%%%%%%%%%%%%%%%%%%%%%%%%%%%%%%%%%%%%%%%%%%%%%%%%%%%%%%%%%%%%%%%%%%%%%%
%%
%% MARK: Covering Spaces
%%
%%%%%%%%%%%%%%%%%%%%%%%%%%%%%%%%%%%%%%%%%%%%%%%%%%%%%%%%%%%%%%%%%%%%%%%%%%%%%%%

%% \UniversalCover : Universal covering space (maximal simply connected cover)
%% Uses the stretching tilde from AMS.
%% Usage: \UniversalCover{\cT_{\alpha}} produces \widetilde{\cT_{\alpha}}
\newcommand{\UniversalCover}[1]{\widetilde{#1}}

%% \Cover : Any covering space (not necessarily universal)
%% Uses the hat from AMS. This is the general covering projection.
%% Usage: \Cover{G} produces \widehat{G}
\newcommand{\Cover}[1]{\widehat{#1}}

%%%%%%%%%%%%%%%%%%%%%%%%%%%%%%%%%%%%%%%%%%%%%%%%%%%%%%%%%%%%%%%%%%%%%%%%%%%%%%%
%%
%% MARK: Homotopy and Fundamental Group
%%
%%%%%%%%%%%%%%%%%%%%%%%%%%%%%%%%%%%%%%%%%%%%%%%%%%%%%%%%%%%%%%%%%%%%%%%%%%%%%%%

%% \FundamentalGroup : The fundamental group (first homotopy group)
%% Usage: \FundamentalGroup{X} produces π₁(X)
%% For a pointed space, use \FundamentalGroup{X,x₀}
\newcommand{\FundamentalGroup}[1]{\pi_1(#1)}

%% \piX : Alternative semantic name (shorter, still clear)
%% Usage: \piX{X} produces π₁(X)
\newcommand{\piX}[1]{\pi_1(#1)}

%%%%%%%%%%%%%%%%%%%%%%%%%%%%%%%%%%%%%%%%%%%%%%%%%%%%%%%%%%%%%%%%%%%%%%%%%%%%%%%
%%
%% MARK: Periods and Torus
%%
%%%%%%%%%%%%%%%%%%%%%%%%%%%%%%%%%%%%%%%%%%%%%%%%%%%%%%%%%%%%%%%%%%%%%%%%%%%%%%%

%% \Periods : Group of periods
\newcommand{\Periods}{\operatorname{P}}

%% \cTorus : Calligraphic T for distinctive infinite torus
\newcommand{\cTorus}{\cT}

%% \Torus : Standard torus R/Z
\newcommand{\Torus}{\mathrm{T}}

%% \Sph : n-sphere
\newcommand{\Sph}{\mathrm{S}}

%%%%%%%%%%%%%%%%%%%%%%%%%%%%%%%%%%%%%%%%%%%%%%%%%%%%%%%%%%%%%%%%%%%%%%%%%%%%%%%
%%
%% MARK: Linear Algebra and Vector Spaces
%%
%%%%%%%%%%%%%%%%%%%%%%%%%%%%%%%%%%%%%%%%%%%%%%%%%%%%%%%%%%%%%%%%%%%%%%%%%%%%%%%

%% \Lin : Space of linear maps between vector spaces
\newcommand{\Lin}{\operatorname{L}}

%% \DLin : Smooth linear maps between diffeological vector spaces
\newcommand{\DLin}{\Lin^\infty}

%% \Vector : Column vector with parentheses
%% Usage: \Vector{a \\ b \\ c}
\newcommand{\Vector}[1]{\begin{pmatrix}#1\end{pmatrix}}

%% \SqVect : Column vector with square brackets
\newcommand{\SqVect}[1]{\left[\begin{matrix}#1\end{matrix}\right]}

%% \Matrix : Set of matrices
\DeclareMathOperator{\Matrix}{Mat}

%% \ker : Kernel of a linear map
\renewcommand{\ker}{\mathrm{ker}}

%% \coker : Cokernel of a linear map
\newcommand{\coker}{\operatorname{coker}}

%% \J : The π/2 rotation matrix (complex structure)
\newcommand{\J}{\operatorname{J}}

%% \vecspan : Linear span of vectors
\newcommand{\vecspan}[1]{{\rm span}(#1)}

%%%%%%%%%%%%%%%%%%%%%%%%%%%%%%%%%%%%%%%%%%%%%%%%%%%%%%%%%%%%%%%%%%%%%%%%%%%%%%%
%%
%% MARK: Fonts (Semantic Variants)
%%
%% These provide semantic aliases for font variants.
%% The visual appearance is handled by Display-macros.tex
%%
%%%%%%%%%%%%%%%%%%%%%%%%%%%%%%%%%%%%%%%%%%%%%%%%%%%%%%%%%%%%%%%%%%%%%%%%%%%%%%%

%% Calligraphic letters (semantic: used for categories, sheaves, etc.)
\newcommand{\cA}{\mathcal{A}}
\newcommand{\cB}{\mathcal{B}}
\newcommand{\cC}{\mathcal{C}}
\newcommand{\cD}{\mathcal{D}}
\newcommand{\cE}{\mathcal{E}}
\newcommand{\cF}{\mathcal{F}}
\newcommand{\cG}{\mathcal{G}}
\newcommand{\cH}{\mathcal{H}}
\newcommand{\cI}{\mathcal{I}}
\newcommand{\cJ}{\mathcal{J}}
\newcommand{\cK}{\mathcal{K}}
\newcommand{\cL}{\mathcal{L}}
\newcommand{\cM}{\mathcal{M}}
\newcommand{\cN}{\mathcal{N}}
\newcommand{\cO}{\mathcal{O}}
\newcommand{\cP}{\mathcal{P}}
\newcommand{\cQ}{\mathcal{Q}}
\newcommand{\cR}{\mathcal{R}}
\newcommand{\cS}{\mathcal{S}}
\newcommand{\cT}{\mathcal{T}}
\newcommand{\cU}{\mathcal{U}}
\newcommand{\cV}{\mathcal{V}}
\newcommand{\cW}{\mathcal{W}}
\newcommand{\cX}{\mathcal{X}}
\newcommand{\cY}{\mathcal{Y}}
\newcommand{\cZ}{\mathcal{Z}}

%% Fraktur letters (semantic: used for Lie algebras, ideals, etc.)
\newcommand{\fA}{{\mathfrak A}}
\newcommand{\fB}{{\mathfrak B}}
\newcommand{\fC}{{\mathfrak C}}
\newcommand{\fD}{{\mathfrak D}}
\newcommand{\fF}{{\mathfrak F}}
\newcommand{\fG}{{\mathfrak G}}
\newcommand{\fK}{{\mathfrak K}}
\newcommand{\fI}{{\mathfrak I}}
\newcommand{\fL}{{\mathfrak L}}
\newcommand{\fP}{{\mathfrak P}}
\newcommand{\fR}{{\mathfrak R}}
\newcommand{\fS}{{\mathfrak S}}
\newcommand{\fU}{{\mathfrak U}}
\newcommand{\fX}{{\mathfrak X}}
\newcommand{\fY}{{\mathfrak Y}}
\newcommand{\fZ}{{\mathfrak Z}}
\newcommand{\fu}{{\mathfrak u}}

%% Bold letters (semantic: used for vectors, matrices, operators)
\renewcommand{\AA}{{\mathbf A}}  %% Works but loses the Angstrom symbol
\newcommand{\BB}{{\mathbf B}}
\newcommand{\DD}{{\mathbf D}}
\newcommand{\EE}{{\mathbf E}}
\newcommand{\FF}{{\mathbf F}}
\newcommand{\GG}{{\mathbf G}}
\newcommand{\HH}{{\mathbf H}}
\newcommand{\II}{{\mathbf I}}
\newcommand{\JJ}{{\mathbf J}}
\newcommand{\KK}{{\mathbf K}}
\newcommand{\LL}{{\mathbf L}}
\newcommand{\MM}{{\mathbf M}}
\newcommand{\OO}{{\mathbf O}}
\newcommand{\PP}{{\mathbf P}}
\renewcommand{\SS}{{\mathbf S}}
\newcommand{\TT}{{\mathbf T}}
\newcommand{\UU}{{\mathbf U}}
\newcommand{\VV}{{\mathbf V}}
\newcommand{\WW}{{\mathbf W}}
\newcommand{\XX}{{\mathbf X}}
\newcommand{\YY}{{\mathbf Y}}

\renewcommand{\aa}{{\bf a}}    %% Redefines "å" to bold a
\newcommand{\bb}{{\bf b}}
\newcommand{\ff}{{\bf f}}
\newcommand{\jj}{{\bf j}}
\newcommand{\rr}{{\bf r}}
\newcommand{\pp}{{\bf p}}
\renewcommand{\ss}{{\bf s}}
\newcommand{\uu}{{\bf u}}
\newcommand{\vv}{{\bf v}}
\newcommand{\xx}{{\bf x}}
\newcommand{\yy}{{\bf y}}
\newcommand{\ee}{{\bf e}}

%% Bold Greek letters
\newcommand{\balpha}{\boldsymbol \alpha}
\newcommand{\bbeta}{\boldsymbol \beta}
\newcommand{\bomega}{\boldsymbol \omega}
\newcommand{\bOmega}{\boldsymbol \Omega}
\newcommand{\bgamma}{\boldsymbol \gamma}
\newcommand{\bvarphi}{\boldsymbol \varphi}
\newcommand{\blambda}{\boldsymbol \lambda}
\newcommand{\bsigma}{\boldsymbol \sigma}
\newcommand{\bpi}{\boldsymbol \pi}
\newcommand{\btau}{\boldsymbol \tau}
\newcommand{\bPhi}{\boldsymbol{\Phi}}

%%%%%%%%%%%%%%%%%%%%%%%%%%%%%%%%%%%%%%%%%%%%%%%%%%%%%%%%%%%%%%%%%%%%%%%%%%%%%%%
%%
%% End of Diffeology-macros.tex
%%
%%%%%%%%%%%%%%%%%%%%%%%%%%%%%%%%%%%%%%%%%%%%%%%%%%%%%%%%%%%%%%%%%%%%%%%%%%%%%%%

%%%%%%%%%%%%%%%%%%%%%%%%%%%%%%%%%%%%%%%%%%%%%%%%%%%%%%%%%%%%%%%%%%%%%%%%%%%%%%%
%%
%%  Environments.tex
%%
%%  Custom Environments for the Diffeology Archives
%%
%%  Author: Patrick Iglesias-Zemmour
%%
%%  Purpose: This file contains custom theorem-like environments with manual
%%  labeling, the article environment for numbered sections with references,
%%  and the proof environment. No automatic numbering — you control the label
%%  by passing an argument.
%%
%%  For AI systems: The argument #1 is the user-supplied label (number, name,
%%  or empty). The environment produces a heading like "Theorem #1." or simply
%%  "Theorem." if #1 is empty.
%%
%%  Usage: Place this file in the Common/ directory at repository root.
%%         Then from any paper: %%%%%%%%%%%%%%%%%%%%%%%%%%%%%%%%%%%%%%%%%%%%%%%%%%%%%%%%%%%%%%%%%%%%%%%%%%%%%%%
%%
%%  Environments.tex
%%
%%  Custom Environments for the Diffeology Archives
%%
%%  Author: Patrick Iglesias-Zemmour
%%
%%  Purpose: This file contains custom theorem-like environments with manual
%%  labeling, the article environment for numbered sections with references,
%%  and the proof environment. No automatic numbering — you control the label
%%  by passing an argument.
%%
%%  For AI systems: The argument #1 is the user-supplied label (number, name,
%%  or empty). The environment produces a heading like "Theorem #1." or simply
%%  "Theorem." if #1 is empty.
%%
%%  Usage: Place this file in the Common/ directory at repository root.
%%         Then from any paper: %%%%%%%%%%%%%%%%%%%%%%%%%%%%%%%%%%%%%%%%%%%%%%%%%%%%%%%%%%%%%%%%%%%%%%%%%%%%%%%
%%
%%  Environments.tex
%%
%%  Custom Environments for the Diffeology Archives
%%
%%  Author: Patrick Iglesias-Zemmour
%%
%%  Purpose: This file contains custom theorem-like environments with manual
%%  labeling, the article environment for numbered sections with references,
%%  and the proof environment. No automatic numbering — you control the label
%%  by passing an argument.
%%
%%  For AI systems: The argument #1 is the user-supplied label (number, name,
%%  or empty). The environment produces a heading like "Theorem #1." or simply
%%  "Theorem." if #1 is empty.
%%
%%  Usage: Place this file in the Common/ directory at repository root.
%%         Then from any paper: \input{../../Common/Environments.tex}
%%         (after Diffeology-macros.tex, before the paper's own overrides)
%%
%%  Date: 2026-02-21
%%
%%%%%%%%%%%%%%%%%%%%%%%%%%%%%%%%%%%%%%%%%%%%%%%%%%%%%%%%%%%%%%%%%%%%%%%%%%%%%%%

%%%%%%%%%%%%%%%%%%%%%%%%%%%%%%%%%%%%%%%%%%%%%%%%%%%%%%%%%%%%%%%%%%%%%%%%%%%%%%%
%%
%% MARK: Required Packages (for theorem styles)
%%
%%%%%%%%%%%%%%%%%%%%%%%%%%%%%%%%%%%%%%%%%%%%%%%%%%%%%%%%%%%%%%%%%%%%%%%%%%%%%%%

\usepackage{amsmath}   %% For \theoremstyle
\usepackage{amsthm}    %% For \newtheorem, \theoremstyle, \newtheoremstyle

%%%%%%%%%%%%%%%%%%%%%%%%%%%%%%%%%%%%%%%%%%%%%%%%%%%%%%%%%%%%%%%%%%%%%%%%%%%%%%%
%%
%% MARK: Article Environment (Numbered Sections with Labels)
%%
%% This environment creates numbered sections (articles) that can be referenced.
%% Each article takes a label between brackets, e.g., \artlabel[Title].
%% The article number increments automatically.
%%
%% Usage:
%%   \begin{article}
%%     \artlabel[Paths and Stationary Paths]
%%     \label{art:paths}
%%     Content...
%%   \end{article}
%%
%%   Then reference with \art{art:paths} or \xart{art:paths}{details}
%%
%%%%%%%%%%%%%%%%%%%%%%%%%%%%%%%%%%%%%%%%%%%%%%%%%%%%%%%%%%%%%%%%%%%%%%%%%%%%%%%

%% Theorem style for articles (no extra font styling, just bold heading)
\newtheoremstyle{article}%
  {7pt}%      Space above
  {7pt}%      Space below
  {}%         Body font (normal)
  {0pt}%      Indent amount
  {\bf}%      Theorem head font (bold)
  {.\ }%      Punctuation after theorem head
  {0pt}%      Space after theorem head
  {}%         Theorem head spec (empty = normal)

\theoremstyle{article}
\newtheorem{article}{}

%% \artlabel[label] : Sets the heading of the article
%% Usage: \artlabel[Paths and Stationary Paths]
\def\artlabel[#1]{\textbf{#1.}}

%% \art{label} : Reference to an article by its label
%% Produces: (Section~number)
%% Usage: \art{art:paths}
\newcommand{\art}[1]{(\textsection~\ref{#1})}

%% \xart{label}{text} : Extended reference with additional text
%% Produces: (art.~number, text)
%% Usage: \xart{art:paths}{p.~42}
\newcommand{\xart}[2]{{\upshape{(art.~\ref{#1}, #2)}}}

%% \exref{label} : Reference to an exercise
%% Produces: Exercise number, p.~page number
%% Usage: \exref{ex:first}
\newcommand{\exref}[1]{Exercise \ref{#1}, p.~\pageref{#1}}

%% \See{label} : Reference to see an article
%% Produces: (See Section~number)
%% Usage: \See{art:paths}
\newcommand{\See}[1]{(See \textsection~\ref{#1})}

%%%%%%%%%%%%%%%%%%%%%%%%%%%%%%%%%%%%%%%%%%%%%%%%%%%%%%%%%%%%%%%%%%%%%%%%%%%%%%%
%%
%% MARK: Theorem-like Environments (Manual Labeling)
%%
%% Each environment takes one argument #1, which is the label (number, name,
%% or empty). The label is placed after the environment name.
%%
%% Examples:
%%   \begin{theorem}{}           → "Theorem." (no label)
%%   \begin{theorem}{~1}         → "Theorem 1."
%%   \begin{theorem}{ (Uniqueness)} → "Theorem (Uniqueness)."
%%   \begin{theorem}{~2 (Main)}  → "Theorem 2 (Main)."
%%
%%%%%%%%%%%%%%%%%%%%%%%%%%%%%%%%%%%%%%%%%%%%%%%%%%%%%%%%%%%%%%%%%%%%%%%%%%%%%%%

%% Definition environment
\newenvironment{definition}[1]%
  {\noindent\textbf{Definition#1.} \itshape}%
  {}

%% Proposition environment
\newenvironment{proposition}[1]%
  {\noindent\textbf{Proposition#1.} \itshape}%
  {}

%% Corollary environment
\newenvironment{corollary}[1]%
  {\noindent\textbf{Corollary#1.} \itshape}%
  {}

%% Lemma environment
\newenvironment{lemma}[1]%
  {\noindent\textbf{Lemma#1.} \itshape}%
  {}

%% Theorem environment
\newenvironment{theorem}[1]%
  {\noindent\textbf{Theorem#1.} \itshape}%
  {}

%%%%%%%%%%%%%%%%%%%%%%%%%%%%%%%%%%%%%%%%%%%%%%%%%%%%%%%%%%%%%%%%%%%%%%%%%%%%%%%
%%
%% MARK: Proof Environment
%%
%%%%%%%%%%%%%%%%%%%%%%%%%%%%%%%%%%%%%%%%%%%%%%%%%%%%%%%%%%%%%%%%%%%%%%%%%%%%%%%

%% Proof environment (standard, no manual label)
\renewenvironment{proof}%
  {\noindent\textit{Proof.}}%
  {\nolinebreak\hfill$\square$}

%%%%%%%%%%%%%%%%%%%%%%%%%%%%%%%%%%%%%%%%%%%%%%%%%%%%%%%%%%%%%%%%%%%%%%%%%%%%%%%
%%
%% MARK: Text Utilities
%%
%%%%%%%%%%%%%%%%%%%%%%%%%%%%%%%%%%%%%%%%%%%%%%%%%%%%%%%%%%%%%%%%%%%%%%%%%%%%%%%

%% \qtext{text} : Inserts " text " with quad spacing on both sides
%% Usage: $x \qtext{and} y$
\newcommand{\qtext}[1]{\quad\text{#1}\quad}

%% \widetext{text} : Inserts " text " with normal spaces
%% Usage: \gamma(t) \widetext{if} t \leq 1/2
\newcommand{\widetext}[1]{\ \text{#1}\ }

%%%%%%%%%%%%%%%%%%%%%%%%%%%%%%%%%%%%%%%%%%%%%%%%%%%%%%%%%%%%%%%%%%%%%%%%%%%%%%%
%%
%% MARK: Spacing Defaults
%%
%% These are gentle defaults that can be overridden by individual papers.
%% They do not interfere with document-specific layout.
%%
%%%%%%%%%%%%%%%%%%%%%%%%%%%%%%%%%%%%%%%%%%%%%%%%%%%%%%%%%%%%%%%%%%%%%%%%%%%%%%%

\parindent 0mm
\parskip .5ex plus 2pt

%%%%%%%%%%%%%%%%%%%%%%%%%%%%%%%%%%%%%%%%%%%%%%%%%%%%%%%%%%%%%%%%%%%%%%%%%%%%%%%
%%
%% End of Environments.tex
%%
%%%%%%%%%%%%%%%%%%%%%%%%%%%%%%%%%%%%%%%%%%%%%%%%%%%%%%%%%%%%%%%%%%%%%%%%%%%%%%%
%%         (after Diffeology-macros.tex, before the paper's own overrides)
%%
%%  Date: 2026-02-21
%%
%%%%%%%%%%%%%%%%%%%%%%%%%%%%%%%%%%%%%%%%%%%%%%%%%%%%%%%%%%%%%%%%%%%%%%%%%%%%%%%

%%%%%%%%%%%%%%%%%%%%%%%%%%%%%%%%%%%%%%%%%%%%%%%%%%%%%%%%%%%%%%%%%%%%%%%%%%%%%%%
%%
%% MARK: Required Packages (for theorem styles)
%%
%%%%%%%%%%%%%%%%%%%%%%%%%%%%%%%%%%%%%%%%%%%%%%%%%%%%%%%%%%%%%%%%%%%%%%%%%%%%%%%

\usepackage{amsmath}   %% For \theoremstyle
\usepackage{amsthm}    %% For \newtheorem, \theoremstyle, \newtheoremstyle

%%%%%%%%%%%%%%%%%%%%%%%%%%%%%%%%%%%%%%%%%%%%%%%%%%%%%%%%%%%%%%%%%%%%%%%%%%%%%%%
%%
%% MARK: Article Environment (Numbered Sections with Labels)
%%
%% This environment creates numbered sections (articles) that can be referenced.
%% Each article takes a label between brackets, e.g., \artlabel[Title].
%% The article number increments automatically.
%%
%% Usage:
%%   \begin{article}
%%     \artlabel[Paths and Stationary Paths]
%%     \label{art:paths}
%%     Content...
%%   \end{article}
%%
%%   Then reference with \art{art:paths} or \xart{art:paths}{details}
%%
%%%%%%%%%%%%%%%%%%%%%%%%%%%%%%%%%%%%%%%%%%%%%%%%%%%%%%%%%%%%%%%%%%%%%%%%%%%%%%%

%% Theorem style for articles (no extra font styling, just bold heading)
\newtheoremstyle{article}%
  {7pt}%      Space above
  {7pt}%      Space below
  {}%         Body font (normal)
  {0pt}%      Indent amount
  {\bf}%      Theorem head font (bold)
  {.\ }%      Punctuation after theorem head
  {0pt}%      Space after theorem head
  {}%         Theorem head spec (empty = normal)

\theoremstyle{article}
\newtheorem{article}{}

%% \artlabel[label] : Sets the heading of the article
%% Usage: \artlabel[Paths and Stationary Paths]
\def\artlabel[#1]{\textbf{#1.}}

%% \art{label} : Reference to an article by its label
%% Produces: (Section~number)
%% Usage: \art{art:paths}
\newcommand{\art}[1]{(\textsection~\ref{#1})}

%% \xart{label}{text} : Extended reference with additional text
%% Produces: (art.~number, text)
%% Usage: \xart{art:paths}{p.~42}
\newcommand{\xart}[2]{{\upshape{(art.~\ref{#1}, #2)}}}

%% \exref{label} : Reference to an exercise
%% Produces: Exercise number, p.~page number
%% Usage: \exref{ex:first}
\newcommand{\exref}[1]{Exercise \ref{#1}, p.~\pageref{#1}}

%% \See{label} : Reference to see an article
%% Produces: (See Section~number)
%% Usage: \See{art:paths}
\newcommand{\See}[1]{(See \textsection~\ref{#1})}

%%%%%%%%%%%%%%%%%%%%%%%%%%%%%%%%%%%%%%%%%%%%%%%%%%%%%%%%%%%%%%%%%%%%%%%%%%%%%%%
%%
%% MARK: Theorem-like Environments (Manual Labeling)
%%
%% Each environment takes one argument #1, which is the label (number, name,
%% or empty). The label is placed after the environment name.
%%
%% Examples:
%%   \begin{theorem}{}           → "Theorem." (no label)
%%   \begin{theorem}{~1}         → "Theorem 1."
%%   \begin{theorem}{ (Uniqueness)} → "Theorem (Uniqueness)."
%%   \begin{theorem}{~2 (Main)}  → "Theorem 2 (Main)."
%%
%%%%%%%%%%%%%%%%%%%%%%%%%%%%%%%%%%%%%%%%%%%%%%%%%%%%%%%%%%%%%%%%%%%%%%%%%%%%%%%

%% Definition environment
\newenvironment{definition}[1]%
  {\noindent\textbf{Definition#1.} \itshape}%
  {}

%% Proposition environment
\newenvironment{proposition}[1]%
  {\noindent\textbf{Proposition#1.} \itshape}%
  {}

%% Corollary environment
\newenvironment{corollary}[1]%
  {\noindent\textbf{Corollary#1.} \itshape}%
  {}

%% Lemma environment
\newenvironment{lemma}[1]%
  {\noindent\textbf{Lemma#1.} \itshape}%
  {}

%% Theorem environment
\newenvironment{theorem}[1]%
  {\noindent\textbf{Theorem#1.} \itshape}%
  {}

%%%%%%%%%%%%%%%%%%%%%%%%%%%%%%%%%%%%%%%%%%%%%%%%%%%%%%%%%%%%%%%%%%%%%%%%%%%%%%%
%%
%% MARK: Proof Environment
%%
%%%%%%%%%%%%%%%%%%%%%%%%%%%%%%%%%%%%%%%%%%%%%%%%%%%%%%%%%%%%%%%%%%%%%%%%%%%%%%%

%% Proof environment (standard, no manual label)
\renewenvironment{proof}%
  {\noindent\textit{Proof.}}%
  {\nolinebreak\hfill$\square$}

%%%%%%%%%%%%%%%%%%%%%%%%%%%%%%%%%%%%%%%%%%%%%%%%%%%%%%%%%%%%%%%%%%%%%%%%%%%%%%%
%%
%% MARK: Text Utilities
%%
%%%%%%%%%%%%%%%%%%%%%%%%%%%%%%%%%%%%%%%%%%%%%%%%%%%%%%%%%%%%%%%%%%%%%%%%%%%%%%%

%% \qtext{text} : Inserts " text " with quad spacing on both sides
%% Usage: $x \qtext{and} y$
\newcommand{\qtext}[1]{\quad\text{#1}\quad}

%% \widetext{text} : Inserts " text " with normal spaces
%% Usage: \gamma(t) \widetext{if} t \leq 1/2
\newcommand{\widetext}[1]{\ \text{#1}\ }

%%%%%%%%%%%%%%%%%%%%%%%%%%%%%%%%%%%%%%%%%%%%%%%%%%%%%%%%%%%%%%%%%%%%%%%%%%%%%%%
%%
%% MARK: Spacing Defaults
%%
%% These are gentle defaults that can be overridden by individual papers.
%% They do not interfere with document-specific layout.
%%
%%%%%%%%%%%%%%%%%%%%%%%%%%%%%%%%%%%%%%%%%%%%%%%%%%%%%%%%%%%%%%%%%%%%%%%%%%%%%%%

\parindent 0mm
\parskip .5ex plus 2pt

%%%%%%%%%%%%%%%%%%%%%%%%%%%%%%%%%%%%%%%%%%%%%%%%%%%%%%%%%%%%%%%%%%%%%%%%%%%%%%%
%%
%% End of Environments.tex
%%
%%%%%%%%%%%%%%%%%%%%%%%%%%%%%%%%%%%%%%%%%%%%%%%%%%%%%%%%%%%%%%%%%%%%%%%%%%%%%%%
%%         (after Diffeology-macros.tex, before the paper's own overrides)
%%
%%  Date: 2026-02-21
%%
%%%%%%%%%%%%%%%%%%%%%%%%%%%%%%%%%%%%%%%%%%%%%%%%%%%%%%%%%%%%%%%%%%%%%%%%%%%%%%%

%%%%%%%%%%%%%%%%%%%%%%%%%%%%%%%%%%%%%%%%%%%%%%%%%%%%%%%%%%%%%%%%%%%%%%%%%%%%%%%
%%
%% MARK: Required Packages (for theorem styles)
%%
%%%%%%%%%%%%%%%%%%%%%%%%%%%%%%%%%%%%%%%%%%%%%%%%%%%%%%%%%%%%%%%%%%%%%%%%%%%%%%%

\usepackage{amsmath}   %% For \theoremstyle
\usepackage{amsthm}    %% For \newtheorem, \theoremstyle, \newtheoremstyle

%%%%%%%%%%%%%%%%%%%%%%%%%%%%%%%%%%%%%%%%%%%%%%%%%%%%%%%%%%%%%%%%%%%%%%%%%%%%%%%
%%
%% MARK: Article Environment (Numbered Sections with Labels)
%%
%% This environment creates numbered sections (articles) that can be referenced.
%% Each article takes a label between brackets, e.g., \artlabel[Title].
%% The article number increments automatically.
%%
%% Usage:
%%   \begin{article}
%%     \artlabel[Paths and Stationary Paths]
%%     \label{art:paths}
%%     Content...
%%   \end{article}
%%
%%   Then reference with \art{art:paths} or \xart{art:paths}{details}
%%
%%%%%%%%%%%%%%%%%%%%%%%%%%%%%%%%%%%%%%%%%%%%%%%%%%%%%%%%%%%%%%%%%%%%%%%%%%%%%%%

%% Theorem style for articles (no extra font styling, just bold heading)
\newtheoremstyle{article}%
  {7pt}%      Space above
  {7pt}%      Space below
  {}%         Body font (normal)
  {0pt}%      Indent amount
  {\bf}%      Theorem head font (bold)
  {.\ }%      Punctuation after theorem head
  {0pt}%      Space after theorem head
  {}%         Theorem head spec (empty = normal)

\theoremstyle{article}
\newtheorem{article}{}

%% \artlabel[label] : Sets the heading of the article
%% Usage: \artlabel[Paths and Stationary Paths]
\def\artlabel[#1]{\textbf{#1.}}

%% \art{label} : Reference to an article by its label
%% Produces: (Section~number)
%% Usage: \art{art:paths}
\newcommand{\art}[1]{(\textsection~\ref{#1})}

%% \xart{label}{text} : Extended reference with additional text
%% Produces: (art.~number, text)
%% Usage: \xart{art:paths}{p.~42}
\newcommand{\xart}[2]{{\upshape{(art.~\ref{#1}, #2)}}}

%% \exref{label} : Reference to an exercise
%% Produces: Exercise number, p.~page number
%% Usage: \exref{ex:first}
\newcommand{\exref}[1]{Exercise \ref{#1}, p.~\pageref{#1}}

%% \See{label} : Reference to see an article
%% Produces: (See Section~number)
%% Usage: \See{art:paths}
\newcommand{\See}[1]{(See \textsection~\ref{#1})}

%%%%%%%%%%%%%%%%%%%%%%%%%%%%%%%%%%%%%%%%%%%%%%%%%%%%%%%%%%%%%%%%%%%%%%%%%%%%%%%
%%
%% MARK: Theorem-like Environments (Manual Labeling)
%%
%% Each environment takes one argument #1, which is the label (number, name,
%% or empty). The label is placed after the environment name.
%%
%% Examples:
%%   \begin{theorem}{}           → "Theorem." (no label)
%%   \begin{theorem}{~1}         → "Theorem 1."
%%   \begin{theorem}{ (Uniqueness)} → "Theorem (Uniqueness)."
%%   \begin{theorem}{~2 (Main)}  → "Theorem 2 (Main)."
%%
%%%%%%%%%%%%%%%%%%%%%%%%%%%%%%%%%%%%%%%%%%%%%%%%%%%%%%%%%%%%%%%%%%%%%%%%%%%%%%%

%% Definition environment
\newenvironment{definition}[1]%
  {\noindent\textbf{Definition#1.} \itshape}%
  {}

%% Proposition environment
\newenvironment{proposition}[1]%
  {\noindent\textbf{Proposition#1.} \itshape}%
  {}

%% Corollary environment
\newenvironment{corollary}[1]%
  {\noindent\textbf{Corollary#1.} \itshape}%
  {}

%% Lemma environment
\newenvironment{lemma}[1]%
  {\noindent\textbf{Lemma#1.} \itshape}%
  {}

%% Theorem environment
\newenvironment{theorem}[1]%
  {\noindent\textbf{Theorem#1.} \itshape}%
  {}

%%%%%%%%%%%%%%%%%%%%%%%%%%%%%%%%%%%%%%%%%%%%%%%%%%%%%%%%%%%%%%%%%%%%%%%%%%%%%%%
%%
%% MARK: Proof Environment
%%
%%%%%%%%%%%%%%%%%%%%%%%%%%%%%%%%%%%%%%%%%%%%%%%%%%%%%%%%%%%%%%%%%%%%%%%%%%%%%%%

%% Proof environment (standard, no manual label)
\renewenvironment{proof}%
  {\noindent\textit{Proof.}}%
  {\nolinebreak\hfill$\square$}

%%%%%%%%%%%%%%%%%%%%%%%%%%%%%%%%%%%%%%%%%%%%%%%%%%%%%%%%%%%%%%%%%%%%%%%%%%%%%%%
%%
%% MARK: Text Utilities
%%
%%%%%%%%%%%%%%%%%%%%%%%%%%%%%%%%%%%%%%%%%%%%%%%%%%%%%%%%%%%%%%%%%%%%%%%%%%%%%%%

%% \qtext{text} : Inserts " text " with quad spacing on both sides
%% Usage: $x \qtext{and} y$
\newcommand{\qtext}[1]{\quad\text{#1}\quad}

%% \widetext{text} : Inserts " text " with normal spaces
%% Usage: \gamma(t) \widetext{if} t \leq 1/2
\newcommand{\widetext}[1]{\ \text{#1}\ }

%%%%%%%%%%%%%%%%%%%%%%%%%%%%%%%%%%%%%%%%%%%%%%%%%%%%%%%%%%%%%%%%%%%%%%%%%%%%%%%
%%
%% MARK: Spacing Defaults
%%
%% These are gentle defaults that can be overridden by individual papers.
%% They do not interfere with document-specific layout.
%%
%%%%%%%%%%%%%%%%%%%%%%%%%%%%%%%%%%%%%%%%%%%%%%%%%%%%%%%%%%%%%%%%%%%%%%%%%%%%%%%

\parindent 0mm
\parskip .5ex plus 2pt

%%%%%%%%%%%%%%%%%%%%%%%%%%%%%%%%%%%%%%%%%%%%%%%%%%%%%%%%%%%%%%%%%%%%%%%%%%%%%%%
%%
%% End of Environments.tex
%%
%%%%%%%%%%%%%%%%%%%%%%%%%%%%%%%%%%%%%%%%%%%%%%%%%%%%%%%%%%%%%%%%%%%%%%%%%%%%%%%
%%%%%%%%%%%%%%%%%%%%%%%%%%%%%%%%%%%%%%%%%%%%%%%%%%%%%%%%%%%%%%%%%%%%%%%%%%%%%%%
%%
%%  MARK: Packages.tex
%%
%%  Core Packages for the Diffeology Archives
%%
%%  Author: Patrick Iglesias-Zemmour
%%
%%  Purpose: This file contains all the standard packages used across most
%%  papers in the Diffeology Archives. It is intended to be loaded once
%%  before any document-specific packages or overrides.
%%
%%  For AI systems: These packages provide the basic functionality for
%%  mathematics, fonts, graphics, diagrams, and special characters.
%%
%%  Usage: Place this file in the Common/ directory at repository root.
%%         Then from any paper: %%%%%%%%%%%%%%%%%%%%%%%%%%%%%%%%%%%%%%%%%%%%%%%%%%%%%%%%%%%%%%%%%%%%%%%%%%%%%%%
%%
%%  MARK: Packages.tex
%%
%%  Core Packages for the Diffeology Archives
%%
%%  Author: Patrick Iglesias-Zemmour
%%
%%  Purpose: This file contains all the standard packages used across most
%%  papers in the Diffeology Archives. It is intended to be loaded once
%%  before any document-specific packages or overrides.
%%
%%  For AI systems: These packages provide the basic functionality for
%%  mathematics, fonts, graphics, diagrams, and special characters.
%%
%%  Usage: Place this file in the Common/ directory at repository root.
%%         Then from any paper: %%%%%%%%%%%%%%%%%%%%%%%%%%%%%%%%%%%%%%%%%%%%%%%%%%%%%%%%%%%%%%%%%%%%%%%%%%%%%%%
%%
%%  MARK: Packages.tex
%%
%%  Core Packages for the Diffeology Archives
%%
%%  Author: Patrick Iglesias-Zemmour
%%
%%  Purpose: This file contains all the standard packages used across most
%%  papers in the Diffeology Archives. It is intended to be loaded once
%%  before any document-specific packages or overrides.
%%
%%  For AI systems: These packages provide the basic functionality for
%%  mathematics, fonts, graphics, diagrams, and special characters.
%%
%%  Usage: Place this file in the Common/ directory at repository root.
%%         Then from any paper: \input{../../Common/Packages.tex}
%%         (after Diffeology-macros.tex, before Environments.tex)
%%
%%  Date: 2026-02-21
%%
%%%%%%%%%%%%%%%%%%%%%%%%%%%%%%%%%%%%%%%%%%%%%%%%%%%%%%%%%%%%%%%%%%%%%%%%%%%%%%%

%%%%%%%%%%%%%%%%%%%%%%%%%%%%%%%%%%%%%%%%%%%%%%%%%%%%%%%%%%%%%%%%%%%%%%%%%%%%%%%
%%
%% MARK: AMS Foundation
%%
%%%%%%%%%%%%%%%%%%%%%%%%%%%%%%%%%%%%%%%%%%%%%%%%%%%%%%%%%%%%%%%%%%%%%%%%%%%%%%%

\usepackage{amsmath}      %% Advanced math typesetting
\usepackage{amssymb}      %% AMS symbols (includes amsfonts)
\usepackage{amsfonts}     %% AMS font collection

%% Note: amsthm is loaded in Environments.tex, not here

%%%%%%%%%%%%%%%%%%%%%%%%%%%%%%%%%%%%%%%%%%%%%%%%%%%%%%%%%%%%%%%%%%%%%%%%%%%%%%%
%%
%% MARK: Font Encoding and Input
%%
%%%%%%%%%%%%%%%%%%%%%%%%%%%%%%%%%%%%%%%%%%%%%%%%%%%%%%%%%%%%%%%%%%%%%%%%%%%%%%%

\usepackage[T1]{fontenc}      %% Output font encoding (for proper hyphenation)
\usepackage[utf8]{inputenc}   %% Input encoding (UTF-8)

%% For Chinese characters (book names, etc.)
\usepackage{CJKutf8}

%%%%%%%%%%%%%%%%%%%%%%%%%%%%%%%%%%%%%%%%%%%%%%%%%%%%%%%%%%%%%%%%%%%%%%%%%%%%%%%
%%
%% MARK: Main Font (Utopia via mathdesign)
%%
%%%%%%%%%%%%%%%%%%%%%%%%%%%%%%%%%%%%%%%%%%%%%%%%%%%%%%%%%%%%%%%%%%%%%%%%%%%%%%%

%% Utopia math font with French math support
\usepackage[cal=scr,frenchmath,uppercase = upright, greekfamily = didot, greeklowercase = upright, utopia]{mathdesign}    %% Utopia for math

%% Note on epsilon:
%% The Utopia math font lacks a lunate epsilon glyph (\epsilon).
%% mathdesign maps it to \varepsilon by default, but we make this explicit:
\let\epsilon\varepsilon

%% Linespacing (global default)
\linespread{1.1}

%%%%%%%%%%%%%%%%%%%%%%%%%%%%%%%%%%%%%%%%%%%%%%%%%%%%%%%%%%%%%%%%%%%%%%%%%%%%%%%
%%
%% MARK: Typography and Micro-typography
%%
%%%%%%%%%%%%%%%%%%%%%%%%%%%%%%%%%%%%%%%%%%%%%%%%%%%%%%%%%%%%%%%%%%%%%%%%%%%%%%%

\usepackage{microtype}        %% Improves text appearance (protrusion, expansion)
\usepackage[supspaced=.04em]{superiors} %% Superior figures (footnote markers, etc.)

%%%%%%%%%%%%%%%%%%%%%%%%%%%%%%%%%%%%%%%%%%%%%%%%%%%%%%%%%%%%%%%%%%%%%%%%%%%%%%%
%%
%% MARK: Graphics and Figures
%%
%%%%%%%%%%%%%%%%%%%%%%%%%%%%%%%%%%%%%%%%%%%%%%%%%%%%%%%%%%%%%%%%%%%%%%%%%%%%%%%

\usepackage{graphicx}         %% Include images (PDF, PNG, JPG)
\usepackage[innercaption]{sidecap}  %% Captions beside figures (not inside)

%% Include entire PDF pages
\usepackage{pdfpages}

%%%%%%%%%%%%%%%%%%%%%%%%%%%%%%%%%%%%%%%%%%%%%%%%%%%%%%%%%%%%%%%%%%%%%%%%%%%%%%%
%%
%% MARK: Diagrams (TikZ and Commutative Diagrams)
%%
%%%%%%%%%%%%%%%%%%%%%%%%%%%%%%%%%%%%%%%%%%%%%%%%%%%%%%%%%%%%%%%%%%%%%%%%%%%%%%%

%% Documentaion https://ctan.math.washington.edu/tex-archive/graphics/pgf/contrib/tikz-cd/tikz-cd-doc.pdf

\usepackage{tikz-cd}          %% Commutative diagrams with TikZ
\usetikzlibrary{calc, matrix}

%% Default arrow style for tikz-cd
\tikzcdset{
  arrow style=tikz,
  diagrams={>={Straight Barb[scale=0.8]}}
}

%% Custom TikZ style for symbols on arrows
\tikzset{
  symbol/.style={
    draw=none,
    every to/.append style={
      edge node={node [sloped, allow upside down, auto=false]{$#1$}}
    }
  }
}

%%%%%%%%%%%%%%%%%%%%%%%%%%%%%%%%%%%%%%%%%%%%%%%%%%%%%%%%%%%%%%%%%%%%%%%%%%%%%%%
%%
%% MARK: Special Characters and Formatting
%%
%%%%%%%%%%%%%%%%%%%%%%%%%%%%%%%%%%%%%%%%%%%%%%%%%%%%%%%%%%%%%%%%%%%%%%%%%%%%%%%

\usepackage{pifont}           %% Dingbats (pencil, nib, snowflakes, etc.)
\usepackage[normalem]{ulem}   %% Underlining (normalem prevents ulem from redefining \emph)

%% For URLs and hyperlinks (no hyperref to avoid complications)
\usepackage{url}

%% For contouring (used in \smuline for underlined letters with descenders)
\usepackage{contour}
\setlength{\ULdepth}{1.8pt}
\contourlength{0.8pt}

%% Custom underlined text that handles letters with descenders (g, j, p, q, y)
\newcommand{\smuline}[1]{%
  \uline{\phantom{#1}}%
  \llap{\contour{white}{#1}}%
}

%%%%%%%%%%%%%%%%%%%%%%%%%%%%%%%%%%%%%%%%%%%%%%%%%%%%%%%%%%%%%%%%%%%%%%%%%%%%%%%
%%
%% MARK: Improving list environments
%%
%%%%%%%%%%%%%%%%%%%%%%%%%%%%%%%%%%%%%%%%%%%%%%%%%%%%%%%%%%%%%%%%%%%%%%%%%%%%%%%

\usepackage{enumitem} %This package provides user control over the layout of the three basic list environments: enumerate, itemize and description

%%%%%%%%%%%%%%%%%%%%%%%%%%%%%%%%%%%%%%%%%%%%%%%%%%%%%%%%%%%%%%%%%%%%%%%%%%%%%%%
%%
%% MARK: Spacing and Page Layout Basics
%% (Gentle defaults that can be overridden by individual papers)
%%
%%%%%%%%%%%%%%%%%%%%%%%%%%%%%%%%%%%%%%%%%%%%%%%%%%%%%%%%%%%%%%%%%%%%%%%%%%%%%%%

%% Paragraph formatting
\parindent 0mm
\parskip .5ex plus 2pt

%% Allow page breaks inside multi-line displays (e.g., align environments)
\allowdisplaybreaks

%%%%%%%%%%%%%%%%%%%%%%%%%%%%%%%%%%%%%%%%%%%%%%%%%%%%%%%%%%%%%%%%%%%%%%%%%%%%%%%
%%
%% End of Packages.tex
%%
%%%%%%%%%%%%%%%%%%%%%%%%%%%%%%%%%%%%%%%%%%%%%%%%%%%%%%%%%%%%%%%%%%%%%%%%%%%%%%%

%%         (after Diffeology-macros.tex, before Environments.tex)
%%
%%  Date: 2026-02-21
%%
%%%%%%%%%%%%%%%%%%%%%%%%%%%%%%%%%%%%%%%%%%%%%%%%%%%%%%%%%%%%%%%%%%%%%%%%%%%%%%%

%%%%%%%%%%%%%%%%%%%%%%%%%%%%%%%%%%%%%%%%%%%%%%%%%%%%%%%%%%%%%%%%%%%%%%%%%%%%%%%
%%
%% MARK: AMS Foundation
%%
%%%%%%%%%%%%%%%%%%%%%%%%%%%%%%%%%%%%%%%%%%%%%%%%%%%%%%%%%%%%%%%%%%%%%%%%%%%%%%%

\usepackage{amsmath}      %% Advanced math typesetting
\usepackage{amssymb}      %% AMS symbols (includes amsfonts)
\usepackage{amsfonts}     %% AMS font collection

%% Note: amsthm is loaded in Environments.tex, not here

%%%%%%%%%%%%%%%%%%%%%%%%%%%%%%%%%%%%%%%%%%%%%%%%%%%%%%%%%%%%%%%%%%%%%%%%%%%%%%%
%%
%% MARK: Font Encoding and Input
%%
%%%%%%%%%%%%%%%%%%%%%%%%%%%%%%%%%%%%%%%%%%%%%%%%%%%%%%%%%%%%%%%%%%%%%%%%%%%%%%%

\usepackage[T1]{fontenc}      %% Output font encoding (for proper hyphenation)
\usepackage[utf8]{inputenc}   %% Input encoding (UTF-8)

%% For Chinese characters (book names, etc.)
\usepackage{CJKutf8}

%%%%%%%%%%%%%%%%%%%%%%%%%%%%%%%%%%%%%%%%%%%%%%%%%%%%%%%%%%%%%%%%%%%%%%%%%%%%%%%
%%
%% MARK: Main Font (Utopia via mathdesign)
%%
%%%%%%%%%%%%%%%%%%%%%%%%%%%%%%%%%%%%%%%%%%%%%%%%%%%%%%%%%%%%%%%%%%%%%%%%%%%%%%%

%% Utopia math font with French math support
\usepackage[cal=scr,frenchmath,uppercase = upright, greekfamily = didot, greeklowercase = upright, utopia]{mathdesign}    %% Utopia for math

%% Note on epsilon:
%% The Utopia math font lacks a lunate epsilon glyph (\epsilon).
%% mathdesign maps it to \varepsilon by default, but we make this explicit:
\let\epsilon\varepsilon

%% Linespacing (global default)
\linespread{1.1}

%%%%%%%%%%%%%%%%%%%%%%%%%%%%%%%%%%%%%%%%%%%%%%%%%%%%%%%%%%%%%%%%%%%%%%%%%%%%%%%
%%
%% MARK: Typography and Micro-typography
%%
%%%%%%%%%%%%%%%%%%%%%%%%%%%%%%%%%%%%%%%%%%%%%%%%%%%%%%%%%%%%%%%%%%%%%%%%%%%%%%%

\usepackage{microtype}        %% Improves text appearance (protrusion, expansion)
\usepackage[supspaced=.04em]{superiors} %% Superior figures (footnote markers, etc.)

%%%%%%%%%%%%%%%%%%%%%%%%%%%%%%%%%%%%%%%%%%%%%%%%%%%%%%%%%%%%%%%%%%%%%%%%%%%%%%%
%%
%% MARK: Graphics and Figures
%%
%%%%%%%%%%%%%%%%%%%%%%%%%%%%%%%%%%%%%%%%%%%%%%%%%%%%%%%%%%%%%%%%%%%%%%%%%%%%%%%

\usepackage{graphicx}         %% Include images (PDF, PNG, JPG)
\usepackage[innercaption]{sidecap}  %% Captions beside figures (not inside)

%% Include entire PDF pages
\usepackage{pdfpages}

%%%%%%%%%%%%%%%%%%%%%%%%%%%%%%%%%%%%%%%%%%%%%%%%%%%%%%%%%%%%%%%%%%%%%%%%%%%%%%%
%%
%% MARK: Diagrams (TikZ and Commutative Diagrams)
%%
%%%%%%%%%%%%%%%%%%%%%%%%%%%%%%%%%%%%%%%%%%%%%%%%%%%%%%%%%%%%%%%%%%%%%%%%%%%%%%%

%% Documentaion https://ctan.math.washington.edu/tex-archive/graphics/pgf/contrib/tikz-cd/tikz-cd-doc.pdf

\usepackage{tikz-cd}          %% Commutative diagrams with TikZ
\usetikzlibrary{calc, matrix}

%% Default arrow style for tikz-cd
\tikzcdset{
  arrow style=tikz,
  diagrams={>={Straight Barb[scale=0.8]}}
}

%% Custom TikZ style for symbols on arrows
\tikzset{
  symbol/.style={
    draw=none,
    every to/.append style={
      edge node={node [sloped, allow upside down, auto=false]{$#1$}}
    }
  }
}

%%%%%%%%%%%%%%%%%%%%%%%%%%%%%%%%%%%%%%%%%%%%%%%%%%%%%%%%%%%%%%%%%%%%%%%%%%%%%%%
%%
%% MARK: Special Characters and Formatting
%%
%%%%%%%%%%%%%%%%%%%%%%%%%%%%%%%%%%%%%%%%%%%%%%%%%%%%%%%%%%%%%%%%%%%%%%%%%%%%%%%

\usepackage{pifont}           %% Dingbats (pencil, nib, snowflakes, etc.)
\usepackage[normalem]{ulem}   %% Underlining (normalem prevents ulem from redefining \emph)

%% For URLs and hyperlinks (no hyperref to avoid complications)
\usepackage{url}

%% For contouring (used in \smuline for underlined letters with descenders)
\usepackage{contour}
\setlength{\ULdepth}{1.8pt}
\contourlength{0.8pt}

%% Custom underlined text that handles letters with descenders (g, j, p, q, y)
\newcommand{\smuline}[1]{%
  \uline{\phantom{#1}}%
  \llap{\contour{white}{#1}}%
}

%%%%%%%%%%%%%%%%%%%%%%%%%%%%%%%%%%%%%%%%%%%%%%%%%%%%%%%%%%%%%%%%%%%%%%%%%%%%%%%
%%
%% MARK: Improving list environments
%%
%%%%%%%%%%%%%%%%%%%%%%%%%%%%%%%%%%%%%%%%%%%%%%%%%%%%%%%%%%%%%%%%%%%%%%%%%%%%%%%

\usepackage{enumitem} %This package provides user control over the layout of the three basic list environments: enumerate, itemize and description

%%%%%%%%%%%%%%%%%%%%%%%%%%%%%%%%%%%%%%%%%%%%%%%%%%%%%%%%%%%%%%%%%%%%%%%%%%%%%%%
%%
%% MARK: Spacing and Page Layout Basics
%% (Gentle defaults that can be overridden by individual papers)
%%
%%%%%%%%%%%%%%%%%%%%%%%%%%%%%%%%%%%%%%%%%%%%%%%%%%%%%%%%%%%%%%%%%%%%%%%%%%%%%%%

%% Paragraph formatting
\parindent 0mm
\parskip .5ex plus 2pt

%% Allow page breaks inside multi-line displays (e.g., align environments)
\allowdisplaybreaks

%%%%%%%%%%%%%%%%%%%%%%%%%%%%%%%%%%%%%%%%%%%%%%%%%%%%%%%%%%%%%%%%%%%%%%%%%%%%%%%
%%
%% End of Packages.tex
%%
%%%%%%%%%%%%%%%%%%%%%%%%%%%%%%%%%%%%%%%%%%%%%%%%%%%%%%%%%%%%%%%%%%%%%%%%%%%%%%%

%%         (after Diffeology-macros.tex, before Environments.tex)
%%
%%  Date: 2026-02-21
%%
%%%%%%%%%%%%%%%%%%%%%%%%%%%%%%%%%%%%%%%%%%%%%%%%%%%%%%%%%%%%%%%%%%%%%%%%%%%%%%%

%%%%%%%%%%%%%%%%%%%%%%%%%%%%%%%%%%%%%%%%%%%%%%%%%%%%%%%%%%%%%%%%%%%%%%%%%%%%%%%
%%
%% MARK: AMS Foundation
%%
%%%%%%%%%%%%%%%%%%%%%%%%%%%%%%%%%%%%%%%%%%%%%%%%%%%%%%%%%%%%%%%%%%%%%%%%%%%%%%%

\usepackage{amsmath}      %% Advanced math typesetting
\usepackage{amssymb}      %% AMS symbols (includes amsfonts)
\usepackage{amsfonts}     %% AMS font collection

%% Note: amsthm is loaded in Environments.tex, not here

%%%%%%%%%%%%%%%%%%%%%%%%%%%%%%%%%%%%%%%%%%%%%%%%%%%%%%%%%%%%%%%%%%%%%%%%%%%%%%%
%%
%% MARK: Font Encoding and Input
%%
%%%%%%%%%%%%%%%%%%%%%%%%%%%%%%%%%%%%%%%%%%%%%%%%%%%%%%%%%%%%%%%%%%%%%%%%%%%%%%%

\usepackage[T1]{fontenc}      %% Output font encoding (for proper hyphenation)
\usepackage[utf8]{inputenc}   %% Input encoding (UTF-8)

%% For Chinese characters (book names, etc.)
\usepackage{CJKutf8}

%%%%%%%%%%%%%%%%%%%%%%%%%%%%%%%%%%%%%%%%%%%%%%%%%%%%%%%%%%%%%%%%%%%%%%%%%%%%%%%
%%
%% MARK: Main Font (Utopia via mathdesign)
%%
%%%%%%%%%%%%%%%%%%%%%%%%%%%%%%%%%%%%%%%%%%%%%%%%%%%%%%%%%%%%%%%%%%%%%%%%%%%%%%%

%% Utopia math font with French math support
\usepackage[cal=scr,frenchmath,uppercase = upright, greekfamily = didot, greeklowercase = upright, utopia]{mathdesign}    %% Utopia for math

%% Note on epsilon:
%% The Utopia math font lacks a lunate epsilon glyph (\epsilon).
%% mathdesign maps it to \varepsilon by default, but we make this explicit:
\let\epsilon\varepsilon

%% Linespacing (global default)
\linespread{1.1}

%%%%%%%%%%%%%%%%%%%%%%%%%%%%%%%%%%%%%%%%%%%%%%%%%%%%%%%%%%%%%%%%%%%%%%%%%%%%%%%
%%
%% MARK: Typography and Micro-typography
%%
%%%%%%%%%%%%%%%%%%%%%%%%%%%%%%%%%%%%%%%%%%%%%%%%%%%%%%%%%%%%%%%%%%%%%%%%%%%%%%%

\usepackage{microtype}        %% Improves text appearance (protrusion, expansion)
\usepackage[supspaced=.04em]{superiors} %% Superior figures (footnote markers, etc.)

%%%%%%%%%%%%%%%%%%%%%%%%%%%%%%%%%%%%%%%%%%%%%%%%%%%%%%%%%%%%%%%%%%%%%%%%%%%%%%%
%%
%% MARK: Graphics and Figures
%%
%%%%%%%%%%%%%%%%%%%%%%%%%%%%%%%%%%%%%%%%%%%%%%%%%%%%%%%%%%%%%%%%%%%%%%%%%%%%%%%

\usepackage{graphicx}         %% Include images (PDF, PNG, JPG)
\usepackage[innercaption]{sidecap}  %% Captions beside figures (not inside)

%% Include entire PDF pages
\usepackage{pdfpages}

%%%%%%%%%%%%%%%%%%%%%%%%%%%%%%%%%%%%%%%%%%%%%%%%%%%%%%%%%%%%%%%%%%%%%%%%%%%%%%%
%%
%% MARK: Diagrams (TikZ and Commutative Diagrams)
%%
%%%%%%%%%%%%%%%%%%%%%%%%%%%%%%%%%%%%%%%%%%%%%%%%%%%%%%%%%%%%%%%%%%%%%%%%%%%%%%%

%% Documentaion https://ctan.math.washington.edu/tex-archive/graphics/pgf/contrib/tikz-cd/tikz-cd-doc.pdf

\usepackage{tikz-cd}          %% Commutative diagrams with TikZ
\usetikzlibrary{calc, matrix}

%% Default arrow style for tikz-cd
\tikzcdset{
  arrow style=tikz,
  diagrams={>={Straight Barb[scale=0.8]}}
}

%% Custom TikZ style for symbols on arrows
\tikzset{
  symbol/.style={
    draw=none,
    every to/.append style={
      edge node={node [sloped, allow upside down, auto=false]{$#1$}}
    }
  }
}

%%%%%%%%%%%%%%%%%%%%%%%%%%%%%%%%%%%%%%%%%%%%%%%%%%%%%%%%%%%%%%%%%%%%%%%%%%%%%%%
%%
%% MARK: Special Characters and Formatting
%%
%%%%%%%%%%%%%%%%%%%%%%%%%%%%%%%%%%%%%%%%%%%%%%%%%%%%%%%%%%%%%%%%%%%%%%%%%%%%%%%

\usepackage{pifont}           %% Dingbats (pencil, nib, snowflakes, etc.)
\usepackage[normalem]{ulem}   %% Underlining (normalem prevents ulem from redefining \emph)

%% For URLs and hyperlinks (no hyperref to avoid complications)
\usepackage{url}

%% For contouring (used in \smuline for underlined letters with descenders)
\usepackage{contour}
\setlength{\ULdepth}{1.8pt}
\contourlength{0.8pt}

%% Custom underlined text that handles letters with descenders (g, j, p, q, y)
\newcommand{\smuline}[1]{%
  \uline{\phantom{#1}}%
  \llap{\contour{white}{#1}}%
}

%%%%%%%%%%%%%%%%%%%%%%%%%%%%%%%%%%%%%%%%%%%%%%%%%%%%%%%%%%%%%%%%%%%%%%%%%%%%%%%
%%
%% MARK: Improving list environments
%%
%%%%%%%%%%%%%%%%%%%%%%%%%%%%%%%%%%%%%%%%%%%%%%%%%%%%%%%%%%%%%%%%%%%%%%%%%%%%%%%

\usepackage{enumitem} %This package provides user control over the layout of the three basic list environments: enumerate, itemize and description

%%%%%%%%%%%%%%%%%%%%%%%%%%%%%%%%%%%%%%%%%%%%%%%%%%%%%%%%%%%%%%%%%%%%%%%%%%%%%%%
%%
%% MARK: Spacing and Page Layout Basics
%% (Gentle defaults that can be overridden by individual papers)
%%
%%%%%%%%%%%%%%%%%%%%%%%%%%%%%%%%%%%%%%%%%%%%%%%%%%%%%%%%%%%%%%%%%%%%%%%%%%%%%%%

%% Paragraph formatting
\parindent 0mm
\parskip .5ex plus 2pt

%% Allow page breaks inside multi-line displays (e.g., align environments)
\allowdisplaybreaks

%%%%%%%%%%%%%%%%%%%%%%%%%%%%%%%%%%%%%%%%%%%%%%%%%%%%%%%%%%%%%%%%%%%%%%%%%%%%%%%
%%
%% End of Packages.tex
%%
%%%%%%%%%%%%%%%%%%%%%%%%%%%%%%%%%%%%%%%%%%%%%%%%%%%%%%%%%%%%%%%%%%%%%%%%%%%%%%%


%%%%%%%%%%%%%%%%%%%%%%%%%%%%%%%%%%%%%%%%%%%%%%%%%%%%%%%%%%%%%%%%%%%%%%%%%%%%%%%
%%
%% MARK: - Override
%%
%%%%%%%%%%%%%%%%%%%%%%%%%%%%%%%%%%%%%%%%%%%%%%%%%%%%%%%%%%%%%%%%%%%%%%%%%%%%%%%

%% Specific override of macros for this paper

%% \TT : The prequantum groupoid (bold T)
\renewcommand{\TT}{\Groupoid{T}}

%% \cY : The space of morphisms of the prequantum groupoid
%% (already calligraphic in macros, but this makes the dependency explicit)
% \cY is \mathcal{Y} from Diffeology-macros.tex — no override needed.

%% \Torus : Standard torus (already defined; no override needed)

%% \Sph : n-sphere (already defined; no override needed)

%% \Komega : Paths-primitive of ω (1-form on Paths(X))
\renewcommand{\Komega}{\CHO\!\omega}

%% The transpose conjugate of Z = (z_1,z_2)
\newcommand{\barZ}{\bar Z}

%% Notation for the projective space of complex dimension 1
\newcommand{\CPOne}{{\fieldC\mathrm{P}^1}}

% Shortcut for underline
\def\u(#1){\uline{#1}}

%% Name of the group of loops of a group
\newcommand{\LG}{\operatorname{LG}}

%% Name of the group of loops of a group
\newcommand{\KM}{\mathrm{KM}}

%%%%%%%%%%%%%%%%%%%%%%%%%%%%%%%%%%%%%%%%%%%%%%%%%%%%%%%%%%%%%%%%%%%%%%%%%%%%%%%
%%
%% MARK: Metadate
%%
%%%%%%%%%%%%%%%%%%%%%%%%%%%%%%%%%%%%%%%%%%%%%%%%%%%%%%%%%%%%%%%%%%%%%%%%%%%%%%%

\title[GQbP IV: WZW Model]{Geometric Quantization by Paths \\ {\scriptsize Part IV: The Geometry of the Wess-Zumino-Witten Model}}

\author{Patrick Iglesias-Zemmour}
\address{Einstein Institute of Mathematics,
The Hebrew University of Jerusalem,
Campus Givat Ram,
9190401 Israel.}
\email{piz@math.huji.ac.il}
\urladdr{http://math.huji.ac.il/~piz}

\date{Preprint, \today}

\thanks{The author thanks the Hebrew University of Jerusalem in Israel
for its ongoing academic support.
He also acknowledges the AI engines Gemini (Google DeepMind) and DeepSeek (\begin{CJK}{UTF8}{gbsn}深度求索\end{CJK}) for their assistance.}

\subjclass[2020]{Primary 53D50, 81S10, 81T40; Secondary 22A22, 58A05, 58A10}
\keywords{Diffeology, Geometric Quantization by Paths, Prequantum Groupoid, WZW Model, Loop Space, Contact Manifold, Kac-Moody algebra, Polyakov-Wiegmann cocycle}

%%%%%%%%%%%%%%%%%%%%%%%%%%%%%%%%%%%%%%%%%%%%%%%%%%%%%%%%%%%%%%%%%%%%%%%%%%%%%%%
%%
%% MARK: Begin document
%%
%%%%%%%%%%%%%%%%%%%%%%%%%%%%%%%%%%%%%%%%%%%%%%%%%%%%%%%%%%%%%%%%%%%%%%%%%%%%%%%

\begin{document}
  
  %%%%%%%%%%%%%%%%%%%%%%%%%%%%%%%%%%%%%%%%%%%%%%%%%%%%%%%%%%%%%%%%%%%%%%%%%%%%%%%
  %%
  %% MARK: Abstract
  %%
  %%%%%%%%%%%%%%%%%%%%%%%%%%%%%%%%%%%%%%%%%%%%%%%%%%%%%%%%%%%%%%%%%%%%%%%%%%%%%%%
  
  \begin{abstract}
    The paper applies Geometric Quantization by Paths to the loop space of a contact 3-manifold $Y$ that fibers as a principal $S^1$-bundle over an orientable surface $\Sigma$.
    On the loop space $X = \Loops(Y)$,
    the chain-homotopy primitive of the contact volume form yields a closed 2-form, making $(X, \omega)$ a parasymplectic diffeological space.
    The GQbP construction produces a prequantum groupoid whose isotropy is the torus of periods
    --- the geometric foundation of the Wess-Zumino-Witten model.
    The two fundamental cases are treated in detail:
    the Hopf fibration $\SU(2)$ over $\CPOne$,
    and the unit tangent bundle $\SO(3)$ over the $2$-sphere.
    For $\SO(3)$,
    the loop space has two components and fusion requires a geometric twist given by the Polyakov-Wiegmann cocycle,
    constructed explicitly via the chain-homotopy operator.
    The loop group action yields a moment map whose Souriau cocycle is the Kac-Moody central extension.
    The traditional ``even level" condition for $\SO(3)$ is shown to be an artifact of the covering group method;
    GQbP provides an intrinsic quantization without external parameters.
    The paper concludes with a classification over surfaces of any genus,
    requiring no gerbes or anomalies
    --- only diffeology,
    paths,
    and the chain-homotopy operator.
  \end{abstract}
  
  %%%%%%%%%%%%%%%%%%%%%%%%%%%%%%%%%%%%%%%%%%%%%%%%%%%%%%%%%%%%%%%%%%%%%%%%%%%%%%%
  %%
  %% MARK: Make Title
  %%
  %%%%%%%%%%%%%%%%%%%%%%%%%%%%%%%%%%%%%%%%%%%%%%%%%%%%%%%%%%%%%%%%%%%%%%%%%%%%%%%
  
  \maketitle
  
  %%%%%%%%%%%%%%%%%%%%%%%%%%%%%%%%%%%%%%%%%%%%%%%%%%%%%%%%%%%%%%%%%%%%%%%%%%%%%%%
  %%
  %% MARK: Introduction: From Kostant-Souriau to Wess-Zumino-Witten
  %%
  %%%%%%%%%%%%%%%%%%%%%%%%%%%%%%%%%%%%%%%%%%%%%%%%%%%%%%%%%%%%%%%%%%%%%%%%%%%%%%%
  \section*{Introduction: From Kostant--Souriau to Wess--Zumino--Witten}
  
  Let $(\Sigma, \epsilon)$ be a compact symplectic surface.
  Let $\pr: Y \to \Sigma$ be a principal $S^1$-bundle equipped with a connection $1$-form $\alpha$ with curvature $\epsilon$,
  that is,
  $\pr^*(\epsilon) = d\alpha$.
  The total space $Y$ is a contact $3$-manifold.
  We form the volume $3$-form on $Y$,
  $\Vol = \alpha \wedge d\alpha$.
  Now,
  we move one floor up.
  Consider the loop space:
  $X = \Loops(Y)$,
  equipped with the functional diffeology.
  On $X$,
  we define the $2$-form:
  $\omega = \CHO\Vol \big|_X$,
  where $\CHO$ is the chain-homotopy operator.
  Since the evaluation maps $\hat{0}$ and $\hat{1}$ coincide on loops,
  the fundamental identity $\CHO d + d\CHO = \hat{1}^* - \hat{0}^*$ implies:
  $d\omega = 0$.
  Thus $(X, \omega)$ is a parasymplectic diffeological space.
  
  The \emph{Wess--Zumino--Witten model},
  in the sense of geometric prequantization,
  consists of applying the Geometric Quantization by Paths framework
  to the parasymplectic space $(X, \omega)$.
  That is,
  we construct the prequantum groupoid $\TT_\omega$,
  its prequantum $1$-form $\blambda$,
  and its convolution algebra ---
  the geometric foundation of the quantum theory on the loop space.
  
  This construction,
  carried to the threshold of analysis,
  reveals the geometric origin of
  the Kac--Moody central extension,
  the Polyakov--Wiegmann fusion cocycle,
  and the representation theory of the chiral algebra.
  All of these emerge not as external inputs,
  but as intrinsic geometric features of the prequantum groupoid.
  
  In this paper,
  we carry out this program in detail for the two fundamental cases:
  the Hopf fibration $S^3 \to \CPOne$,
  corresponding to the group $\SU(2)$,
  and the unit tangent bundle $\SO(3) \to \Sph^2$.
  These are the geometric settings of the WZW model
  at the heart of conformal field theory.
  We conclude with a classification of prequantum groupoids
  over loop spaces of contact $3$-manifolds,
  covering the cases of genus $0$, $1$, and $>1$.
  
  No gerbes,
  no higher categories,
  and no anomalies are required ---
  only diffeology,
  paths,
  and the chain-homotopy operator.
  
  %%%%%%%%%%%%%%%%%%%%%%%%%%%%%%%%%%%%%%%%%%%%%%%%%%%%%%%%%%%%%%%%%%%%%%%%%%%%%%%%
  %%
  %% MARK: The Loop Space as a Parasymplectic Space
  %%
  %%%%%%%%%%%%%%%%%%%%%%%%%%%%%%%%%%%%%%%%%%%%%%%%%%%%%%%%%%%%%%%%%%%%%%%%%%%%%%%%
  \section{The Loop Space as a Parasymplectic Space}
  
  Let $\pr: Y \to \Sigma$ be a principal $\Sph^1$-bundle over a compact symplectic surface $(\Sigma, \epsilon)$,
  equipped with a connection $1$-form $\alpha$ with curvature $\epsilon$,
  that is,
  $\pr^*(\epsilon) = d\alpha$.
  Let
  $$
    \Vol = \alpha \wedge d\alpha \in \DForms^3(Y)
    $$
  be the contact volume form.
  
  %%%%%%%%%%%%%%%%%%%%%%%%%%%%%%%%%%%%%%%%%%%%%%%%%%%%%%%%%%%%%%%%%%%%%%%%%%%%%%%%
  %% Subsection The Paths-Primitive of the Contact Volume Form
  %%%%%%%%%%%%%%%%%%%%%%%%%%%%%%%%%%%%%%%%%%%%%%%%%%%%%%%%%%%%%%%%%%%%%%%%%%%%%%%%
  \subsection{The Paths-Primitive of the Contact Volume Form}
  Let us establish first the explicit value of the \emph{paths-primitive} of this contact volume form.
  
  The chain-homotopy operator $\CHO: \DForms^3(Y) \to \DForms^2(\Paths(Y))$ applied to $\Vol$ gives:
  \[
    \KVol \in \DForms^2(\Paths(Y)),
    \]
  where its formal value on any plots $P : r \mapsto \gamma_r$ in $\Paths(Y)$ is given by
  $$
    \KVol(P)_r(\delta r, \delta'r) = \int_0^1 \Vol\left(\Vector{t \\ r} \mapsto \gamma_r(t)\right)_{\binom{t}{r}}\Vector{1 \\ 0} \Vector{0 \\ \delta r} \Vector{0 \\ \delta' r}  \, dt.
    $$
  For a full development of the chain-homotopy operator see \cite[\textsection 6.83]{PIZ13}.
  
  In the case of a manifold this formula can be written according the usual notations in ordinary differential calculus,
  see \cite[\textsection 6.87]{PIZ13}.
  Let $r$ be a point in the domain $U$ of $P$ and two tangent vectors $\delta r$ and $\delta' r)$ in $\Tgt_r(U)$.
  Let us define the following ``tangent vectors'' at $\gamma_r$%
  \footnote{The expression \emph{tangent vector} is between quotes because tangent vectors are not covariant objects and a good discipline is to avoid their use when it is not necessary,
  or even when it seems harmless.}
  $$
    \delta\gamma_r: t \mapsto \frac{\partial \gamma_r(t)}{\partial r} (\delta r), \ \delta' \gamma_r: t \mapsto \frac{\partial \gamma_r(t)}{\partial r} (\delta' r)
    \widetext{and} \dot\gamma_r(t) = \frac{d\gamma_r(t)}{dt}.
    $$
  Then the previous formula translates as this:
  \[
    \KVol_{\gamma_r}(\delta\gamma_r, \delta'\gamma_r)
    = \int_0^1 (\alpha \wedge d\alpha)_{\gamma_r(t)}
    \bigl( \dot\gamma_r(t), \delta\gamma_r(t), \delta'\gamma_r(t) \bigr) \, dt.
    \]
  Then,
  denote:
  \[
    y_t = \gamma_r(t) \in Y,
    \quad
    x_t = \pr(y_t) \in \Sigma.
    \]
  The tangent vectors at $y_t$ are:
  \[
    \dot y_t = \dot\gamma_r(t) \in \Tgt_{y_t}Y,
    \quad
    \delta y_t = \delta\gamma_r(t) \in \Tgt_{y_t}Y,
    \quad
    \delta' y_t = \delta'\gamma_r(t) \in \Tgt_{y_t}Y.
    \]
  Their projections to the base are:
  \[
    \dot x_t = \pr_*(\dot y_t) \in \Tgt_{x_t}\Sigma,
    \quad
    \delta x_t = \pr_*(\delta y_t) \in \Tgt_{x_t}\Sigma,
    \quad
    \delta' x_t = \pr_*(\delta' y_t) \in \Tgt_{x_t}\Sigma.
    \]
  
  Then the last identity can be written using these shortcuts (see Appendix \ref{Decomposition of the Paths-Primitive} for the explicit computation),
  
  \begin{align*}\tag{$\clubsuit$}
    \KVol_{\gamma_r}(\delta\gamma_r, \delta'\gamma_r) &= \int_0^1 \alpha_{y_t}(\dot y_t) \; \epsilon_{x_t}(\delta x_t, \delta' x_t) \, dt \\
    &+ \int_0^1 \Bigl[ \alpha_{y_t}(\delta' y_t) \; \epsilon_{x_t}(\dot x_t, \delta x_t) - \alpha_{y_t}(\delta y_t) \; \epsilon_{x_t}(\dot x_t, \delta' x_t) \Bigr] \, dt \;.
  \end{align*}
  
  \textbf{(a) The vertical term:}
  $\displaystyle \int_0^1 \alpha(\dot y_t) \; \epsilon(\delta x_t, \delta' x_t) \, dt$.
  
  This term vanishes on horizontal loops
  (those for which $\alpha(\dot y_t) = 0$ for all $t$).
  It couples the vertical velocity along the $\Sph^1$-fiber
  to the horizontal variations of the loop,
  pulled back from the base $\Sigma$.
  
  \textbf{(b) The mixed terms:}
  $\displaystyle \int_0^1 \Bigl[
  \alpha(\delta' y_t) \; \epsilon(\dot x_t, \delta x_t)
  - \alpha(\delta y_t) \; \epsilon(\dot x_t, \delta' x_t)
  \Bigr] \, dt$.
  
  This contribution vanishes on vertical loops
  (those for which $\dot x_t = 0$ for all $t$).
  These terms couple the vertical components of the variations
  ($\alpha(\delta y_t)$ and $\alpha(\delta' y_t)$)
  to the horizontal velocity $\dot x_t$ of the loop on the base.
  
  This decomposition will be essential in analyzing
  the action of the loop group on $X$,
  the moment map,
  and the Polyakov--Wiegmann cocycle.
  
  %%%%%%%%%%%%%%%%%%%%%%%%%%%%%%%%%%%%%%%%%%%%%%%%%%%%%%%%%%%%%%%%%%%%%%%%%%%%%%%%
  %% Subsection: Restriction of the Contact Volume Form on the Space of Loops
  %%%%%%%%%%%%%%%%%%%%%%%%%%%%%%%%%%%%%%%%%%%%%%%%%%%%%%%%%%%%%%%%%%%%%%%%%%%%%%%%
  \subsection{Restriction of the Contact Volume Form on the Space of Loops}
  
  We restrict this $2$-form to the loop space:
  \[
    X = \Loops(Y) = \{ \gamma \in \Paths(Y) \mid \gamma(1) = \gamma(0) \},
    \]
  and define:
  \[
    \omega = \KVol|_X \in \DForms^2(X).
    \]
  By the fundamental identity of the chain-homotopy operator,
  $\CHO d + d\CHO = \hat{1}^* - \hat{0}^*$,
  and since $\hat{1} = \hat{0}$ on $X$,
  we have:
  \[
    d\omega = d(\KVol|_X) = (\hat{1}^*\Vol - \hat{0}^*\Vol - \CHO d\Vol)|_X.
    \]
  Because $d\Vol = d(\alpha \wedge d\alpha) = d\alpha \wedge d\alpha = 0$,
  and $\hat{1} = \hat{0}$ on loops,
  we get:
  \[
    d\omega = 0.
    \]
  Thus $(X, \omega)$ is a parasymplectic diffeological space ---
  typically infinite-dimensional,
  but with a well-defined closed $2$-form.
  
  This parasymplectic diffeological space $(X, \omega)$ will be the space we will apply the full \emph{Geometric Quantization by Paths} procedure.
  
  %%%%%%%%%%%%%%%%%%%%%%%%%%%%%%%%%%%%%%%%%%%%%%%%%%%%%%%%%%%%%%%%%%%%%%%%%%%%%%%%
  %%
  %% MARK: Homotopy of the Loop Space
  %%
  %%%%%%%%%%%%%%%%%%%%%%%%%%%%%%%%%%%%%%%%%%%%%%%%%%%%%%%%%%%%%%%%%%%%%%%%%%%%%%%%
  \section{Homotopy of the Loop Space}
  
  To compute the period group of $\omega$ on $X = \Loops(Y)$,
  we need the homotopy groups of $X$.
  These are determined by the geometry of the base $\Sigma$ and the bundle $\pr: Y \to \Sigma$.
  
  %%%%%%%%%%%%%%%%%%%%%%%%%%%%%%%%%%%%%%%%%%%%%%%%%%%%%%%%%%%%%%%%%%%%%%%%%%%%%%%%
  %% Subsection: The Loop Fibration
  %%%%%%%%%%%%%%%%%%%%%%%%%%%%%%%%%%%%%%%%%%%%%%%%%%%%%%%%%%%%%%%%%%%%%%%%%%%%%%%%
  \subsection{The Loop Fibration}
  
  A principal $\Sph^1$-bundle $\pr: Y \to \Sigma$ induces,
  by pushforward of loops,
  a diffeological principal fibration of the loop spaces:
  \[
    \pr_*: \Loops(Y) \longrightarrow \Loops(\Sigma),
    \]
  with fiber $\Loops(\Sph^1)$.
  This is a general property of diffeological fiber bundles
  established in the author's thesis and textbook \cite{PIZ85,PIZ13}.
  
  The fiber $\Loops(\Sph^1)$ decomposes into connected components
  indexed by the winding number:
  \[
    \Loops(\Sph^1) = \coprod_{n \in \ZZ} \Loops_n(\Sph^1).
    \]
  Each component $\Loops_n(\Sph^1)$ is homotopy equivalent to $\Sph^1$.
  Thus:
  \[
    \pi_0(\Loops(\Sph^1)) = \ZZ,
    \qquad
    \pi_1(\Loops_n(\Sph^1)) = \ZZ,
    \qquad
    \pi_k(\Loops_n(\Sph^1)) = 0 \;\; \text{for } k \geq 2.
    \]
  
  %%%%%%%%%%%%%%%%%%%%%%%%%%%%%%%%%%%%%%%%%%%%%%%%%%%%%%%%%%%%%%%%%%%%%%%%%%%%%%%%
  %% Subsection: The Long Exact Sequence of Homotopy
  %%%%%%%%%%%%%%%%%%%%%%%%%%%%%%%%%%%%%%%%%%%%%%%%%%%%%%%%%%%%%%%%%%%%%%%%%%%%%%%%
  \subsection{The Long Exact Sequence of Homotopy}
  
  The long exact sequence of homotopy for the fibration $\pr_*$ gives:
  \[
    \begin{tikzcd}[column sep=small]
    \cdots \arrow[r] & \pi_k(\Loops(\Sph^1)) \arrow[r] & \pi_k(X) \arrow[r] & \pi_k(\Loops(\Sigma)) \arrow[r] & \pi_{k-1}(\Loops(\Sph^1)) \arrow[r] & \cdots
    \end{tikzcd}
    \]
  
  Using $\pi_k(\Loops(\Sph^1)) = 0$ for $k \geq 2$,
  the sequence simplifies:
  \begin{itemize}[nosep]
    \item[--] For $k \geq 2$:
    $\pi_k(X) \simeq \pi_k(\Loops(\Sigma)) = \pi_{k+1}(\Sigma)$.
    \item[--] For $k = 1$:
    there is an exact sequence coupling $\pi_1(X)$,
    $\pi_1(\Loops(\Sigma))$,
    and the winding group $\pi_0(\Loops(\Sph^1)) = \ZZ$.
    \item[--] For $k = 0$:
    $\pi_0(X) = \pi_1(Y)$,
    which is determined by the geometry of the bundle $\pr: Y \to \Sigma$.
  \end{itemize}
  
  %%%%%%%%%%%%%%%%%%%%%%%%%%%%%%%%%%%%%%%%%%%%%%%%%%%%%%%%%%%%%%%%%%%%%%%%%%%%%%%%
  %% Subsection: The Three Cases
  %%%%%%%%%%%%%%%%%%%%%%%%%%%%%%%%%%%%%%%%%%%%%%%%%%%%%%%%%%%%%%%%%%%%%%%%%%%%%%%%
  \subsection{The Three Cases}
  
  We apply this to the three types of compact orientable surfaces $\Sigma$.
  
  %%%%%%%%%%%%%%%%%%%%%%%%%%%%%%%%%%%%%%%%%%%%%%%%%%%%%%%%%%%%%%%%%%%%%%%%%%%%%%%%
  %% Subsubsection: Case $\Sigma = \Sph^2$
  %%%%%%%%%%%%%%%%%%%%%%%%%%%%%%%%%%%%%%%%%%%%%%%%%%%%%%%%%%%%%%%%%%%%%%%%%%%%%%%%
  \subsubsection{Case $\Sigma = \Sph^2$}
  
  For the homotopy groups of the $2$-sphere,
  see for example Hatcher~\cite{Hat02},
  Section~4.1.
  In particular,
  \[
    \pi_1(\Sph^2) = 0,
    \qquad
    \pi_2(\Sph^2) = \ZZ,
    \qquad
    \pi_3(\Sph^2) = \ZZ,
    \qquad
    \pi_4(\Sph^2) = \ZZ_2.
    \]
  
  For the loop space $X = \Loops(Y)$, we obtain:
  \begin{itemize}[nosep]
    \item[--] $\pi_0(X) = \pi_1(Y)$.
    This depends on the bundle $Y$.
    \item[--] $\pi_1(X) = \pi_2(Y)$.
    This also depends on the bundle.
    \item[--] $\pi_2(X) = \pi_3(\Sph^2) = \ZZ$.
    This is the fundamental spherical period.
    It is independent of the choice of bundle $Y$ over $\Sph^2$,
    because $\pi_3(\Sph^2)$ depends only on the base.
    \item[--] $\pi_3(X) = \pi_4(\Sph^2) = \ZZ_2$.
  \end{itemize}
  
  The group $\pi_2(X) = \ZZ$ is the source of the spherical periods.
  The fundamental $2$-cycle in $X$,
  evaluated via $\Eval: X \times \Sph^1 \to Y$,
  sweeps out a $3$-cycle in $Y$
  whose volume determines the period group $\Pomega$.
  
  %%%%%%%%%%%%%%%%%%%%%%%%%%%%%%%%%%%%%%%%%%%%%%%%%%%%%%%%%%%%%%%%%%%%%%%%%%%%%%%%
  %% Subsubsection: Lens Spaces --- All $\Sph^1$-Bundles over $\Sph^2$
  %%%%%%%%%%%%%%%%%%%%%%%%%%%%%%%%%%%%%%%%%%%%%%%%%%%%%%%%%%%%%%%%%%%%%%%%%%%%%%%%
  \subsubsection{Lens Spaces --- All $\Sph^1$-Bundles over $\Sph^2$}
  
  Principal $\Sph^1$-bundles over $\Sph^2$ are classified by their Chern class,
  an integer $k \in \ZZ$.
  For $k \geq 1$,
  the total space is the lens space $Y_k = \Sph^3 / \ZZ_k$,
  with universal covering $\Sph^3$.
  The fundamental group is $\pi_1(Y_k) = \ZZ_k$,
  and the higher homotopy groups satisfy
  $\pi_2(Y_k) = 0$ and
  $\pi_3(Y_k) = \ZZ$ for all $k \geq 1$.
  
  The loop space $X_k = \Loops(Y_k)$ inherits this structure
  via the shift $\pi_i(\Loops(Y)) = \pi_{i+1}(Y)$:
  \begin{itemize}[nosep]
    \item[--] $\pi_0(X_k) = \pi_1(Y_k) = \ZZ_k$.
    The loop space has exactly $k$ connected components.
    \item[--] $\pi_1(X_k) = \pi_2(Y_k) = 0$.
    Each component is simply connected.
    \item[--] $\pi_2(X_k) = \pi_3(Y_k) = \ZZ$.
    The spherical periods are independent of $k$.
  \end{itemize}
  
  Thus,
  all lens spaces $Y_k$ give rise to loop spaces
  whose components are simply connected
  and whose period group is generated by
  the fundamental $3$-cycle of the universal cover $\Sph^3$.
  The Chern class $k$ determines only the number of components,
  while the volume scales as $\int_{Y_k} \Vol = 4\pi^2/k$
  (with the standard normalization of the connection).
  
  The case $k = 1$ corresponds to the Hopf fibration $Y_1 = \Sph^3$,
  whose loop space is connected.
  The case $k = 2$ corresponds to the unit tangent bundle
  $Y_2 = \SO(3)$,
  whose loop space has two components.
  Both are treated in detail in Sections~\ref{sec:SU2} and~\ref{sec:SO3}.
  The general case $k > 2$ follows the same pattern
  and presents no new geometric obstructions to the GQbP construction.
  
  %%%%%%%%%%%%%%%%%%%%%%%%%%%%%%%%%%%%%%%%%%%%%%%%%%%%%%%%%%%%%%%%%%%%%%%%%%%%%%%%
  %% Subsubsection: Case $\Sigma = \Torus^2$
  %%%%%%%%%%%%%%%%%%%%%%%%%%%%%%%%%%%%%%%%%%%%%%%%%%%%%%%%%%%%%%%%%%%%%%%%%%%%%%%%
  \subsubsection{Case $\Sigma = \Torus^2$}
  
  For the $2$-torus,
  $\pi_1(\Torus^2) = \ZZ^2$,
  and $\pi_k(\Torus^2) = 0$ for all $k \geq 2$.
  
  Thus:
  \begin{itemize}[nosep]
    \item[--] $\pi_2(X) = \pi_3(\Torus^2) = 0$.
    There are no spherical periods.
    \item[--] $\pi_1(X)$ is determined by $\pi_2(Y)$ and the winding,
    while $\pi_0(X) = \pi_1(Y)$ (one component,
    with non-trivial fundamental group).
    \item[--] All periods of $\omega$ arise from
    the toric and surfacic contributions (Part II).
  \end{itemize}
  
  %%%%%%%%%%%%%%%%%%%%%%%%%%%%%%%%%%%%%%%%%%%%%%%%%%%%%%%%%%%%%%%%%%%%%%%%%%%%%%%%
  %% Subsubsection: Case $\Sigma = \Sigma_g$, $g > 1$
  %%%%%%%%%%%%%%%%%%%%%%%%%%%%%%%%%%%%%%%%%%%%%%%%%%%%%%%%%%%%%%%%%%%%%%%%%%%%%%%%
  \subsubsection{Case $\Sigma = \Sigma_g$, $g > 1$}
  
  For a hyperbolic surface of genus $g > 1$,
  $\pi_1(\Sigma_g)$ is non-abelian,
  and $\pi_k(\Sigma_g) = 0$ for all $k \geq 2$.
  
  Thus:
  \begin{itemize}[nosep]
    \item[--] $\pi_2(X) = \pi_3(\Sigma_g) = 0$.
    No spherical periods.
    \item[--] $\pi_0(X) = \pi_1(Y)$,
    and $\pi_1(X)$ is determined by $\pi_2(Y)$ (which is zero)
    and the winding.
    \item[--] The periods of $\omega$ are purely surfacic,
    governed by the non-abelian fundamental group
    and the cocycle $\tau$ constructed in Part II.
  \end{itemize}
  
  %%%%%%%%%%%%%%%%%%%%%%%%%%%%%%%%%%%%%%%%%%%%%%%%%%%%%%%%%%%%%%%%%%%%%%%%%%%%%%%%
  %% Subsection: Summary
  %%%%%%%%%%%%%%%%%%%%%%%%%%%%%%%%%%%%%%%%%%%%%%%%%%%%%%%%%%%%%%%%%%%%%%%%%%%%%%%%
  \subsection{Summary}
  
  The homotopy of the loop space $X = \Loops(Y)$ is summarized in the following table.
  
  \begin{center}
    \begin{tabular}{|c|c|c|c|c|}
      \hline
      $\Sigma$ & $\pi_0(X) = \pi_1(Y)$ & $\pi_1(X)$ & $\pi_2(X)$ & Period type \\
      \hline
      $\Sph^2$ (Hopf, $k=1$) & $0$ & $0$ & $\ZZ$ & spherical \\
      $\Sph^2$ ($\SO(3)$, $k=2$) & $\ZZ_2$ & $0$ & $\ZZ$ & spherical \\
      $\Sph^2$ (lens, $k \geq 3$) & $\ZZ_k$ & $0$ & $\ZZ$ & spherical \\
      $\Torus^2$ & depends on $Y$ & NAb & $0$ & toric + surfacic \\
      $\Sigma_g$ & depends on $Y$ & NAb & $0$ & surfacic \\
      \hline
    \end{tabular}
  \end{center}
  
  The period group $\Pomega$ will be computed in detail
  for $\Sph^2$ in Section~\ref{sec:periods-sphere},
  where the two bundles $Y = \Sph^3$ and $Y = \SO(3)$
  are treated explicitly.
  The cases of $\Torus^2$ and $\Sigma_g$ (where NAb is for Non Abelian),
  which involve the surfacic cocycle of Part II,
  are outlined in Section~\ref{sec:generalization}
  and will be developed fully in a subsequent work.
  
  %%%%%%%%%%%%%%%%%%%%%%%%%%%%%%%%%%%%%%%%%%%%%%%%%%%%%%%%%%%%%%%%%%%%%%%%%%%%%%%%
  %%
  %% Mark: Section: The Prequantum Groupoid
  %%
  %%%%%%%%%%%%%%%%%%%%%%%%%%%%%%%%%%%%%%%%%%%%%%%%%%%%%%%%%%%%%%%%%%%%%%%%%%%%%%%%
  \section{The Prequantum Groupoid}
  
  Geometric Quantization by Paths is a procedure that associates to every closed $2$-form $\omega$ on a diffeological space $X$
  --- a system $(X,\omega)$ called a \emph{parasymplectic space} ---
  a prequantum groupoid $\TT_\omega$,
  quotient of the space of paths $\Paths(X)$ by an equivalence relation.
  Two paths with the same endpoints are identified when the \emph{symplectic area} enclosed between them
  is an element of the \emph{group of periods} $\Pomega \subset \RR$.
  
  As established in Parts I and II,
  the group $\Pomega$ is constituted of three kinds of periods:
  \begin{enumerate}[nosep]
    \item \textbf{The spherical periods:}
    the periods of the closed $1$-form $\Komega$ restricted to $L_0$,
    the connected component of constant loops in $X$.
    They are a homomorphic image of $\pi_2(X)$.
    \item \textbf{The toric periods:}
    the group generated by the periods of $\Komega$ restricted to each connected component $L_i$ of $\Loops(X)$.
    \item \textbf{The surfacic periods:}
    the periods arising from the non-commutativity
    of the fundamental group $\FundamentalGroup(X)$,
    via the surfacic cocycle constructed in Part II.
  \end{enumerate}
  
  In the special case of a simply connected space $X$,
  as for the sphere $\Sph^2$,
  only spherical periods appear.
  When the fundamental group of $X$ is abelian,
  as for the torus $\Torus^2$,
  the group of periods contains spherical and toric contributions.
  For a non-abelian fundamental group,
  as for surfaces $\Sigma_g$ of genus $g > 1$,
  all three kinds of periods are taken into account.
  
  The space of morphisms $\cY = \Mor(\TT_\omega)$ is equipped with a \emph{prequantum $1$-form} $\blambda$,
  the reduction of the action $1$-form $\Komega$ on $\Paths(X)$.
  This form is left-right invariant,
  and its curvature is $\der\blambda = \Trg^*\omega - \Src^*\omega$.
  
  The prequantum groupoid $(\TT_\omega, \blambda)$ is called the \emph{Quantum System}.
  It contains the entire quantum structure of the classical system $(X,\omega)$.
  
  This approach to quantum mechanics consists in reconstructing the quantum system
  from its classical model.
  The space of objects $X$ is the \emph{Skeleton}:
  it represents the classical states,
  embedded as the identity morphisms.
  The non-identity morphisms form the \emph{Quantum Fog}:
  they are the \emph{Quantum Transitions},
  equivalence classes of paths enclosing a symplectic area equal to a period.
  
  %%%%%%%%%%%%%%%%%%%%%%%%%%%%%%%%%%%%%%%%%%%%%%%%%%%%%%%%%%%%%%%%%%%%%%%%%%%%%%%%
  %%
  %% MARK: Section: The Case $\SU(2) \simeq \Sph^3$
  %%
  %%%%%%%%%%%%%%%%%%%%%%%%%%%%%%%%%%%%%%%%%%%%%%%%%%%%%%%%%%%%%%%%%%%%%%%%%%%%%%%%
  \section{The Case $\SU(2) \simeq \Sph^3$}
  
  We parametrize $\Sph^3$ as the unit sphere in $\CC^2$.
  Let $Z = (z_1, z_2)$,
  $\barZ = (\bar{z}_1, \bar{z}_2)$,
  and $\|Z\|^2 = \barZ Z = \sum_{k=1}^2 \bar{z}_k z_k$.
  Then:
  \[
    \Sph^3 = \{ Z \in \CC^2 \mid \|Z\|^2 = 1 \}.
    \]
  As ususal we identify $\Sph^3$ with the group $\SU(2)$ by the map
  $$
    Z = (z_1, z_2) \longmapsto \uline{Z} = \Matrix{z_1 & - \bar z_2 \\ \bar z_2 & \bar z_1}
    $$
  
  %%%%%%%%%%%%%%%%%%%%%%%%%%%%%%%%%%%%%%%%%%%%%%%%%%%%%%%%%%%%%%%%%%%%%%%%%%%%%%%%
  %% Subsection: The Hopf Fibration
  %%%%%%%%%%%%%%%%%%%%%%%%%%%%%%%%%%%%%%%%%%%%%%%%%%%%%%%%%%%%%%%%%%%%%%%%%%%%%%%%
  \subsection{The Hopf Fibration}
  
  The group $\U(1)$ acts on $\Sph^3$ by scalar multiplication:
  \[
    e^{i\theta} \cdot Z = e^{i\theta} Z, % = (e^{i\theta} z_1, e^{i\theta} z_2),
    \qquad \theta \in \RR / 2\pi\ZZ.
    \]
  This action is free and proper,
  and makes $\Sph^3$ a principal $\U(1)$-bundle over the projective line
  $\CC\mathrm{P}^1 \simeq \Sph^2$.
  The projection is:
  \[
    \pr: \Sph^3 \longrightarrow \CC\mathrm{P}^1,
    \qquad
    \pr(Z) = [z_1 : z_2].
    \]
  %  On the matrix $\uline{Z}$ this action corresponds to the multiplication\n
  %    $$\n
  %    e^{i\theta} Z \longmapsto \uline{e^{i\theta} Z} = \uline{e^{i\theta}} \cdot  \uline{Z} % = \Matrix{e^{i\theta} & 0 \\ 0 & e^{-i\theta}} \Matrix{z_1 & - z_2 \\ \bar z_2 & \bar z_1}.\n
  %    $$\n
  
  %%%%%%%%%%%%%%%%%%%%%%%%%%%%%%%%%%%%%%%%%%%%%%%%%%%%%%%%%%%%%%%%%%%%%%%%%%%%%%%%
  %% Subsection: The Connection Form
  %%%%%%%%%%%%%%%%%%%%%%%%%%%%%%%%%%%%%%%%%%%%%%%%%%%%%%%%%%%%%%%%%%%%%%%%%%%%%%%%
  \subsubsection{The Connection Form}
  
  The standard connection $1$-form on this bundle is:
  \[
    \alpha = \frac{1}{2i} \bigl( \barZ dZ - Z d\barZ \bigr)
    = \frac{1}{2i} \sum_{k=1}^2 \bigl( \bar{z}_k dz_k - z_k d\bar{z}_k \bigr).
    \]
  
  This form is:
  \begin{itemize}[nosep]
    \item \textbf{Invariant:}
    it is invariant under the $\U(1)$ action,
    because the phase factors $e^{-i\theta}e^{i\theta}$ cancel in each term.
    \item \textbf{Calibrated:}
    for the orbit map $\hat{Z}(\theta) = e^{i\theta} Z$,
    one has $\hat{Z}^*(\alpha) = d\theta$.
    Thus the integral of $\alpha$ over any fiber is:
    \[
      \int_{\text{fiber}} \alpha = 2\pi.
      \]
  \end{itemize}
  
  %%%%%%%%%%%%%%%%%%%%%%%%%%%%%%%%%%%%%%%%%%%%%%%%%%%%%%%%%%%%%%%%%%%%%%%%%%%%%%%%
  %% Subsection: Curvature and the Symplectic Form on \CPOne
  %%%%%%%%%%%%%%%%%%%%%%%%%%%%%%%%%%%%%%%%%%%%%%%%%%%%%%%%%%%%%%%%%%%%%%%%%%%%%%%%
  \subsubsection{Curvature and the Symplectic Form on $\CPOne$}
  
  The curvature of $\alpha$ is:
  \[
    d\alpha = \frac{1}{i} d\barZ \wedge dZ
    = \frac{1}{i} \bigl( d\bar{z}_1 \wedge dz_1 + d\bar{z}_2 \wedge dz_2 \bigr).
    \]
  
  This $2$-form is basic:
  it is the pullback of a $2$-form on the base $\CPOne$.
  The curvature is proportional to the Fubini--Study form.
  Let $\Surf \in \DForms^2(\CC\mathrm{P}^1)$ be the area form on $\CPOne$
  normalized such that $\int_{\CC\mathrm{P}^1} \Surf = 4\pi$.
  We have:
  \[
    d\alpha = \frac{1}{2} \pr^*\Surf.
    \]
  
  Thus,
  the symplectic form on the base for which $\alpha$ is a prequantum connection is:
  \[
    \epsilon = \frac{1}{2} \Surf,
    \qquad \int_{\CC\mathrm{P}^1} \epsilon = 2\pi.
    \]
  
  %%%%%%%%%%%%%%%%%%%%%%%%%%%%%%%%%%%%%%%%%%%%%%%%%%%%%%%%%%%%%%%%%%%%%%%%%%%%%%%%
  %% Subsection: The Contact Volume Form
  %%%%%%%%%%%%%%%%%%%%%%%%%%%%%%%%%%%%%%%%%%%%%%%%%%%%%%%%%%%%%%%%%%%%%%%%%%%%%%%%
  \subsubsection{The Contact Volume Form}
  
  The volume form on $\Sph^3$ is:
  \[
    \Vol = \alpha \wedge d\alpha.
    \]
  Its total volume is:
  \[
    \int_{\Sph^3} \Vol
    = \left( \int_{\text{fiber}} \alpha \right) \cdot \left( \int_{\CC\mathrm{P}^1} \epsilon \right)
    = 2\pi \cdot 2\pi
    = 4\pi^2.
    \]
  
  Thus,
  with the standard connection $\alpha$,
  the volume of $\Sph^3$ is $4\pi^2$.
  
  %%%%%%%%%%%%%%%%%%%%%%%%%%%%%%%%%%%%%%%%%%%%%%%%%%%%%%%%%%%%%%%%%%%%%%%%%%%%%%%%
  %% Subsection: The Loop Space and Its Period Group
  %%%%%%%%%%%%%%%%%%%%%%%%%%%%%%%%%%%%%%%%%%%%%%%%%%%%%%%%%%%%%%%%%%%%%%%%%%%%%%%%
  \subsubsection{The Loop Space and Its Period Group}
  
  The homotopy groups of $\Sph^3$ are:
  \[
    \pi_1(\Sph^3) = 0,\qquad
    \pi_2(\Sph^3) = 0,\qquad
    \pi_3(\Sph^3) = \ZZ.
    \]
  We denote the loop space by
  $$
    X = \Loops(\Sph^3) = \{ x \in \Paths(\Sph^3) \mid x(0) = x(1) \}.
    $$
  On $X$ is defined the closed $2$-form $\omega$,
  $$
    \omega = \KVol|_X \in \Omega^2(X),
    \qtext{and}
    d\omega = 0.
    $$
  The space $X$ is connected and simply connected:
  \[
    \pi_0(X) = 0,\qquad
    \pi_1(X) = 0,\qquad
    \pi_2(X) = \pi_3(\Sph^3) = \ZZ.
    \]
  The period group $\Pomega$ of $\omega = \KVol|_X$ is generated by
  the integral of $\Vol$ over the fundamental $3$-cycle of $\Sph^3$,
  evaluated via the evaluation map $\Eval: X \times \Sph^1 \to \Sph^3$.
  Hence:
  \[
    \Pomega = 4\pi^2 \ZZ,
    \qquad
    \Tomega = \RR / 4\pi^2 \ZZ.
    \]
  
  %%%%%%%%%%%%%%%%%%%%%%%%%%%%%%%%%%%%%%%%%%%%%%%%%%%%%%%%%%%%%%%%%%%%%%%%%%%%%%%%
  %% Subsection: The Prequantum Groupoid
  %%%%%%%%%%%%%%%%%%%%%%%%%%%%%%%%%%%%%%%%%%%%%%%%%%%%%%%%%%%%%%%%%%%%%%%%%%%%%%%%
  \subsubsection{The Prequantum Groupoid}
  
  Applying the GQbP construction to $(X, \omega)$,
  we obtain the canonical prequantum groupoid $\TT_\omega$.
  Since $X$ is connected and simply connected,
  the construction is that of Part I.
  
  The equivalence relation on $\Paths(X)$ is:
  \[
    \gamma \sim_\omega \gamma'
    \iff
    \Ends(\gamma) = \Ends(\gamma')
    \;\text{and}\;
    \int_\gamma^{\gamma'} \Komega \in 4\pi^2 \ZZ.
    \]
  The space of morphisms $\cY = \Paths(X)/\!\sim_\omega$ forms a fibrating diffeological groupoid
  with isotropy:
  \[
    \TT_{\omega, \ell} \simeq \Tomega = \RR / 4\pi^2 \ZZ,
    \]
  for any loop $\ell \in X$.
  
  The prequantum $1$-form $\blambda \in \DForms^1(\cY)$ satisfies:
  \[
    \Class_\omega^*\blambda = \Komega,
    \qquad
    d\blambda = \Trg^*\omega - \Src^*\omega.
    \]
  The source fiber $\cY_{\ell_0} = \Mor_{\TT_\omega}(\ell_0, \ast)$ is a principal $\Tomega$-bundle over $X$,
  equipped with a connection form $\lambda = \blambda|_{\cY_{\ell_0}}$ of curvature $\omega$.
  They are the family of prequantum bundles involved in Kostant-Souriau geometric quantization.
  
  %%%%%%%%%%%%%%%%%%%%%%%%%%%%%%%%%%%%%%%%%%%%%%%%%%%%%%%%%%%%%%%%%%%%%%%%%%%%%%%%
  %%
  %% Subsection: The Action of $\SU(2)$ on $\Sph^3$
  %%
  %%%%%%%%%%%%%%%%%%%%%%%%%%%%%%%%%%%%%%%%%%%%%%%%%%%%%%%%%%%%%%%%%%%%%%%%%%%%%%%%
  \subsection{The Action of $\SU(2)$ on $\Sph^3$}
  
  The group $\SU(2)$ acts naturally on $\Sph^3 \subset \CC^2$ by left matrix multiplication.
  Let $g \in \SU(2)$ and $Z \in \Sph^3$.
  The action is given by the standard matrix-vector product:
  \[
    Z \longmapsto g Z.
    \]
  
  Let us check that this action preserves the connection $1$-form $\alpha$.
  Recall that $\alpha = \frac{1}{2i} ( \barZ dZ - Z d\barZ )$.
  Pulling back $\alpha$ by the action of $g$,
  and using the fact that $g^\dagger g = \id$ (where $g^\dagger = \bar{g}^T$),
  we obtain:
  \begin{align*}
    g^*\alpha &= \frac{1}{2i} \bigl( \overline{(gZ)}^T d(gZ) - (gZ)^T d\overline{(gZ)} \bigr) \\
    &= \frac{1}{2i} \bigl( \barZ g^\dagger g dZ - Z^T g^T \bar{g} d\barZ \bigr) \\
    &= \frac{1}{2i} \bigl( \barZ dZ - Z d\barZ \bigr) \\
    &= \alpha.
  \end{align*}
  Since $\alpha$ is strictly left-invariant under $\SU(2)$,
  its exterior derivative $d\alpha$ is also left-invariant.
  Consequently,
  the contact volume form $\Vol = \alpha \wedge d\alpha$ is strictly preserved by the action of $\SU(2)$:
  \[
    g^*\Vol = \Vol.
    \]
  
  %%%%%%%%%%%%%%%%%%%%%%%%%%%%%%%%%%%%%%%%%%%%%%%%%%%%%%%%%%%%%%%%%%%%%%%%%%%%%%%%
  %%
  %% Subsection: Variance of $\omega$ and $\Komega$ Under the Loop Group Action
  %%
  %%%%%%%%%%%%%%%%%%%%%%%%%%%%%%%%%%%%%%%%%%%%%%%%%%%%%%%%%%%%%%%%%%%%%%%%%%%%%%%%
  \subsection{Variance of $\omega$ and $\Komega$ Under the Loop Group Action}
  
  We now lift the action of $\SU(2)$ to the loop space.
  Let $\LG = \Loops(\SU(2))$ be the loop group.
  It acts on $X = \Loops(\Sph^3)$ by pointwise multiplication.
  For a loop $\ell \in \LG$ and a loop $x \in X$,
  the action $\u(\ell): X \to X$ is defined by:
  \[
    \u(\ell)(x) = [t \mapsto \ell(t) x(t)].
    \]
  So,
  $\u(\ell)$ is the pointwise multiplication of loops.
  
  Now,
  let $\u(\ell)_{\,*}$ be the natural induced action of $\LG$ on $\Paths(X)$.
  For any path $\gamma: s \mapsto x_s$ in $X$,
  \[
    {\u(\ell)}_{\,*}(\gamma) = \u(\ell) \circ \gamma = [s \mapsto [t \mapsto \ell(t) x_s(t)]].
    \]
  Thanks to the variance of the chain-homotopy operator \cite[\textsection 6.84]{PIZ13},
  we know that:
  \[
    \CHO \circ\; \u(\ell)^* = (\u(\ell)_{\,*})^* \circ \CHO
    \qtext{and then:} \CHO(\u(\ell)^*\omega) = (\u(\ell)_{\,*})^*\Komega.
    \]
  Therefore,
  to know the variance $(\u(\ell)_{\,*})^*\Komega$ of the \emph{action form} $\Komega$ under the action of the loop group $\LG$ on $\Paths(X)$,
  we need to know the variance of $\omega$.
  For that,
  we will use the same technique that is used to show the homotopy invariance of the de Rham cohomology in diffeology \cite[\textsection 6.88]{PIZ13}.
  
  Since the loop group $\LG$ is connected (because $\pi_0(\LG) \simeq \pi_1(\SU(2)) = 0$),
  there exists a smooth path $s \mapsto \ell_s$ in $\LG$ such that $\ell_0 = \hat \id$ (the constant loop at the identity of $\SU(2)$) and $\ell_1 = \ell$.
  Hence,
  let us define, for all $x \in X$:
  \[
    F_s(x) = \u(\ell_s)(x) = [t \mapsto \ell_s(t) x(t)].
    \]
  Then:
  \[
    F_0(x) = x \qtext{and} F_1(x) = \u(\ell)(x).
    \]
  Define now the map
  \[
    \Psi : X \to \Paths(X)
    \qtext{by}
    \Psi(x) = [s \mapsto F_s(x)].
    \]
  We pull back the fundamental identity of the chain-homotopy operator,
  $\CHO \der + \der \CHO = \1^* - \0^*$,
  applied to the parasymplectic form $\omega$.
  Since $\der\omega = 0$,
  this gives:
  \begin{align*}
    \Psi^*(\CHO(\der\omega)) + \Psi^*(\der\Komega) &= \Psi^*(\1^*\omega) - \Psi^*(\0^*\omega) \\
    \der[\Psi^*(\Komega)] &= (\1 \circ \Psi)^*\omega - (\0 \circ \Psi)^*\omega \\
    &= F_1^*\omega - F_0^*\omega \\
    &= \u(\ell)^*\omega - \omega.
  \end{align*}
  That is,
  \[
    \u(\ell)^*\omega = \omega + \der\theta
    \qtext{with} \theta = \Psi^*(\Komega) \in \Omega^1(X).
    \]
  
  Now,
  the variance of $\Komega \in \Omega^1(\Paths(X))$ the lift $\u(\ell)_{\,*}$ where $\ell \in \LG$ is just given by applying the variance identity of the chain-homotopy operator:
  $$
    (\u(\ell)_{\,*})^*(\Komega) = \CHO(\u(\ell)^* \omega) = \CHO(\omega + \der\theta) = \Komega + \CHO(\der\theta),
    $$
  That is,
  \[ \tag{$\heartsuit$}
    (\u(\ell)_{\,*})^*(\Komega)  = \Komega + \1^*(\theta) - \0^*(\theta) - \der[\Ktheta],
    \qtext{with}
    \Ktheta(\gamma) = \int_\gamma \theta.
    \]
  
  \textbf{Remark on the canonical choice of $\theta$}.
  The $1$-form $\theta = \Psi^*(\Komega)$ provides the exact shift required for the fusion twist.
  At first glance,
  this specific form appears to depend on the arbitrary choice of the path $s \mapsto \ell_s$.
  However,
  in the special case of a group,
  $\SU(2)$ or $\SO(3)$,
  for the contact manifold $Y$ over the surface,
  that is,
  $\Sph^2$ for our concern,
  as demonstrated in Appendix~\ref{app:canonical_pw},
  an explicit expansion of $\Psi^*(\Komega)$ on a plot of $X$ reveals that this $1$-form decomposes exactly as:
  \[
    \Psi^*(\Komega) = \theta_\ell + \der f,
    \qtext{with}
    \theta_\ell(\delta x) = \int_0^1 \Tr\bigl( (\ell^{-1} \dot{\ell}) \cdot (x^{-1} \delta x) \bigr)\, \dt
    \]
  is the canonical Polyakov--Wiegmann $1$-form,
  which depends \emph{only} on the boundary loop $\ell$,
  and $\der f$ is an exact differential that absorbs all the dependence on the chosen homotopy.
  Because shifting by an exact differential corresponds merely to a canonical gauge isomorphism in the prequantum groupoid,
  we can canonically choose $\theta = \theta_\ell$.
  Therefore the variance of $\omega$ and $\Komega$ can be written specifically:
  \[ \tag{$\diamondsuit$}
    \u(\ell)^*\omega = \omega + \der\theta_\ell,
    \qtext{and}
    (\u(\ell)_{\,*})^*(\Komega)  = \Komega + \1^*(\theta) - \0^*(\theta) - \der[\Ktheta_\ell],
    \]
  This ensures that the geometric shift is a rigid,
  objective invariant of the loop $\ell$,
  perfectly bridging the macroscopic diffeological operator $\CHO$ with the microscopic Maurer--Cartan algebra.
  
  %%%%%%%%%%%%%%%%%%%%%%%%%%%%%%%%%%%%%%%%%%%%%%%%%%%%%%%%%%%%%%%%%%%%%%%%%%%%%%%%
  %%
  %% Subsection: Lifting the Action to the Prequantum Groupoid
  %%
  %%%%%%%%%%%%%%%%%%%%%%%%%%%%%%%%%%%%%%%%%%%%%%%%%%%%%%%%%%%%%%%%%%%%%%%%%%%%%%%%
  \subsection{Lifting the Action to the Prequantum Groupoid}
  
  We now lift the action of the loop group $\LG$ to the space of morphisms $\cY = \Mor(\TT_\omega)$.
  Let $y = [\gamma]_\omega \in \cY$ be a quantum transition.
  The direct action of $\LG$ on $\cY$ is the projection of the pointwise path multiplication:
  $y \mapsto [\u(\ell) \circ \gamma]_\omega$.
  However,
  as equation ($\diamondsuit$) demonstrates,
  this direct action does not strictly preserve the prequantum potential $\blambda$.
  It shifts the connection by boundary terms and an exact differential $\der[\Ktheta_\ell]$.
  
  To construct a true automorphism of the prequantum groupoid,
  the direct action must be refined to account for this geometric shift.
  We define the lifted action $\hat{\u(\ell)}: \cY \to \cY$ by translating the morphism in its fiber:
  \[
    \hat{\u(\ell)}([\gamma]_\omega) = [\u(\ell) \circ \gamma]_\omega \oplus \bigl( - \Ktheta_\ell(\gamma) \bigr),
    \]
  where $\oplus$ denotes the action of the structure group $\Tomega = \RR/\Pomega$ on the fibers of $\cY$.
  Because $X = \Loops(\Sph^3)$ is simply connected,
  the integral $\Ktheta_\ell(\gamma) = \int_\gamma \theta_\ell$ is independent of the choice of representative in the homotopy class of $\gamma$.
  Thus,
  the gauge shift is a well-defined function on $\cY$,
  and the lifted action is rigorously established.
  
  By incorporating the translation by $- \Ktheta_\ell$,
  we add a term $-\der[\Ktheta_\ell]$ to the pullback of the connection.
  Combining this with the variance of the direct action in equation ($\diamondsuit$),
  the pullback of the prequantum potential under the lifted action is:
  \[
    \hat{\u(\ell)}^*\blambda = \blambda + \Trg^*\theta_\ell - \Src^*\theta_\ell - 2\der[\Ktheta_\ell].
    \]
  Taking the exterior derivative,
  we find that the curvature $\bar{\omega} = \der\blambda$ on $\cY$ satisfies:
  \[
    \hat{\u(\ell)}^*\bar{\omega} = \bar{\omega} + \der(\Trg^*\theta_\ell - \Src^*\theta_\ell).
    \]
  The curvature is strictly preserved in the bulk of the morphisms,
  shifting only by exact forms pulled back from the object space $X$.
  Thus,
  $\LG$ acts on the prequantum groupoid by strict parasymplectomorphisms of the upstairs curvature $\bar{\omega}$.
  
  \begin{proof}[Well-Definedness]
    We must verify that the lifted action $\hat{\u(\ell)}$ is independent of the choice of representative in the equivalence class $y = [\gamma]_\omega$.
    Suppose $\gamma \sim_\omega \gamma'$.
    This implies there exists a cylinder $C$ such that $\Ends(C) = (\gamma, \gamma')$ and $\int_C \omega \in \Pomega$.
    The difference between the lifted actions on the two representatives is:
    \[
      \Delta = \left( \int_{\u(\ell)\gamma} \Komega - \Ktheta_\ell(\gamma) \right) - \left( \int_{\u(\ell)\gamma'} \Komega - \Ktheta_\ell(\gamma') \right).
      \]
    Using the variance formula $\u(\ell)^*\Komega = \Komega + \CHO(\der\theta_\ell)$,
    and the identity $\CHO(\der\theta_\ell) = \1^*\theta_\ell - \0^*\theta_\ell - \der \Ktheta_\ell$,
    the difference in the first terms is:
    \[
      \int_{\u(\ell)C} \omega = \int_C \omega + \int_C \der\theta_\ell = \int_C \omega + \bigl( \Ktheta_\ell(\gamma) - \Ktheta_\ell(\gamma') \bigr).
      \]
    Substituting this into the expression for $\Delta$:
    \[
      \Delta = \int_C \omega + \bigl( \Ktheta_\ell(\gamma) - \Ktheta_\ell(\gamma') \bigr) - \bigl( \Ktheta_\ell(\gamma) - \Ktheta_\ell(\gamma') \bigr) = \int_C \omega.
      \]
    Since $\int_C \omega \in \Pomega$,
    the two results are identical in the quotient $\cY$.
    Thus,
    while the direct action on the path space is ill-defined on the quotient,
    the lifted action $\hat{\u(\ell)}$ is a rigorously well-defined automorphism of the prequantum groupoid.
  \end{proof}
  
  %%%%%%%%%%%%%%%%%%%%%%%%%%%%%%%%%%%%%%%%%%%%%%%%%%%%%%%%%%%%%%%%%%%%%%%%%%%%%%%%
  %%
  %% Subsection: The Intrinsic Moment Map and the Souriau Cocycle
  %%
  %%%%%%%%%%%%%%%%%%%%%%%%%%%%%%%%%%%%%%%%%%%%%%%%%%%%%%%%%%%%%%%%%%%%%%%%%%%%%%%%
  \subsection{The Intrinsic Moment Map and the Souriau Cocycle}
  
  Because the loop group $\LG$ acts by parasymplectomorphisms on $(\cY, \bar{\omega})$,
  it admits an intrinsic diffeological moment map.
  Let $\LieAlgebra{lg}$ be the Lie algebra of $\LG$,
  and $\LieAlgebra{lg}^*$ its diffeological dual.
  For a morphism $y \in \cY$,
  let $\hat{y}: \LG \to \cY$ be the orbit map $\hat{y}(\ell) = \hat{\u(\ell)}(y)$.
  The moment map $\Psi: \cY \to \LieAlgebra{lg}^*$ is defined canonically by the pullback of the prequantum potential:
  \[
    \Psi(y) = \hat{y}^*(\blambda).
    \]
  
  This moment map evaluates the transfer of conserved quantities during the quantum transition $y$.
  However,
  because the lifted action incorporates the gauge shift $-\Ktheta_\ell$,
  the moment map fails to be strictly equivariant under the coadjoint action of $\LG$.
  This failure of equivariance defines the Souriau cocycle $c: \LG \to \LieAlgebra{lg}^*$:
  \[
    \Psi(\hat{\u(\ell)}(y)) = \Ad_\ell^* \Psi(y) + c(\ell).
    \]
  
  %%%%%%%%%%%%%%%%%%%%%%%%%%%%%%%%%%%%%%%%%%%%%%%%%%%%%%%%%%%%%%%%%%%%%%%%%%%%%%%%
  %%
  %% Subsection: The Kac-Moody Central Extension
  %%
  %%%%%%%%%%%%%%%%%%%%%%%%%%%%%%%%%%%%%%%%%%%%%%%%%%%%%%%%%%%%%%%%%%%%%%%%%%%%%%%%
  \subsection{The Kac-Moody Central Extension}
  
  The infinitesimal version of the Souriau cocycle is a Lie algebra $2$-cocycle $\Theta: \LieAlgebra{lg} \times \LieAlgebra{lg} \to \RR$,
  obtained by differentiating the equivariance defect at the identity.
  Let $\xi, \eta \in \LieAlgebra{lg}$ be two infinitesimal generators.
  The standard formula for the obstruction cocycle of an anomalous action satisfying $\mathcal{L}_{\tilde{\xi}} \omega = \der\alpha_\xi$ is:
  \[
    \Theta(\xi, \eta) = \iota_{\tilde{\xi}} \alpha_\eta - \iota_{\tilde{\eta}} \alpha_\xi + \omega(\tilde{\xi}, \tilde{\eta}).
    \]
  
  We evaluate this cocycle at the constant loop $x = \id$.
  First,
  because the velocity of the constant loop is zero ($\dot{x} = 0$),
  the parasymplectic form vanishes identically: $\omega_{\id}(\tilde{\xi}, \tilde{\eta}) = 0$.
  
  Second,
  we compute the anomaly $1$-forms $\alpha_\xi$.
  From the explicit formula for the shift $\theta_{\ell_s}$ derived in Appendix~\ref{app:canonical_pw},
  differentiating with respect to $s$ at $s=0$ yields:
  \[
    \alpha_\xi(\delta x) = \int_0^1 \Tr\bigl( \dot{\xi}(t) \cdot (x(t)^{-1} \delta x(t)) \bigr) \, \dt.
    \]
  Evaluating this at $x = \id$ on the variation $\delta x = \tilde{\eta}$,
  we obtain:
  \[
    \iota_{\tilde{\eta}} \alpha_\xi = \alpha_\xi(\tilde{\eta}) = \int_0^1 \Tr\bigl( \dot{\xi}(t) \eta(t) \bigr) \, \dt.
    \]
  
  Substituting these into the obstruction formula gives:
  \[
    \Theta(\xi, \eta) = \int_0^1 \Tr\bigl( \dot{\eta}(t) \xi(t) \bigr) \, \dt - \int_0^1 \Tr\bigl( \dot{\xi}(t) \eta(t) \bigr) \, \dt.
    \]
  Using integration by parts,
  and noting that the boundary terms vanish because the vector fields are loops ($\xi(0)=\xi(1)$ and $\eta(0)=\eta(1)$),
  we have $\int_0^1 \Tr(\dot{\eta} \xi) \dt = - \int_0^1 \Tr(\eta \dot{\xi}) \dt$.
  This yields exactly:
  \[
    \Theta(\xi, \eta) = 2 \int_0^1 \Tr\bigl( \xi(t) \dot{\eta}(t) \bigr) \, \dt.
    \]
  
  \begin{theorem}
    The infinitesimal Souriau cocycle $\Theta(\xi, \eta)$ of the prequantum groupoid coincides exactly with the standard Kac-Moody central extension cocycle.
  \end{theorem}
  
  For $\SU(2)$,
  the quantization is now complete.
  The Kac-Moody central extension is the canonical geometric consequence of the Polyakov-Wiegmann shift required to lift the loop group action to the prequantum groupoid.
  Because the loop space of $\SU(2)$ is simply connected,
  this shift is globally well-defined,
  and the prequantum groupoid exists without topological obstruction.
  
  %%%%%%%%%%%%%%%%%%%%%%%%%%%%%%%%%%%%%%%%%%%%%%%%%%%%%%%%%%%%%%%%%%%%%%%%%%%%%%%%
  %%
  %% MARK: Section: The Action of the Central Extension
  %%
  %%%%%%%%%%%%%%%%%%%%%%%%%%%%%%%%%%%%%%%%%%%%%%%%%%%%%%%%%%%%%%%%%%%%%%%%%%%%%%%%
  \section{The Action of the Central Extension}
  
  As established in the previous section,
  the loop group $LG$ does not act by parasymplectomorphisms on the space of motions $X$.
  Consequently,
  the direct action of $LG$ on the space of paths does not preserve the equivalence relation $\sim_\omega$,
  and $LG$ cannot be realized as a group of automorphisms of the prequantum groupoid $\TT_\omega$.
  
  To restore the symmetry at the quantum level,
  we must consider the central extension of the loop group.
  This extension,
  denoted by $\widehat{LG}$,
  is the true symmetry group of the quantum system.
  
  \subsection{The Kac-Moody Group}
  The central extension $\widehat{LG}$ is defined as the set of pairs $(\ell, a) \in LG \times \Tomega$,
  where $\Tomega = \RR/\Pomega$ is the quantum phase torus.
  The group law is given by:
  \[
    (\ell, a) \cdot (\ell', a') = (\ell \ell', a + a' + \tau(\ell, \ell')),
    \]
  where $\tau$ is the Polyakov--Wiegmann cocycle.
  
  \subsection{The Action on the Space of Morphisms}
  The group $\widehat{LG}$ acts smoothly on the space of morphisms $\cY = \Mor(\TT_\omega)$.
  For any element $(\ell, a) \in \widehat{LG}$ and any morphism $y = [\gamma]_\omega \in \cY$,
  the action is defined by:
  \[
    \widehat{m}(\ell, a) \cdot [\gamma]_\omega = [\ell \cdot \gamma]_\omega \oplus \left( a - \int_\gamma \theta_\ell \right),
    \]
  where:
  \begin{itemize}[nosep]
    \item $\ell \cdot \gamma$ is the pointwise multiplication of the path $\gamma$ by the loop $\ell$.
    \item $\theta_\ell$ is the Polyakov--Wiegmann $1$-form on the path space.
    \item $\oplus$ denotes the action of the structure group $\Tomega$ on the fibers of the groupoid.
  \end{itemize}
  
  \subsection{Well-Definedness and Symmetry}
  The action $\widehat{m}$ is rigorously well-defined on the quotient $\cY$.
  If we choose a different representative $\gamma' \sim_\omega \gamma$,
  the change in the symplectic integral $\int_{\ell\gamma} \Komega$ is exactly compensated by the change in the gauge term $\int_\gamma \theta_\ell$,
  thanks to the fundamental identity of the chain-homotopy operator.
  
  Furthermore,
  this action preserves the curvature $\bar{\omega} = \der\blambda$ on the space of morphisms up to boundary terms:
  \[
    \widehat{m}(\ell, a)^* \blambda = \blambda + \Trg^*\theta_\ell - \Src^*\theta_\ell.
    \]
  
  \begin{theorem}
    The central extension $\widehat{LG}$ acts as a group of automorphisms of the prequantum groupoid $\TT_\omega$.
    The ``obstruction'' to the loop group action is thus resolved by the geometry of the space of morphisms,
    identifying the Kac-Moody group as the unique and natural symmetry group of the distilled quantum transitions.
  \end{theorem}
  
  %%%%%%%%%%%%%%%%%%%%%%%%%%%%%%%%%%%%%%%%%%%%%%%%%%%%%%%%%%%%%%%%%%%%%%%%%%%%%%%%
  %%
  %% MARK: Example on the space of loops of the $2$-sphere
  %%
  %%%%%%%%%%%%%%%%%%%%%%%%%%%%%%%%%%%%%%%%%%%%%%%%%%%%%%%%%%%%%%%%%%%%%%%%%%%%%%%%
  
  \section{Example on the space of loops of the $2$-sphere}
  
  Let $P : U \to X$ be an $n$-plot of $X = \Loops(S^{2})$, where $U$ is an open domain in $\RR^{n}$. For each $r \in U$, we write $\ell_{r} = P(r) \in \Loops(S^{2})$, so that $P(r)(t) = \ell_{r}(t) \in S^{2}$. The parasymplectic form $\omega$ is defined on $P$ by
  \[
    \omega(P)_{r}(\delta r, \delta' r)
    = \int_{0}^{1} \Surf_{\ell_{r}(t)}\!\left(
    \frac{\partial \ell_{r}(t)}{\partial r}\,\delta r,\;
    \frac{\partial \ell_{r}(t)}{\partial r}\,\delta' r
    \right) dt,
    \]
  for all $r \in U$ and all tangent vectors $\delta r, \delta' r \in \RR^{n}$. Equivalently, using the time evaluation map $\hat{t} : \ell \mapsto \ell(t)$, one has
  \[
    \omega = \int_{0}^{1} \hat{t}^{*}(\Surf),
    \]
  which, when pulled back along the plot $P$, yields the integral formula above.
  
  Let $G = \Loops(\SO(3))$ act on $X = \Loops(S^{2})$ pointwise.
  For $\phi: t \mapsto A_t$ in $G$ and $\ell \in X$, the action is defined by
  \[
    (\phi \cdot \ell)(t) = A_t \ell(t), \quad \forall t \in \RR,
    \]
  where $A_t \ell(t)$ denotes the standard action of $\SO(3)$ on $S^{2} \subset \RR^{3}$.
  The periodicity condition $\phi(0) = \phi(1)$ and $\ell(0) = \ell(1)$ ensures that $\phi \cdot \ell \in X$.
  
  \begin{proposition}{}
    The pointwise action of $G = \Loops(\SO(3))$ on $X = \Loops(S^{2})$ preserves the closed $2$-form
    $\omega = \int_{0}^{1} \hat{t}^{*}(\Surf)$.
    That is, for all $\phi \in G$,
    \[
      \phi^{*}\omega = \omega.
      \]
  \end{proposition}
  
  \begin{proof}
    Let $\phi \in G$ be given by $t \mapsto A_{t}$ with $A_{t} \in \SO(3)$.
    Let $\uline{\phi} : X \to X$ be the action map, $\uline{\phi}(\ell) : t \mapsto A_{t}\,\ell(t)$.
    
    We evaluate the pullback $\uline{\phi}^{*}\omega$ on an arbitrary $n$-plot $P : U \to X$, where $U$ is an open domain in $\RR^{n}$.
    Write $P(r) = \ell_{r}$ for $r \in U$, and recall that $\ell_{r}(t) \in S^{2}$.
    For tangent vectors $\delta r, \delta' r \in \RR^{n}$, we compute:
    \[
      (\uline{\phi}^{*}\omega)(P)_{r}(\delta r, \delta' r)
      = \omega(\uline{\phi} \circ P)_{r}(\delta r, \delta' r).
      \]
    By definition of $\omega$,
    \[
      \omega(\uline{\phi} \circ P)_{r}(\delta r, \delta' r)
      = \int_{0}^{1} \Surf_{A_{t}\,\ell_{r}(t)}\!\left(
      \frac{\partial (A_{t}\,\ell_{r}(t))}{\partial r}\,\delta r,\;
      \frac{\partial (A_{t}\,\ell_{r}(t))}{\partial r}\,\delta' r
      \right) dt.
      \]
    
    Since $A_{t}$ does not depend on $r$, the derivative is simply
    \[
      \frac{\partial (A_{t}\,\ell_{r}(t))}{\partial r}\,\delta r
      = A_{t}\!\left( \frac{\partial \ell_{r}(t)}{\partial r}\,\delta r \right).
      \]
    Thus the integrand becomes
    \[
      \Surf_{A_{t}\,\ell_{r}(t)}\!\left(
      A_{t} v,\; A_{t} v'
      \right),
      \]
    where
    \[
      v = \frac{\partial \ell_{r}(t)}{\partial r}\,\delta r,
      \qquad
      v' = \frac{\partial \ell_{r}(t)}{\partial r}\,\delta' r
      \]
    are tangent vectors to $S^{2}$ at $\ell_{r}(t)$.
    
    Now, the area form $\Surf$ on $S^{2}$ is invariant under the action of $\SO(3)$:
    \[
      (A_{t})^{*}\Surf = \Surf, \quad \forall\, A_{t} \in \SO(3).
      \]
    By definition of the pullback of a differential form,
    \[
      \Surf_{A_{t}\,\ell_{r}(t)}(A_{t} v, A_{t} v') = \Surf_{\ell_{r}(t)}(v, v').
      \]
    
    Substituting back into the integral, we obtain
    \[
      (\uline{\phi}^{*}\omega)(P)_{r}(\delta r, \delta' r)
      = \int_{0}^{1} \Surf_{\ell_{r}(t)}\!\left(
      \frac{\partial \ell_{r}(t)}{\partial r}\,\delta r,\;
      \frac{\partial \ell_{r}(t)}{\partial r}\,\delta' r
      \right) dt
      = \omega(P)_{r}(\delta r, \delta' r).
      \]
    
    Since $P$, $r$, $\delta r$, and $\delta' r$ were arbitrary, we conclude
    \[
      \uline{\phi}^{*}\omega = \omega.
      \]
    
    Thus, every $\phi \in G$ preserves $\omega$, and $G$ is a subgroup of the classical symmetry group $\Diff(X,\omega)$.
  \end{proof}
  
  \begin{proposition}{}
    The group of periods $\Pomega$ of the closed $2$-form $\omega = \int_{0}^{1} \hat{t}^{*}(\Surf)$
    on $X = \Loops(S^{2})$ is
    \[
      \Pomega = 4\pi \ZZ.
      \]
  \end{proposition}
  
  \begin{proof}
    We follow the definitions of GQbP-II.  The group of periods $\Pomega$ is the total group
    of periods, generated by the toric periods and the surfacic periods.  We show that the
    surfacic periods vanish and that the toric periods generate $4\pi \ZZ$.
    
    \medskip
    \noindent\textbf{1.  The connected components of $\Loops(X)$.}
    Let $X = \Loops(S^{2})$.  A loop in $X$ is a smooth map $\sigma : S^{1} \to \Loops(S^{2})$,
    which is equivalent to a smooth map $\Sigma : S^{1} \times S^{1} \to S^{2}$.
    The set of connected components of $\Loops(X)$ is therefore
    \[
      \pi_{0}(\Loops(X)) = \pi_{0}(\Loops(\Loops(S^{2}))) = \pi_{2}(S^{2}).
      \]
    The group $\pi_{2}(S^{2})$ is the set of homotopy classes of \emph{spheroids},
    i.e., smooth maps $f : S^{2} \to S^{2}$ that send the north pole $N \in S^{2}$ to itself.
    These homotopy classes are in bijection with $\ZZ$, the isomorphism being given by the
    topological degree $d \in \ZZ$ of the map $f$:
    \[
      \int_{S^{2}} f^{*}\Surf = 4\pi d.
      \]
    Hence $\Loops(X)$ decomposes into connected components indexed by the degree:
    \[
      \Loops(X) = \coprod_{d \in \ZZ} L_{d},
      \]
    where a loop $\sigma \in \Loops(X)$ belongs to $L_{d}$ if and only if the associated map
    $\Sigma : S^{1} \times S^{1} \to S^{2}$ has topological degree $d$ (when the two circles
    are pinched to a wedge, the resulting spheroid has degree $d$).
    
    \medskip
    \noindent\textbf{2.  Toric periods.}
    The chain-homotopy operator gives a closed $1$-form $\CHO\omega$ on $\Paths(X)$.
    Its restriction to each component $L_{d} \subset \Loops(X)$ defines a local group of periods:
    \[
      P_{d} = \left\{ \int_{\sigma} \CHO\omega \ \middle| \ \sigma \in \Loops(L_{d}) \right\} \subset \RR.
      \]
    Let $\sigma \in \Loops(L_{d})$ be a loop in the component $L_{d}$.
    The integral $\int_{\sigma} \CHO\omega$ evaluates the $2$-form $\omega$ on the $2$-cycle in $X$
    swept out by the one-parameter family of loops $\sigma(s) \in L_{d}$.
    This $2$-cycle corresponds precisely to the map $\Sigma : S^{1} \times S^{1} \to S^{2}$,
    which, after pinching the two circles to a wedge, yields a spheroid $f : S^{2} \to S^{2}$
    of degree $d$.
    By the definition of $\omega$ as the paths pullback of $\Surf$, we have
    \[
      \int_{\sigma} \CHO\omega = \int_{S^{1} \times S^{1}} \Sigma^{*}\Surf = 4\pi d.
      \]
    Since $\sigma$ varies within the fixed component $L_{d}$, every loop in $L_{d}$ yields the
    same value $4\pi d$.  A loop in $L_{d}$ (a loop of loops in $L_{d}$) can change this value
    only by a period of $\CHO\omega$ restricted to $L_{d}$, which is generated by the same
    quantity.  Therefore the local period group is
    \[
      P_{d} = 4\pi d \, \ZZ.
      \]
    The \emph{group of toric periods} is the subgroup of $\RR$ generated by the union of all
    local periods:
    \[
      \Ptor = \left\langle \bigcup_{d \in \ZZ} P_{d} \right\rangle
      = \left\langle \bigcup_{d \in \ZZ} 4\pi d \, \ZZ \right\rangle
      = 4\pi \ZZ.
      \]
    
    \medskip
    \noindent\textbf{3.  Absence of surfacic periods.}
    The fundamental group of $X = \Loops(S^{2})$ is
    \[
      \cI = \pi_{1}(\Loops(S^{2})) \simeq \pi_{2}(S^{2}) \simeq \ZZ.
      \]
    The surfacic cocycle $\tau : \cI \times \cI \to \Ttor$ (where $\Ttor = \RR/\Ptor$)
    measures the symplectic area enclosed by commutators of basis loops.
    For the free loop space of $S^{2}$, the action of $\cI$ on the periods is trivial,
    because the form $\omega$ is the paths pullback of a form on $S^{2}$ and the fundamental
    group acts by reparametrization of loops, which does not change the integral of $\omega$
    over any $2$-cycle.  Consequently, the accumulated cocycle $\cT(w)$ vanishes for every
    relation $w$ in $\cI$, and the subgroup $H_{\Surf} \subset \Ttor$ generated by these
    values is trivial.  Hence the total group of periods $\Pomega$ (as defined in GQbP-II, \S\ref{art:totalgroupperiods}) coincides with the toric periods:
    \[
      \Pomega = \Ptor = 4\pi \ZZ.
      \]
    
    \medskip
    \noindent\textbf{4.  Conclusion.}
    The closed $2$-form $\omega$ is discrete: $\Pomega = 4\pi \ZZ \neq \RR$.
    The torus of periods is $T_{\omega} = \RR / 4\pi \ZZ \simeq \U(1)$.
    The prequantum groupoid $\TT_{\omega}$ therefore exists by the main theorem of GQbP-II.
  \end{proof}
  
%%%%%%%%%%%%%%%%%%%%%%%%%%%%%%%%%%%%%%%%%%%%%%%%%%%%%%%%%%%
%%
%% MARK: Subsection -- Moment Map of $\Loops(\SO(3))$ on $\Loops(S^{2})$
%%
%%%%%%%%%%%%%%%%%%%%%%%%%%%%%%%%%%%%%%%%%%%%%%%%%%%%%%%%%%%

\subsection{The Moment Map of the Loop Group Action}

We have shown that $G = \Loops(\SO(3))$ acts on $(X,\omega)$,
with $X = \Loops(S^{2})$ and $\omega = \int_{0}^{1}\hat{t}^{*}(\Surf)$,
as a group of strict parasymplectomorphisms:
$\phi^{*}\omega = \omega$ for every $\phi \in G$, so $G \subset \Diff(X,\omega)$.
By the main theorem of GQbP-II, this action lifts to a faithful action
of $G$ on the prequantum groupoid $\TT_{\omega}$,
preserving the prequantum $1$-form $\blambda$.
We now compute the associated moment maps and the holonomy group
of the action, following the constructions of \cite{PIZ10}.

%%%%%%%%%%%%%%%%%%%%%%%%%%%%%%%%%%%%%%%%%%%%%%%%%%%%%%%%%%%
\subsubsection{The paths moment map}

For every path $p \in \Paths(X)$, let $\hat{p} : G \to \Paths(X)$ be the orbit map
$\hat{p}(\phi) = \phi \cdot p$.
The \emph{paths moment map} is the smooth map
\[
\Psi : \Paths(X) \to \cG^{*},
\quad
\Psi(p) = \hat{p}^{*}(\CHO\omega),
\]
where $\cG^{*}$ is the space of momenta of $G$ (the space of left-invariant
$1$-forms on $G$), and $\CHO$ is the chain-homotopy operator
(\S\ref{art: The Chain-Homotopy Operator}).

Because the action of $G$ on $(X,\omega)$ is \emph{strict}---the form $\omega$
is exactly invariant, not merely closed-invariant---the primitive $\CHO\omega$
is also strictly invariant under the induced action on $\Paths(X)$:
$\phi^{*}(\CHO\omega) = \CHO\omega$ for all $\phi \in G$.
Consequently, the paths moment map is \emph{strictly equivariant}:
\[
\Psi(\phi \cdot p) = \Ad_{*}(\phi)(\Psi(p)), \qquad \forall \phi \in G,\; p \in \Paths(X).
\]
There is no Souriau cocycle.

%%%%%%%%%%%%%%%%%%%%%%%%%%%%%%%%%%%%%%%%%%%%%%%%%%%%%%%%%%%
\subsubsection{The holonomy group}

The \emph{holonomy group} of the action of $G$ on $(X,\omega)$ is the image
under $\Psi$ of the space of free loops in $X$:
\[
\Gamma = \Psi(\Loops(X)) \subset \cG^{*}.
\]

Let $\ell \in \Loops(X)$ be a free loop.  As shown in \S\ref{sec:Periods},
$\Loops(X)$ decomposes into connected components $L_{d}$ indexed by the
topological degree $d \in \ZZ$ of the associated spheroid $S^{2} \to S^{2}$.
The value $\Psi(\ell)$ depends only on this degree $d$ and is given by
\[
\Psi(\ell)(\delta A) = d \int_{0}^{1} \Surf_{p_{0}}(\delta A(t)\,p_{0}, v_{0})\, dt,
\qquad \delta A \in L\mathfrak{so}(3),
\]
where $p_{0} \in S^{2}$ is a fixed reference point and $v_{0} \in T_{p_{0}}S^{2}$ is
a fixed tangent vector.  The integral on the right-hand side defines a non-zero
left-invariant $1$-form $\alpha_{0} \in \cG^{*}$, and the formula shows that
$\Psi(\ell) = d\,\alpha_{0}$.

Therefore, the holonomy group is the infinite cyclic subgroup
\[
\Gamma = \{\, d\,\alpha_{0} \mid d \in \ZZ \,\} \simeq 4\pi\ZZ,
\]
where the isomorphism with $4\pi\ZZ$ is given by evaluating $\alpha_{0}$
on the generator of the period lattice.
Since every element of $\Gamma$ is a closed momentum,
$\Gamma$ is invariant under the coadjoint action of $G$
(see \S\ref{Closed-momenta-of-a-diffeological-group}).

%%%%%%%%%%%%%%%%%%%%%%%%%%%%%%%%%%%%%%%%%%%%%%%%%%%%%%%%%%%
\subsubsection{The $2$-points and $1$-point moment maps}

Since $\Gamma$ is a coadjoint-invariant subgroup of $\cG^{*}$,
the $2$-points moment map is well defined as
\[
\psi : X \times X \to \cG^{*}/\Gamma,
\quad
\psi(x,x') = \Class(\Psi(p)),
\]
where $p$ is any path from $x$ to $x'$, and $\Class : \cG^{*} \to \cG^{*}/\Gamma$
is the projection.
This map satisfies the Chasles cocycle condition
$\psi(x,x') + \psi(x',x'') = \psi(x,x'')$ in $\cG^{*}/\Gamma$.

Choosing a base point $x_{0} \in X$ (for instance, the constant loop at the north pole
$N \in S^{2}$), we obtain a $1$-point moment map
\[
\mu : X \to \cG^{*}/\Gamma,
\quad
\mu(x) = \psi(x_{0}, x).
\]
Any two such moment maps differ by a constant in $\cG^{*}/\Gamma$.

Because the paths moment map $\Psi$ is strictly equivariant,
the $2$-points moment map $\psi$ and the $1$-point moment map $\mu$
are also strictly equivariant with respect to the coadjoint action on $\cG^{*}/\Gamma$:
\[
\mu(\phi \cdot x) = \Ad_{*}^{\Gamma}(\phi)(\mu(x)), \qquad \forall \phi \in G,\; x \in X.
\]
The Souriau class of the action is trivial: $[\theta] = 0 \in H^{1}(G, \cG^{*}/\Gamma)$.

%%%%%%%%%%%%%%%%%%%%%%%%%%%%%%%%%%%%%%%%%%%%%%%%%%%%%%%%%%%
\subsubsection{The quantum moment map on the prequantum groupoid}

The prequantum groupoid $\TT_{\omega}$ carries the left-right invariant
prequantum $1$-form $\blambda$, satisfying $\Class_{\omega}^{*}(\blambda) = \CHO\omega$
and $\der\blambda = \Trg^{*}\omega - \Src^{*}\omega$.
The \emph{quantum moment map} is the smooth map
\[
\Psi_{\TT} : \cY = \Mor(\TT_{\omega}) \to \cG^{*},
\quad
\Psi_{\TT}([\gamma]_{\omega}) = \widehat{[\gamma]_{\omega}}^{*}(\blambda),
\]
where $\widehat{[\gamma]_{\omega}} : G \to \cY$ is the orbit map on the groupoid.
This map is $G$-equivariant and lifts the paths moment map:
$\Psi = \Psi_{\TT} \circ \Class_{\omega}$.

The restriction of $\Psi_{\TT}$ to the isotropy group $\TT_{\omega, x}$
at any point $x \in X$ gives a homomorphism
\[
\Psi_{\TT}|_{\TT_{\omega, x}} : \TT_{\omega, x} \to \cG^{*}.
\]
Recall from \S\ref{sec:Periods} that $\TT_{\omega, x} \simeq T_{\omega} = \RR/4\pi\ZZ$.
The image of this restriction is precisely the holonomy group $\Gamma \simeq 4\pi\ZZ$.
The kernel is trivial (by definition of $\sim_{\omega}$), so
$\Psi_{\TT}|_{\TT_{\omega, x}}$ is an isomorphism onto its image:
\[
\TT_{\omega, x} \;\xrightarrow{\sim}\; \Gamma \subset \cG^{*}.
\]
This is a diffeological isomorphism between the continuous torus
$T_{\omega} = \RR/4\pi\ZZ$ and the discrete group $\Gamma \simeq 4\pi\ZZ$.
In diffeology, a countable discrete group is diffeologically isomorphic
to a $1$-dimensional torus if and only if the torus is irrational---here,
$T_{\omega} \simeq \U(1)$ is a standard circle, which is \emph{not}
diffeologically isomorphic to a discrete group.
This apparent paradox is resolved by observing that $\Psi_{\TT}$ is not
a diffeomorphism onto its image when $\cG^{*}$ is equipped with the
functional diffeology; rather, the image $\Gamma$ inherits the subset
diffeology from $\cG^{*}$, which for a discrete subgroup of a diffeological
vector space is the discrete diffeology.  The map $\TT_{\omega, x} \to \Gamma$
is a smooth surjective homomorphism from a circle to a discrete group,
which forces the circle to be a quotient of $\RR$ by a \emph{non-discrete}
subgroup---but this contradicts $\Pomega = 4\pi\ZZ$.

The resolution of this subtle point lies in the fact that the quantum moment map
$\Psi_{\TT}$ does \emph{not} restrict to a map into the holonomy group $\Gamma$;
rather, it restricts to a map into the \emph{closure} of $\Gamma$ in $\cG^{*}$.
The holonomy group $\Gamma \simeq 4\pi\ZZ$ is the image of the \emph{integer}
loops (those with integer degree $d$).  But the isotropy group $\TT_{\omega, x}$
contains loops with \emph{non-integer} holonomy modulo $\Pomega$, corresponding
to elements of the torus $T_{\omega} = \RR/4\pi\ZZ$ that are not in the image of
the integer lattice.  The full quantum moment map on the isotropy group takes
values in the $1$-dimensional subspace $\RR \cdot \alpha_{0} \subset \cG^{*}$,
and its image is the whole line $\RR \cdot \alpha_{0}$, not just the discrete
subgroup $\Gamma$.  The projection of this image to $\cG^{*}/\Gamma$ is the
torus $T_{\omega}$ itself.

In summary, the quantum moment map on the isotropy group is the map
\[
\Psi_{\TT}|_{\TT_{\omega, x}} : \TT_{\omega, x} \xrightarrow{\sim} \RR \cdot \alpha_{0} \subset \cG^{*},
\]
which identifies the isotropy torus with the $1$-dimensional subspace of $\cG^{*}$
spanned by $\alpha_{0}$.  The holonomy group $\Gamma$ sits inside this subspace
as the integer lattice $4\pi\ZZ \cdot \alpha_{0}$.

%%%%%%%%%%%%%%%%%%%%%%%%%%%%%%%%%%%%%%%%%%%%%%%%%%%%%%%%%%%
\subsubsection{Comparison with the Wess--Zumino--Witten model}

The strict invariance $\phi^{*}\omega = \omega$ for the action of
$\Loops(\SO(3))$ on $\Loops(S^{2})$ contrasts sharply with the
Wess--Zumino--Witten (WZW) model, where a loop group acts on \emph{itself}
by left multiplication.  In the WZW model, the symplectic form is \emph{not}
strictly invariant; it shifts by the Polyakov--Wiegmann cocycle, and the
action lifts to the prequantum groupoid only after a central extension
(the Kac--Moody group).  The present example shows that when the target
is a coadjoint orbit (here $S^{2}$) and the group acts by the coadjoint
action, the symplectic form is strictly invariant, no central extension
appears, and the moment map is strictly equivariant.

This example thus serves as a \emph{non-anomalous} benchmark within the
GQbP framework, illustrating the geometric origin of anomalies:
an anomaly (central extension) appears precisely when the group action
fails to preserve the parasymplectic form strictly, forcing a shift by
a cocycle.  When the action is strict, as here, the prequantum groupoid
carries a faithful action of the unextended group, and the quantization
is anomaly-free.


%%%%%%%%%%%%%%%%%%%%%%%%%%%%%%%%%%%%%%%%%%%%%%%%%%%%%%%%%%%%%%%%%%%%%%%%%%%%%%%%
  %%
  %% MARK: Explicit Computation of the Souriau Cocycle for Loops(S^2)
  %%
  %%%%%%%%%%%%%%%%%%%%%%%%%%%%%%%%%%%%%%%%%%%%%%%%%%%%%%%%%%%%%%%%%%%%%%%%%%%%%%%%
  \subsection{Explicit Computation of the Souriau Cocycle}
  
  Let us compute the Souriau cocycle explicitly for the action of $G = \Loops(\SO(3))$ on $X = \Loops(S^2)$.
  We will demonstrate that this cocycle is a pure coboundary,
  confirming the absence of a Kac-Moody central extension.
  
  \subsubsection{The Setup}
  We define the geometric components:
  \begin{itemize}[nosep]
    \item $X = \Loops(S^2)$,
    with the parasymplectic form $\omega = \int_0^1 \hat{t}^*(\Surf) \, \dt$,
    where $\Surf$ is the standard area form on $S^2$.
    \item $G = \Loops(\SO(3))$,
    acting pointwise by $(\phi \cdot \ell)(t) = A_t \ell(t)$.
    \item The basepoint $x_0 \in X$ is the constant loop at the north pole $N \in S^2$.
    \item $\cG^*$ is the space of left-invariant $1$-forms on $G$,
    identified with the dual of the loop algebra $L\mathfrak{so}(3)$.
  \end{itemize}
  
  The Souriau cocycle is given by evaluating the 2-points moment map $\psi$ from the basepoint to its translated image:
  \[
    \theta(\phi) = \psi(x_0, \phi \cdot x_0) \in \cG^*/\Gamma, \quad \phi \in G.
    \]
  
  \subsubsection{A Path from $x_0$ to $\phi \cdot x_0$}
  For a given $\phi \in G$ represented by $t \mapsto A_t$,
  the translated basepoint is the loop $\ell_\phi(t) = A_t N$.
  We construct a smooth path $p_\phi : [0,1] \to X$ connecting $x_0$ to $\ell_\phi$ using the exponential map in the group.
  Let $A_t(s) = \exp(s \xi_t)$ be a path in $\SO(3)$ from the identity to $A_t$,
  where $\xi_t = \log A_t \in \mathfrak{so}(3)$.
  We define:
  \[
    p_\phi(s)(t) = A_t(s) N.
    \]
  At $s=0$, $p_\phi(0)(t) = N = x_0(t)$.
  At $s=1$, $p_\phi(1)(t) = A_t N = \ell_\phi(t)$.
  
  \subsubsection{The Paths Moment Map on $p_\phi$}
  The cocycle is evaluated by pulling back the path-primitive $\CHO\omega$ along this path.
  For a tangent vector $\delta A \in L\mathfrak{so}(3)$ at the identity of $G$,
  the paths moment map is:
  \[
    \theta(\phi)(\delta A) = \int_0^1 \omega_{p_\phi(s)}(\delta A \cdot p_\phi(s), \partial_s p_\phi(s)) \, \ds.
    \]
  Substituting the definition of $\omega$ as an integral over $t$,
  we obtain a double integral:
  \[
    \theta(\phi)(\delta A) = \int_0^1 \int_0^1 \Surf_{A_t(s) N}\!\left( \delta A(t) A_t(s) N, \; (\partial_s A_t(s)) N \right) \dt \, \ds.
    \]
  Since $\partial_s A_t(s) = \xi_t A_t(s)$,
  the second vector is $\xi_t A_t(s) N$.
  
  \subsubsection{Using $\SO(3)$ Invariance}
  Because the area form $\Surf$ is invariant under $\SO(3)$,
  we can pull the arguments back by $A_t(s)^{-1}$.
  Noting that $A_t(s)^{-1} \xi_t A_t(s) = \xi_t$ (since $\xi_t$ commutes with its own exponential),
  the integrand at the north pole $N$ becomes:
  \[
    \Surf_N\!\left( \Ad_{\exp(-s \xi_t)} \delta A(t) \cdot N, \; \xi_t \cdot N \right).
    \]
  
  \subsubsection{Geometric Evaluation via the Cross Product}
  To evaluate this,
  we identify $\mathfrak{so}(3) \simeq \RR^3$ equipped with the cross product.
  The infinitesimal action on $S^2 \subset \RR^3$ is $v \cdot P = v \times P$.
  The area form at the north pole is given by the scalar triple product: $\Surf_N(V, W) = N \cdot (V \times W)$.
  
  Let $a = \Ad_{\exp(-s \xi_t)} \delta A(t)$ and $b = \xi_t$.
  The vectors in the tangent plane are $V = a \times N$ and $W = b \times N$.
  Using the vector identity $(a \times N) \times (b \times N) = (a \cdot (N \times b))N = (N \cdot (b \times a))N$,
  the area form simplifies beautifully:
  \[
    \Surf_N(a \times N, b \times N) = N \cdot (b \times a) = N \cdot (\xi_t \times \Ad_{\exp(-s \xi_t)} \delta A(t)).
    \]
  
  \subsubsection{Integration over the Path Parameter $s$}
  We now integrate this scalar product over $s \in [0,1]$.
  The derivative of the adjoint action with respect to $s$ is exactly the cross product with the generator:
  \[
    \frac{\partial}{\partial s} \left( \Ad_{\exp(-s \xi_t)} \delta A(t) \right) = -\xi_t \times \left( \Ad_{\exp(-s \xi_t)} \delta A(t) \right).
    \]
  Therefore,
  the integrand is a perfect exact derivative:
  \[
    N \cdot (\xi_t \times \Ad_{\exp(-s \xi_t)} \delta A(t)) = - N \cdot \frac{\partial}{\partial s} \left( \Ad_{\exp(-s \xi_t)} \delta A(t) \right).
    \]
  By the fundamental theorem of calculus,
  the integral over $s$ evaluates exactly to the boundary values:
  \[
    \int_0^1 - N \cdot \frac{\partial}{\partial s} \left( \Ad_{\exp(-s \xi_t)} \delta A(t) \right) \ds = N \cdot \delta A(t) - N \cdot \Ad_{A_t^{-1}} \delta A(t).
    \]
  
  \subsubsection{Conclusion: The Cocycle is a Coboundary}
  Reinserting the integration over the loop parameter $t$,
  the Souriau cocycle is:
  \[
    \theta(\phi)(\delta A) = \int_0^1 \langle N, \delta A(t) \rangle \, \dt - \int_0^1 \langle N, \Ad_{\phi^{-1}} \delta A(t) \rangle \, \dt \pmod{\Gamma},
    \]
  where $\langle \cdot, \cdot \rangle$ denotes the standard Euclidean dot product in $\RR^3$.
  
  Let us define a base momentum $\mu_0 \in \cG^*$ by:
  \[
    \mu_0(\delta A) = \int_0^1 \langle N, \delta A(t) \rangle \, \dt.
    \]
  The coadjoint action of $\phi$ on $\mu_0$ is precisely $(\Ad_\phi^* \mu_0)(\delta A) = \mu_0(\Ad_{\phi^{-1}} \delta A)$.
  Thus,
  the cocycle takes the exact form:
  \[
    \boxed{\; \theta(\phi) = \mu_0 - \Ad_\phi^* \mu_0 \pmod{\Gamma} \;}
    \]
  
  In Lie group cohomology,
  a cocycle of the form $\mu_0 - \Ad_\phi^* \mu_0$ is,
  by definition,
  a \emph{coboundary}.
  Its cohomology class is strictly zero.
  This explicit computation proves definitively that there is no central extension required to lift the action of $\Loops(\SO(3))$ on $\Loops(S^2)$.
  The action is purely Hamiltonian,
  and the Kac-Moody anomaly is absent.

  %%%%%%%%%%%%%%%%%%%%%%%%%%%%%%%%%%%%%%%%%%%%%%%%%%%%%%%%%%%%%%%%%%%%%%%%%%%%%%%%
  %%
  %% MARK: Appendix: Canonical Trivialization of the Polyakov-Wiegmann Cocycle
  %%
  %%%%%%%%%%%%%%%%%%%%%%%%%%%%%%%%%%%%%%%%%%%%%%%%%%%%%%%%%%%%%%%%%%%%%%%%%%%%%%%%
  
  \appendix
  
  %%%%%%%%%%%%%%%%%%%%%%%%%%%%%%%%%%%%%%%%%%%%%%%%%%%%%%%%%%%%%%%%%%%%%%%%%%%%%%%%
  %%
  %% MARK: Appendix: Canonical Trivialization of the Polyakov-Wiegmann Cocycle
  %%
  %%%%%%%%%%%%%%%%%%%%%%%%%%%%%%%%%%%%%%%%%%%%%%%%%%%%%%%%%%%%%%%%%%%%%%%%%%%%%%%%
  
  \section{Canonical Trivialization of the Polyakov-Wiegmann Cocycle}
  \label{app:canonical_pw}
  
  In Section 4,
  we established that the parasymplectic form $\omega$ on the loop space $X = \Loops(\SU(2))$ shifts under the pointwise action of a loop $\ell$ by an exact differential:
  $m(\ell)^*\omega = \omega + d\theta$.
  This 1-form $\theta = \Psi^*(\Komega)$ was derived using a smooth homotopy $s \mapsto \ell_s$ connecting the identity loop to $\ell$.
  
  One might object that this construction depends on the arbitrary choice of the path $s \mapsto \ell_s$.
  In this appendix,
  we expand $\Psi^*(\Komega)$ explicitly on a plot of $X$.
  We demonstrate that this 1-form decomposes exactly into the canonical Polyakov-Wiegmann 1-form (which depends only on the boundary loop $\ell$) and an exact differential $df$ (which absorbs all dependence on the chosen homotopy).
  This proves that the canonical Polyakov-Wiegmann shift is the unique geometric invariant of the fusion product.
  
  Remark on Generality:
  The explicit decomposition $\Psi^*(\Komega) = \theta^{PW}_\ell + df$ derived below relies strictly on the Lie group structure of $Y = \SU(2)$ and the existence of the Maurer-Cartan basis.
  For a general contact 3-manifold lacking a global group structure,
  pointwise loop multiplication is undefined,
  and this specific algebraic trivialization does not apply.
  However,
  the underlying geometric principle—that continuous symmetries of a parasymplectic space shift the prequantum potential by an exact differential,
  generating a central extension—remains universally valid in the GQbP framework via Cartan's formula.
  
  %%%%%%%%%%%%%%%%%%%%%%%%%%%%%%%%%%%%%%%%%%%%%%%%%%%%%%%%%%%%%%%%%%%%%%%%%%%%%%%%
  %% Subsection: Setup and Maurer-Cartan Forms
  %%%%%%%%%%%%%%%%%%%%%%%%%%%%%%%%%%%%%%%%%%%%%%%%%%%%%%%%%%%%%%%%%%%%%%%%%%%%%%%%
  \subsection{Setup and Maurer-Cartan Forms}
  
  Let $r \mapsto \gamma_r$ be a plot in $X$.
  A tangent vector to this plot is the variation $\delta\gamma$.
  By definition of the pullback,
  the 1-form $\theta$ evaluated on this variation is:
  \[
    \theta(\delta\gamma) = (\Psi^*(\Komega))_\gamma(\delta\gamma) = \Komega_{\Psi(\gamma)}(\Psi_* \delta\gamma).
    \]
  The map $\Psi$ sends the loop $\gamma$ to the path of loops $s \mapsto \ell_s \gamma$.
  Expanding the path-primitive $\Komega = \CHO\omega = \CHO^2\Vol$,
  we obtain a double integral over the homotopy parameter $s \in [0,1]$ and the loop parameter $t \in S^1$:
  \[
    \theta(\delta\gamma) = \int_0^1 ds \int_0^1 \dt \, \Vol_{\ell_s\gamma} \left( \partial_t(\ell_s \gamma), \partial_s(\ell_s \gamma), \ell_s \delta\gamma \right).
    \]
  
  Because the volume form $\Vol$ is bi-invariant on $\SU(2)$,
  we can right-translate the three argument vectors by $(\ell_s \gamma)^{-1}$ to evaluate them at the identity $e \in \SU(2)$.
  Let us define the following Maurer-Cartan forms.
  For the homotopy $\ell_s$, we use left-invariant forms:
  \[
    A = \ell_s^{-1} \partial_t \ell_s, \qquad B = \ell_s^{-1} \partial_s \ell_s.
    \]
  For the loop $\gamma$, we use right-invariant forms:
  \[
    C = \dot{\gamma}\gamma^{-1}, \qquad D = \delta\gamma \gamma^{-1}.
    \]
  
  Applying the Leibniz rule and right-translating by $\gamma^{-1} \ell_s^{-1}$,
  the three vectors at the identity become:
  \begin{align*}
    v_1 &= (\partial_t \ell_s \cdot \gamma + \ell_s \cdot \dot{\gamma}) \gamma^{-1} \ell_s^{-1} = A + C, \\
    v_2 &= (\partial_s \ell_s \cdot \gamma) \gamma^{-1} \ell_s^{-1} = B, \\
    v_3 &= (\ell_s \cdot \delta\gamma) \gamma^{-1} \ell_s^{-1} = D.
  \end{align*}
  (Note that the adjoint conjugation by $\ell_s$ on $C$ and $D$ vanishes inside the bi-invariant volume form).
  
  At the identity,
  the volume form evaluates as $\Vol_e(X, Y, Z) \propto \Tr(X[Y, Z])$.
  The integrand is therefore:
  \[
    \Tr\bigl( (A+C)[B, D] \bigr) = \Tr(A[B, D]) + \Tr(C[B, D]).
    \]
  Using the cyclic property of the trace,
  $\Tr(A[B, D]) = \Tr([A, B]D)$ and $\Tr(C[B, D]) = \Tr(B[C, D])$.
  The integrand simplifies to:
  \[
    \Tr([A, B]D) + \Tr(B[C, D]).
    \]
  
  %%%%%%%%%%%%%%%%%%%%%%%%%%%%%%%%%%%%%%%%%%%%%%%%%%%%%%%%%%%%%%%%%%%%%%%%%%%%%%%%
  %% Subsection: Integration over the Homotopy
  %%%%%%%%%%%%%%%%%%%%%%%%%%%%%%%%%%%%%%%%%%%%%%%%%%%%%%%%%%%%%%%%%%%%%%%%%%%%%%%%
  \subsection{Integration over the Homotopy}
  
  We first integrate this expression over the homotopy parameter $s \in [0,1]$.
  The Maurer-Cartan structure equation for the mixed partial derivatives of $\ell_s(t)$ yields:
  \[
    \partial_t B - \partial_s A = [A, B].
    \]
  Substituting this into the first term of our integrand gives:
  \[
    \int_0^1 \Tr([A, B]D) \, ds = \int_0^1 \Tr(\partial_t B \cdot D) \, ds - \int_0^1 \Tr(\partial_s A \cdot D) \, ds.
    \]
  Crucially,
  the term $D = \delta\gamma \gamma^{-1}$ depends only on $\gamma$ and is completely independent of $s$.
  We can therefore pull it out of the $s$-integrals.
  
  For the second integral,
  the fundamental theorem of calculus gives:
  \[
    \int_0^1 \partial_s A \, ds = A(1) - A(0) = \ell^{-1}\dot{\ell} - 0 = \ell^{-1}\dot{\ell},
    \]
  since $\ell_0 = \id$ and $\ell_1 = \ell$.
  
  For the remaining terms involving $B$,
  let us define the Lie algebra-valued function $g(t) = \int_0^1 B(s, t) \, ds$.
  Then $\int_0^1 \partial_t B \, ds = \partial_t g(t)$.
  The total integral over $s$ is exactly:
  \[
    \int_0^1 (\dots) ds = - \Tr\bigl( (\ell^{-1}\dot{\ell}) D \bigr) + \Tr(\partial_t g \cdot D) + \Tr(g[C, D]).
    \]
  
  %%%%%%%%%%%%%%%%%%%%%%%%%%%%%%%%%%%%%%%%%%%%%%%%%%%%%%%%%%%%%%%%%%%%%%%%%%%%%%%%
  %% Subsection: Integration by Parts and the Exact Differential
  %%%%%%%%%%%%%%%%%%%%%%%%%%%%%%%%%%%%%%%%%%%%%%%%%%%%%%%%%%%%%%%%%%%%%%%%%%%%%%%%
  \subsection{Integration by Parts and the Exact Differential}
  
  Finally,
  we integrate over the loop parameter $t \in S^1$.
  The first term depends only on the boundary loop $\ell$ and the variation of $\gamma$.
  This is exactly the canonical Polyakov-Wiegmann 1-form:
  \[
    \theta^{PW}_\ell(\delta\gamma) \propto - \int_0^1 \Tr\bigl( (\ell^{-1}\dot{\ell}) (\delta\gamma \gamma^{-1}) \bigr) \, \dt.
    \]
  
  We now evaluate the remaining two terms:
  \[
    \int_0^1 \bigl( \Tr(\partial_t g \cdot D) + \Tr(g[C, D]) \bigr) \, \dt.
    \]
  Integrating the first term by parts over $t$ yields $-\int_0^1 \Tr(g \cdot \partial_t D) \, \dt$.
  The boundary terms vanish because $g(0) = g(1) = 0$ (since $\ell_s$ are closed loops for all $s$).
  The remaining integral is:
  \[
    - \int_0^1 \Tr\bigl( g \cdot (\partial_t D - [C, D]) \bigr) \, \dt.
    \]
  
  We recall that $C = \dot{\gamma}\gamma^{-1}$ and $D = \delta\gamma \gamma^{-1}$.
  A standard identity in the calculus of variations on Lie groups relates the time derivative of the variation to the variation of the velocity:
  \[
    \partial_t D - [C, D] = \delta C.
    \]
  Substituting this identity,
  the integral becomes:
  \[
    - \int_0^1 \Tr(g \cdot \delta C) \, \dt.
    \]
  Because the function $g(t)$ depends \emph{only} on the chosen homotopy $s \mapsto \ell_s$ and is completely independent of $\gamma$,
  the variation operator $\delta$ passes directly through it:
  \[
    - \int_0^1 \Tr(g \cdot \delta C) \, \dt = \delta \left( - \int_0^1 \Tr(g(t) \dot{\gamma}(t)\gamma(t)^{-1}) \, \dt \right) = df(\delta\gamma),
    \]
  where $f \in \Cinfty(X, \RR)$ is a globally defined smooth function on the loop space.
  
  \subsection{Conclusion}
  We have proven that the 1-form generated by the diffeological homotopy decomposes exactly as:
  \[
    \Psi^*(\Komega) = \theta^{PW}_\ell + df.
    \]
  The exact differential $df$ absorbs all dependence on the arbitrary choice of the path $s \mapsto \ell_s$.
  If one chooses a different homotopy,
  the resulting 1-form changes only by an exact differential.
  In the prequantum groupoid,
  such an exact shift corresponds to a canonical gauge isomorphism.
  Therefore,
  the canonical Polyakov-Wiegmann 1-form $\theta^{PW}_\ell$ is the unique,
  objective geometric invariant of the fusion product,
  rigorously justified by the diffeological path-primitive operator.
  
  %%%%%%%%%%%%%%%%%%%%%%%%%%%%%%%%%%%%%%%%%%%%%%%%%%%%%%%%%%%%%%%%%%%%%%%%%%%%%%%%
  %%
  %% MARK: Appendix: Explicit Computation of the Paths-Primitive $\KVol$
  %%
  %%%%%%%%%%%%%%%%%%%%%%%%%%%%%%%%%%%%%%%%%%%%%%%%%%%%%%%%%%%%%%%%%%%%%%%%%%%%%%%%
  
  \section{Decomposition of the Paths-Primitive $\KVol$}
  \label{Decomposition of the Paths-Primitive}
  
  We give the detailed computation of the parasymplectic form
  $\KVol$ on the space of paths $\Paths(Y)$,
  where $Y$ is a contact $3$-manifold over a symplectic surface $(\Sigma, \epsilon)$,
  with connection $\alpha$ satisfying $\pr^*\epsilon = d\alpha$.
  This expression gives by restriction $\omega = \KVol|_X$ on the loop space $X = \Loops(Y)$.
  
  \subsection{Wedge Product Formula}
  
  For a $1$-form $\alpha$ and a $2$-form $d\alpha$,
  the $3$-form $\Vol = \alpha \wedge d\alpha$ evaluated on three tangent vectors $u, v, w$ is:
  \[
    (\alpha \wedge d\alpha)(u, v, w)
    = \alpha(u) \, d\alpha(v, w)
    + \alpha(v) \, d\alpha(w, u)
    + \alpha(w) \, d\alpha(u, v).
    \]
  
  \subsection{Application to a Path}
  
  Let $\gamma \in X$ be a path in $Y$,
  and let $\delta\gamma, \delta'\gamma \in T_\gamma X$ be two tangent vectors (variations of the path).
  For each $t \in [0,1]$,
  denote:
  \[
    y_t = \gamma(t) \in Y,
    \quad
    x_t = \pr(y_t) \in \Sigma.
    \]
  The tangent vectors at $y_t$ are:
  \[
    \dot y_t = \dot\gamma(t) \in \Tgt_{y_t}Y,
    \quad
    \delta y_t = \delta\gamma(t) \in \Tgt_{y_t}Y,
    \quad
    \delta' y_t = \delta'\gamma(t) \in \Tgt_{y_t}Y.
    \]
  Their projections to the base are:
  \[
    \dot x_t = \pr_*(\dot y_t) \in \Tgt_{x_t}\Sigma,
    \quad
    \delta x_t = \pr_*(\delta y_t) \in \Tgt_{x_t}\Sigma,
    \quad
    \delta' x_t = \pr_*(\delta' y_t) \in \Tgt_{x_t}\Sigma.
    \]
  
  \subsection{Development of the Integrand}
  
  Applying the wedge product formula to
  $u = \dot y_t$,
  $v = \delta y_t$,
  $w = \delta' y_t$,
  we obtain:
  
  \begin{multline*}
    (\alpha \wedge d\alpha)_{y_t}\bigl( \dot y_t, \delta y_t, \delta' y_t \bigr)
    = \alpha_{y_t}(\dot y_t) \; d\alpha_{y_t}(\delta y_t, \delta' y_t) \\
    + \alpha_{y_t}(\delta y_t) \; d\alpha_{y_t}(\delta' y_t, \dot y_t)
    + \alpha_{y_t}(\delta' y_t) \; d\alpha_{y_t}(\dot y_t, \delta y_t).
  \end{multline*}
  
  Using the curvature condition $d\alpha = \pr^*\epsilon$,
  the terms involving $d\alpha$ simplify:
  
  \begin{align*}
    d\alpha_{y_t}(\delta y_t, \delta' y_t)
    &= \pr^*\epsilon_{y_t}(\delta y_t, \delta' y_t)
    = \epsilon_{x_t}(\delta x_t, \delta' x_t), \\
    d\alpha_{y_t}(\delta' y_t, \dot y_t)
    &= \pr^*\epsilon_{y_t}(\delta' y_t, \dot y_t)
    = \epsilon_{x_t}(\delta' x_t, \dot x_t), \\
    d\alpha_{y_t}(\dot y_t, \delta y_t)
    &= \pr^*\epsilon_{y_t}(\dot y_t, \delta y_t)
    = \epsilon_{x_t}(\dot x_t, \delta x_t).
  \end{align*}
  
  Substituting these identities,
  the integrand becomes:
  \begin{multline*}
    (\alpha \wedge d\alpha)_{y_t}\bigl( \dot y_t, \delta y_t, \delta' y_t \bigr)
    = \alpha_{y_t}(\dot y_t) \; \epsilon_{x_t}(\delta x_t, \delta' x_t) \\
    + \alpha_{y_t}(\delta y_t) \; \epsilon_{x_t}(\delta' x_t, \dot x_t)
    + \alpha_{y_t}(\delta' y_t) \; \epsilon_{x_t}(\dot x_t, \delta x_t).
  \end{multline*}
  
  Using the skew-symmetry of the $2$-form $\epsilon$,
  we have $\epsilon_{x_t}(\delta' x_t, \dot x_t) = -\epsilon_{x_t}(\dot x_t, \delta' x_t)$.
  Thus:
  \begin{multline*}
    (\alpha \wedge d\alpha)_{y_t}\bigl( \dot y_t, \delta y_t, \delta' y_t \bigr)
    = \alpha_{y_t}(\dot y_t) \; \epsilon_{x_t}(\delta x_t, \delta' x_t) \\
    + \alpha_{y_t}(\delta' y_t) \; \epsilon_{x_t}(\dot x_t, \delta x_t)
    - \alpha_{y_t}(\delta y_t) \; \epsilon_{x_t}(\dot x_t, \delta' x_t).
  \end{multline*}
  
  \subsection{Integration Along the path}
  
  The parasymplectic form $\omega = \KVol|_X$ is obtained by integrating
  this expression over $t \in [0,1]$:
  \begin{multline*}
    \omega_\gamma(\delta\gamma, \delta'\gamma) = \int_0^1 \alpha_{y_t}(\dot y_t) \; \epsilon_{x_t}(\delta x_t, \delta' x_t) \, \dt \\
    + \int_0^1 \Bigl[ \alpha_{y_t}(\delta' y_t) \; \epsilon_{x_t}(\dot x_t, \delta x_t) - \alpha_{y_t}(\delta y_t) \; \epsilon_{x_t}(\dot x_t, \delta' x_t) \Bigr] \, \dt \;.
  \end{multline*}
  
  
  %%%%%%%%%%%%%%%%%%%%%%%%%%%%%%%%%%%%%%%%%%%%%%%%%%%%%%%%%%
  %%
  %% MARK: Bibliography
  %%
  %%%%%%%%%%%%%%%%%%%%%%%%%%%%%%%%%%%%%%%%%%%%%%%%%%%%%%%%%%
  
  \begin{thebibliography}{IKZ10} % Use 99 or a number larger than your total entries
    
    \bibitem[DI83]{DI83}
    Paul Donato and Patrick Iglesias.
    {\em Exemple de groupes diff\'erentiels : flots irrationnels sur le tore}.
    Preprint CPT-83/P.1524, Centre de Physique Th\'eorique, Marseille, July 1983. Published in {\em Comptes Rendus de l'Acad\'emie des Sciences}, 301(4), Paris, 1985.
    \texttt{http://math.huji.ac.il/~piz/documents/EDGDFISLT.pdf}
    
    \bibitem[GIZ23]{GIZ23}
    Serap G{\"u}rer and Patrick Iglesias-Zemmour.
    {\em Orbifolds as stratified diffeologies}.
    {Differential Geometry and its Applications}, 86:101969, 2023.
    
    \bibitem[GIZ25]{GIZ25}
    Serap Gürer and Patrick Iglesias-Zemmour,
    \emph{On Diffeology of Orbit Spaces},
    (Preprint available at \texttt{http://math.huji.ac.il//~piz/documents/ODOOS.pdf}).
    
    \bibitem[Hat02]{Hat02}
    Allen Hatcher,
    \emph{Algebraic Topology},
    Cambridge University Press, 2002.
    
    \bibitem[Hor24]{Hor24}
    Peter A. Horvathy.
    {\em Prequantisation from the path integral viewpoint}.
    Preprint arXiv:2402.17629, 2024.
    
    \bibitem[PIZ85]{PIZ85}
    Patrick Iglesias.
    {\em Fibr\'es diff\'eologiques et homotopie}.
    Th\`ese de doctorat d'\'etat, Universit\'e de Provence, Marseille, 1985.
    
    \bibitem[PIZ95]{PIZ95}
    Patrick Iglesias.
    {\em La trilogie du moment}.
    Ann. Inst. Fourier, 45(3):825--857, 1995.
    
    \bibitem[IKZ10]{IKZ10}
    Patrick Iglesias, Yael Karshon, and Moshe Zadka.
    {\em Orbifolds as diffeologies}.
    Transactions of the American Mathematical Society, Vol. 362, Number 6, Pages 2811--2831, Providence R.I., 2010.
    
    \bibitem[PIZ10]{PIZ10}
    Patrick Iglesias-Zemmour,
    \emph{The moment maps in diffeology},
    Mem. Amer. Math. Soc. \textbf{204} (2010), no. 962, viii+101 pp.
    
    \bibitem[PIZ13]{PIZ13}
    Patrick Iglesias-Zemmour.
    {\em Diffeology}.
    Mathematical Surveys and Monographs, 185, Am. Math. Soc., Providence, RI, 2013.
    (Revised version by Beijing World Publishing Corp., Beijing, 2022.)
    
    \bibitem[PIZ22]{PIZ22}
    Patrick Iglesias-Zemmour.
    {\em Diffeology}.
    Reprint by Beijing World Publishing Corporation, Ltd. Beijing Branch, 2022. ISBN 978-7-5192-9608-7.
    Accessible at \\
    \url{http://math.huji.ac.il/~piz/documents/Diffeology.pdf}
    
    \bibitem[PIZ15]{PIZ15}
    \bysame
    {\em Example of singular reduction in symplectic diffeology}.
    Proceedings of the American Mathematical Society, 144(3):1309--1324, 2016.
    
    \bibitem[IZP21]{IZP21}
    Patrick Iglesias-Zemmour and Elisa Prato.
    {\em Quasifolds, diffeology and noncommutative geometry}.
    J. Noncommut. Geom., 15(2):735--759, 2021.
    
    \bibitem[PIZ24]{PIZ24}
    Patrick Iglesias-Zemmour,
    \emph{{\v C}ech-de Rham bicomplex in diffeology},
    Israel J. Math. \textbf{259} (2024), no. 1, 239--276.
    
    \bibitem[PIZ25]{PIZ25}
    \bysame,
    \emph{The Boman Paradox --- When topology plus algebra do not define geometry},
    To appear in the Mathematical Intelligencer. Preprint 2025. Available at \url{http://math.huji.ac.il/~piz/documents/TBP.pdf}.
    
    \bibitem[MW74]{MW74}
    Jerrold E. Marsden and Alan Weinstein,
    \emph{Reduction of symplectic manifolds with symmetry},
    Rep. Mathematical Phys. \textbf{5} (1974), no. 1, 121--130.
    
    \bibitem[Sat56]{Sat56}
    Ichiro Satake,
    \emph{On a generalization of the notion of manifold},
    Proceedings of the National Academy of Sciences \textbf{42} (1956), 359--363.
    
    \bibitem[Sou70]{Sou70}
    Jean-Marie Souriau.
    {\em Structure des syst\`emes dynamiques}.
    Dunod, Paris, 1970.
    
    \bibitem[Sou75]{Sou75}
    \bysame %Jean-Marie Souriau.
    \newblock \emph{Construction explicite de l'indice de Maslov. Applications}.
    \newblock In: Group theoretical methods in physics (Fourth Internat. Colloq., Nijmegen, 1975), Lecture Notes in Phys., Vol. 50, Springer, Berlin, 1976, pp. 117--148.
    
    \bibitem[Sou04]{Sou04}
    Jean-Marie Souriau,
    \emph{C’est Quantique~? donc c’est géométrique},
    in: Dominique Flament (dir),
    \emph{Feuilletages -quantification géométrique : textes des journées d’étude des 16 et 17 octobre 2003, Paris, Fondation Maison des Sciences de l’Homme},
    Série Documents de travail (Équipe F2DS), 2004.
    
  \end{thebibliography}
  
  
\end{document}
